\documentclass[a4paper]{article}

\usepackage{amsmath}
\usepackage{amssymb}
\usepackage{cancel}
\usepackage{stellar}
\usepackage{parskip}
\usepackage{fullpage}
\usepackage{caption}
\usepackage{tikz}
\usepackage{wrapfig}
\usepackage{xfrac}
\usepackage{witharrows}
\usepackage{mathrsfs}

\WithArrowsOptions{displaystyle}

\usetikzlibrary{calc}
\usetikzlibrary{decorations.markings}
\usetikzlibrary{cd}
\usetikzlibrary{positioning}
\usetikzlibrary{angles}

% Cancel color
\renewcommand{\CancelColor}{\color{red}}

\title{Psicologia generale}
\author{Stellar Inc.}
\date{}

\begin{document}

\maketitle
\tableofcontents

\section{Sviluppo storico}

\subsection{Che cos'è la psicologia?}

\begin{snippetdefinition}{Psicologia}
    La \emph{psicologia} è lo studio scientifico del comportamento degli individui e dei loro processi mentali, per questo motivo si fonda su un metodo scientifico per cui è possibile giungere a conclusioni tramite la raccolta di prove. In quanto scienza della salute, la psicologia cerca di migliorare la qualità di vita di ogni individuo.
\end{snippetdefinition}
\begin{snippetdefinition}{Comportamento}
    Il \emph{comportamento} è l'insieme delle azioni attraverso cui gli organismi rispondono agli stimoli ed interagiscono con l'ambiente. Gli psicologi esaminano cosa fa l'individuo e come esso si comporta nel farlo, considerando ciascun fenomeno comportamentale calato nell'ambiente sociale e culturale in cui è prodotto.
\end{snippetdefinition}
Nella maggior parte dei casi l'oggetto dell'analisi psicologica sono gli esseri umani (ad esempio un neonato che fa le prime esperienze, un atleta alla ricerca di motivazione). Non è però possibile capire le azioni degli esseri umani senza capirne anche i \emph{processi mentali}, i meccanismi della mente. Si tratta dell'aspetto più importante dell'indagine psicologica per cui si sono sviluppate tecniche per studiare le manifestazioni ed i processi della mente.
\begin{snippetexample}{Comportamento}
    Facciamo un esempio, ci poniamo di indagare cosa influenza la rapidità di lettura, per cui andremo a chiedere ai partecipanti dell'osservazione di leggere delle parole; in questo caso la \underline{rapidità di lettura} consiste in un comportamento che vogliamo analizzare, mentre l'azione di leggere le parole è il \emph{compito/strumento} attraverso cui verrà condotta l'analisi. In questo caso, variabili come
    \begin{itemize}
        \item \emph{lunghezza delle parole:} \texttt{VIA} vs \texttt{RAGIONE},
        \item \emph{frequenza delle parole:} \texttt{TOPO} vs \texttt{CUPO}
    \end{itemize} sono potenzialmente informative dei \emph{processi mentali} che guidano il comportamento. Parole più frequenti sono meglio discriminabili dalla nostra mente (presenza di un sistema memoria, sensibilità al numero di volte in cui una parola compare), mentre parole più lunghe sono associate a tempi di lettura più lunghi.
\end{snippetexample}

\subsection{Wundt e l'introspezione}
Con \emph{Wilhelm Wundt} nasce la psicologia moderna: egli applica i metodi della fisiologia ai processi e ai contenuti della coscienza umana.
\begin{snippetdefinition}{Esperienza diretta}
    Per Wundt, l'\emph{esperienza diretta}, o immediata, è l'oggetto di studio della psicologia. Il metodo privilegiato per rilevarla è l'\emph{introspezione}.
\end{snippetdefinition}
\begin{snippetdefinition}{Introspezione}
    L'\emph{introspezione} è un metodo scientifico che richiede il controllo accurato dello stimolo capace di produrre l'evento mentale oggetto di osservazione e la stesura di un resoconto subito dopo l'osservazione dell'evento.
\end{snippetdefinition}
Per garantire la validità scientifica dell'introspezione, Wundt formula alcune linee guida. L'osservatore, quando possibile,
\begin{itemize}
    \item deve avere la possibilità di stabilire quando introdurre il processo;
    \item deve trovarsi in una condizione di sforzo attentivo;
    \item deve poter ripetere diverse volte la medesima osservazione.
\end{itemize}
La condizione sperimentale, invece, deve prendere in considerazione variazioni in termini di \emph{intensità} e di \emph{qualità} della stimolazione.
La bontà dei risultati di questo metodo dipende dall'abilità e dall'esperienza dell'osservatore. Il controllo sperimentale è ottenuto non solo attraverso il rispetto delle linee guida, ma anche facendo ricorso a osservatori in grado di raccogliere dati attendibili grazie alle loro risorse attentive e alla rapidità nella fase di osservazione e di stesura del resoconto.
Secondo Wundt, il metodo dell'\emph{auto-osservazione sperimentale} permette di rilevare processi mentali come l'\emph{appercezione}, la \emph{volontà} e le \emph{emozioni}. Tuttavia, egli ritiene che tale metodo non sia adatto allo studio dei processi mentali di ordine superiore.
\begin{snippetexample}{Introspezione}
    Lo sperimentatore presenta al partecipante uno stimolo, ad esempio un suono o un colore, e gli chiede poi di riferire le proprie introspezioni. Il resoconto deve basarsi esclusivamente sui processi percettivi, evitando qualsiasi interpretazione soggettiva. Nel caso di un suono, il partecipante può esprimere un giudizio sulla sua \emph{intensità}, ma non deve riportare reazioni come fastidio, paura o altre valutazioni personali.
\end{snippetexample}
\subsection{Lo strutturalismo}
\begin{snippetdefinition}{Strutturalismo}
    Lo \emph{strutturalismo} è un approccio psicologico fondato sull'\emph{elementarismo} e sul metodo dell'\emph{introspezione}. Compare per la prima volta nel 1898 in un articolo dello psicologo inglese \emph{E. B. Titchener}, allievo di Wundt.
\end{snippetdefinition}
Secondo Titchener, il primo passo per comprendere la mente consiste nello scoprirne la struttura, scomponendola nei suoi elementi primari. È quindi necessario individuare come e perché tali elementi si combinano.
Per Titchener, l'esperienza cosciente è costituita da \emph{percezioni}, \emph{idee}, \emph{emozioni} o \emph{sentimenti}. A queste corrispondono tre componenti fondamentali:
\begin{itemize}
    \item le \emph{sensazioni}, alla base delle percezioni;
    \item le \emph{immagini mentali}, alla base delle idee;
    \item gli \emph{stati affettivi}, alla base delle emozioni.
\end{itemize}
Alle due caratteristiche già individuate da Wundt, cioè \emph{qualità} e \emph{intensità}, Titchener aggiunge la \emph{durata} e la \emph{chiarezza}. Ad esempio, uno stimolo visivo può distinguersi più o meno chiaramente rispetto a un altro e può presentarsi per un certo intervallo di tempo.
\begin{snippetdefinition}{Errore dello stimolo}
    L'\emph{errore dello stimolo} consiste nell'attribuzione di significati e valori soggettivi ai dati oggettivi dell'esperienza.
\end{snippetdefinition}
Lo strutturalismo si fonda sul presupposto che tutte le esperienze mentali degli esseri umani possano essere comprese attraverso la combinazione delle componenti di base.
Molti psicologi criticarono lo strutturalismo principalmente per tre ragioni:
\begin{itemize}
    \item era \emph{riduzionista}, perché riconduceva la complessità dell'esperienza umana ad alcune componenti sottostanti, come sensazioni, immagini e stati affettivi;
    \item era \emph{elementarista}, perché non studiava direttamente la complessità o la totalità di un comportamento;
    \item era \emph{mentalista}, perché analizzava soltanto i resoconti verbali della consapevolezza cosciente, ignorando chi non era in grado di descrivere le proprie introspezioni, come bambini o soggetti mentalmente disturbati.
\end{itemize}

\subsection{James e il funzionalismo}
La psicologia \emph{funzionale} e \emph{descrittiva} si oppone al pensiero di Wundt e Titchener, perché è più interessata alle funzioni della mente che alla sua struttura.
Negli Stati Uniti spicca la figura di \emph{William James}, a cui si ispira la psicologia funzionalista. Il punto di partenza è la teoria dell'evoluzione e lo stretto legame tra ambiente e individuo.
\begin{snippetdefinition}{Funzionalismo}
    Il \emph{funzionalismo} è un orientamento psicologico che attribuisce importanza alle funzioni della mente e alle abitudini apprese, poiché esse permettono all'organismo di adattarsi all'ambiente e di funzionare con efficacia.
\end{snippetdefinition}
Secondo i funzionalisti, la domanda chiave a cui la ricerca deve rispondere è: ``qual è la funzione o lo scopo del comportamento?''.
Secondo James:
\begin{itemize}
    \item la psicologia è la scienza della vita mentale, dei suoi fenomeni e delle sue condizioni;
    \item l'oggetto della psicologia è l'esperienza immediata, cioè i fenomeni, che non sono combinazioni di elementi svelati dall'introspezione;
    \item le sensazioni sono l'esito di un processo articolato di inferenza e di astrazione;
    \item la coscienza è un flusso continuo, condizionato dal cervello;
    \item il metodo comparativo, cioè lo studio di diverse popolazioni di soggetti, permette di compensare la soggettività dell'introspezione;
    \item i fenomeni psichici sono funzioni dell'organismo finalizzate all'adattamento all'ambiente.
\end{itemize}
Mentre lo strutturalismo afferma la necessità di una psicologia \emph{pura}, esclusivamente sperimentale, i funzionalisti sostengono l'esigenza di una psicologia impegnata anche sul versante applicativo. La spiegazione di un comportamento deve quindi considerare anche il contesto, cioè l'ambiente in cui il comportamento stesso si manifesta.
Con il funzionalismo nasce anche un nuovo concetto di esperimento di laboratorio: l'oggetto di indagine non è più la coscienza, ma il comportamento nella sua totalità.
\subsection{La psicologia della Gestalt}
\begin{snippetdefinition}{Gestalt}
    Il termine \emph{Gestalt}, cioè forma, indica una configurazione unitaria e organizzata che differisce dalla semplice somma delle sue singole parti.
\end{snippetdefinition}
La nascita della psicologia della Gestalt viene fatta risalire al 1912, anno in cui \emph{Wertheimer} pubblica il lavoro sul movimento apparente, o \emph{fenomeno phi}. Oltre a Wertheimer, i principali esponenti di questo movimento sono \emph{Koffka} e \emph{Köhler}.
\begin{snippetdefinition}{Fenomeno phi}
    Il \emph{fenomeno phi} è un fenomeno percettivo per cui l'osservatore percepisce un unico stimolo in movimento a partire da stimoli luminosi presentati in successione.
\end{snippetdefinition}
Il fenomeno phi si lega alla \emph{persistenza percettiva dell'oggetto}: tendiamo a integrare gli stimoli che arrivano continuamente agli organi di senso come se fossero originati da singoli oggetti permanenti, anche quando non lo sono.
\begin{snippetexample}{Movimento apparente}
    Il movimento apparente si ottiene, ad esempio, proiettando su uno sfondo scuro uno stimolo luminoso e, dopo un intervallo molto breve, un secondo stimolo luminoso vicino alla posizione del primo. In questo modo l'osservatore percepisce un unico stimolo in movimento.
\end{snippetexample}
Questo fenomeno mette in crisi la corrispondenza tra stimolazione e percezione e risulta inspiegabile sulla base di una teoria elementaristica della percezione, incapace di rendere conto della percezione di un movimento continuo.
Lo studio sperimentale del movimento apparente, o \emph{stroboscopico}, è alla base del movimento cinematografico.
Costruire una teoria sulla percezione significa individuare le precise regole dell'interazione tra le parti presenti nel campo fenomenico. Tali regole sono i \emph{principi di unificazione formale}.
La psicologia della Gestalt rifiuta l'elementarismo e concentra l'attenzione sui concetti di \emph{struttura} e \emph{organizzazione}. L'analisi deve andare dall'intero alle sue parti costituenti, al contrario dell'elementarismo, che si limitava a sommare o combinare le proprietà delle parti elementari.
\begin{snippetdefinition}{Isomorfismo}
    L'\emph{isomorfismo} è il principio secondo cui l'organizzazione degli eventi psicologici è riconducibile alle proprietà strutturali degli eventi neurofisiologici corrispondenti.
\end{snippetdefinition}
\subsection{Il comportamentismo}
Contemporaneamente alla nascita della psicologia della Gestalt in Germania, negli Stati Uniti nasce il \emph{comportamentismo}, una grande scuola psicologica che segna la psicologia fino al secondo dopoguerra. Si tratta di un movimento contrario alla tradizione mentalista, poiché pone al centro della psicologia l'indagine del comportamento.
Secondo \emph{John B. Watson}, la psicologia deve occuparsi del comportamento osservabile e deve avvalersi di metodi suscettibili degli stessi controlli utilizzati nelle scienze naturali. Di conseguenza, deve rifiutare il modello dell'introspezione.
Il comportamentismo introduce un metodo sperimentale basato su due variabili:
\begin{itemize}
    \item \emph{variabile indipendente:} manipolata dallo sperimentatore, corrisponde alle stimolazioni ambientali;
    \item \emph{variabile dipendente:} misurata dallo sperimentatore, corrisponde al comportamento.
\end{itemize}
\begin{snippetdefinition}{Black box}
    Per il comportamentismo, la mente è una \emph{scatola nera}, o \emph{black box}, al cui interno lo psicologo non può né è tenuto a entrare.
\end{snippetdefinition}
L'obiettivo dello psicologo comportamentista è quindi lo studio delle associazioni \texttt{S-R}, dove \texttt{S} indica gli stimoli ambientali che impattano sulla scatola nera e \texttt{R} indica le risposte emesse dall'organismo.
Questo studio viene svolto soprattutto sugli animali. Diventa centrale il modello sperimentale fondato su una variabile indipendente, cioè lo stimolo di cui lo sperimentatore ha pieno controllo, e su una variabile dipendente, cioè la risposta comportamentale. In questo modo si osserva una relazione di causa-effetto.
Secondo questa corrente, gli stimoli ambientali impattano sulla scatola nera determinando specifiche risposte. I comportamentisti studiano quindi le associazioni stimolo-risposta.
Il comportamentismo raggiunge il suo culmine con il \emph{condizionamento operante} di Skinner.
\begin{snippetexample}{Problem box di Thorndike}
    Thorndike utilizza una \emph{Problem box} con gatti posti all'interno di una gabbia. Nella gabbia è presente un meccanismo che consente all'animale di uscire. In questo modo Thorndike studia una particolare forma di apprendimento: l'\emph{apprendimento per prove ed errori}. L'apprendimento è definito \emph{alla cieca}, perché l'animale, tastando casualmente l'ambiente, finisce per azionare il meccanismo che gli permette di uscire.
\end{snippetexample}
L'apprendimento alla cieca può essere spiegato attraverso la \emph{legge dell'effetto} e la \emph{legge dell'esercizio}.
\begin{snippetdefinition}{Legge dell'effetto}
    Secondo la \emph{legge dell'effetto}, l'animale impara che a una certa sua azione consegue un determinato effetto nell'ambiente.
\end{snippetdefinition}
\begin{snippetdefinition}{Legge dell'esercizio}
    La \emph{legge dell'esercizio} spiega l'importanza del numero di prove ed errori, cioè dei tentativi compiuti dall'animale.
\end{snippetdefinition}
Nel caso della Problem box, lo sperimentatore non può controllare direttamente lo stimolo: è l'animale che decide autonomamente quando intraprendere una determinata azione.
\begin{snippetdefinition}{Condizionamento classico}
    Il \emph{condizionamento classico}, proposto da Pavlov, è una forma di apprendimento in cui uno stimolo inizialmente neutro diventa capace di produrre una risposta dopo essere stato associato a uno stimolo incondizionato.
\end{snippetdefinition}
\emph{Pavlov} giunge incidentalmente alla scoperta del condizionamento classico mentre studia la digestione nel cane. Originariamente osserva che l'animale comincia a salivare quando il cibo si trova nelle immediate vicinanze e persino quando i collaboratori addetti alla somministrazione del cibo entrano nella stanza.
La salivazione è una risposta innata del cane e non può essere controllata volontariamente. A partire da queste osservazioni, Pavlov sviluppa un metodo per studiare quanto sia possibile instaurare nel comportamento dell'animale risposte a stimoli neutri.
In una condizione iniziale, alla presentazione visiva del cibo, cioè lo \emph{stimolo incondizionato}, corrisponde una risposta dell'animale, cioè la produzione di saliva. Successivamente Pavlov associa uno stimolo neutro allo stimolo incondizionato: quando lo sperimentatore entra con la ciotola, viene presentato anche un suono.
Tramite la presentazione ripetuta di questi due stimoli in continuità, l'animale forma un'associazione. Lo stimolo inizialmente neutro diventa così uno \emph{stimolo condizionato} e produce una determinata risposta, cioè la \emph{risposta condizionata}. Per questo motivo, quando viene presentato soltanto il suono senza il cibo, l'animale produce comunque saliva.
\begin{center}

    \includegraphics[width=0.75\textwidth]{media_psychology/pavlov-classical-conditioning-overview.jpg}
    \captionof{figure}{Schema del condizionamento classico di Pavlov.}
\end{center}
Le fasi del condizionamento classico sono:
\begin{enumerate}
    \item \emph{acquisizione:} si forma l'associazione tra stimolo condizionato e stimolo incondizionato, con aumento della salivazione;
    \item \emph{estinzione:} viene presentato lo stimolo condizionato senza lo stimolo incondizionato, e la salivazione tende a scomparire;
    \item \emph{riacquisizione:} vengono nuovamente presentati insieme stimolo condizionato e stimolo incondizionato, con un rapido riapprendimento della salivazione.
\end{enumerate}
Pavlov si occupa di risposte riflesse e biologiche dell'individuo, mentre Skinner sviluppa un metodo differente.
\begin{snippetdefinition}{Condizionamento operante}
    Il \emph{condizionamento operante}, elaborato da Skinner, è una forma di apprendimento in cui un comportamento aumenta o diminuisce in funzione delle conseguenze che produce nell'ambiente.
\end{snippetdefinition}
Il condizionamento operante di Skinner, sviluppato tra gli anni Venti e Cinquanta, utilizza la \emph{Skinner Box}, una gabbia in cui la somministrazione degli stimoli è controllata da macchinari. I soggetti di studio sono topi inseriti in una gabbia con leve associate a diversi rinforzi, ad esempio l'acqua.
Se il topo non preme mai una certa leva, essa può essere attivata dallo sperimentatore per mostrare all'animale cosa accade quando viene premuta. Nel pavimento della gabbia è presente anche una griglia elettrica che lo sperimentatore può azionare per ridurre un certo tipo di risposta, funzionando come una forma di punizione.
Questo esperimento introduce il concetto di \emph{rinforzo}.
\begin{snippetdefinition}{Rinforzo}
    Il \emph{rinforzo} è uno stimolo che aumenta la probabilità che un certo comportamento si manifesti nuovamente.
\end{snippetdefinition}
Il condizionamento operante si applica a qualsiasi tipo di risposta, anche complessa, perché ciascuna risposta può essere seguita da un rinforzo.
\emph{Tolman} sostiene invece che si può apprendere anche senza rinforzi, per far fronte a una situazione problematica. Questa forma di apprendimento prende il nome di \emph{apprendimento latente}.
\begin{snippetdefinition}{Apprendimento latente}
    L'\emph{apprendimento latente} è una forma di apprendimento che avviene anche in assenza di rinforzo e che può manifestarsi solo in un momento successivo, quando diventa utile per affrontare una situazione problematica.
\end{snippetdefinition}
Tolman arriva a questa conclusione attraverso studi con topi posti all'interno di un labirinto, in cui devono trovare l'uscita. I topi vengono divisi in tre gruppi:
\begin{enumerate}
    \item nessun rinforzo all'uscita, ad esempio nessun cibo: si osserva una debole riduzione degli errori;
    \item rinforzo all'uscita: si osserva una rapida riduzione degli errori;
    \item nessun rinforzo per dieci giorni e introduzione del rinforzo dall'undicesimo giorno: si osserva una drastica riduzione degli errori all'undicesimo giorno.
\end{enumerate}
Per Tolman, il rinforzo è utile perché il comportamento si manifesti, ma non è necessario affinché venga appreso.
Tolman introduce anche il concetto di \emph{mappa spaziale}, o \emph{mappa cognitiva}. Egli osserva che i topi costruiscono una rappresentazione spaziale del labirinto: quando vengono introdotti blocchi che modificano il percorso, i topi scelgono una strada piuttosto che un'altra.
\subsection{Il cognitivismo}
Lo sviluppo dell'orientamento cognitivista rappresenta lo sbocco naturale di una serie di ricerche che portano al superamento del comportamentismo.
\begin{snippetdefinition}{Cognitivismo}
    Il \emph{cognitivismo} è un approccio psicologico che studia i processi attraverso cui il sistema cognitivo acquisisce, elabora, archivia e recupera le informazioni.
\end{snippetdefinition}
In \emph{Verbal Behavior}, Skinner afferma che l'acquisizione del linguaggio da parte dei bambini avviene attraverso ordinari processi di apprendimento. Il linguista \emph{Noam Chomsky} si oppone a questa posizione parlando di \emph{competenza innata} nell'acquisizione del linguaggio, contribuendo alla nascita del cognitivismo.
Secondo Chomsky, la teoria di Skinner riduce il linguaggio a un insieme di apprendimenti stimolo-risposta, riuscendo a spiegare solo una semplice ripetizione linguistica, ma non l'infinita possibilità combinatoria del linguaggio.
\subsection{La scienza cognitiva}
La \emph{scienza cognitiva} ha origine negli anni Settanta.
I due paradigmi principali della scienza della cognizione sono:
\begin{itemize}
    \item \emph{modularismo:} prospettiva secondo cui la mente è organizzata in moduli specializzati;
    \item \emph{connessionismo:} prospettiva che mette in relazione l'architettura biologica del cervello con l'architettura funzionale dell'attività cognitiva.
\end{itemize}
I modelli del funzionamento mentale della scienza cognitiva possono essere divisi in \emph{modelli computer-style} e \emph{modelli brain-style}. I primi si fondano sull'idea che, per spiegare l'elaborazione dell'informazione da parte della mente umana, occorra considerare le modalità di funzionamento del computer. I secondi partono dal presupposto che, per formulare ipotesi sulla mente umana, si debba fare riferimento al modo di operare del cervello umano.
\subsubsection{Il modularismo e la mente computazionale}
Il modularismo propone un'analogia tra computer e mente umana. Le due entità condividono due abilità basilari: la capacità di usare simboli e la capacità di eseguire calcoli e operazioni.
\emph{Fodor} elabora una visione modulare della mente.
\begin{snippetdefinition}{Modulo}
    Un \emph{modulo} è una struttura specializzata della mente, paragonabile a un cassetto, che funziona come sistema esperto in uno specifico ambito dell'interazione con l'ambiente.
\end{snippetdefinition}
Secondo Fodor:
\begin{itemize}
    \item la mente è organizzata in moduli, ciascuno dotato di una struttura specializzata;
    \item il modularismo postula una \emph{mente adattata}, cioè una mente fondata su una natura umana specie-specifica, universale, fissa e biologicamente determinata.
\end{itemize}
Fodor ipotizza che il sistema cognitivo sia costituito da tre tipi di strutture gerarchicamente organizzate:
\begin{itemize}
    \item i \emph{trasduttori};
    \item i \emph{sistemi di input};
    \item i \emph{sistemi centrali}.
\end{itemize}
I trasduttori registrano informazioni sul mondo e producono rappresentazioni degli stati ambientali in forme tali da renderle disponibili per l'elaborazione successiva. Il prodotto dei trasduttori, ad esempio un impulso neurale, deve essere elaborato dai sistemi di input: i cinque sistemi percettivi più il linguaggio.
I sistemi di input non sono sotto il controllo cosciente dell'individuo e sono definiti \emph{incapsulati}, perché non sono influenzati né dagli altri sistemi di input né dai processi centrali.
I processi centrali, invece, non sono incapsulati, poiché interagiscono tra loro e cercano di produrre una rappresentazione del mondo. Il loro compito è collegare e integrare tra loro i prodotti dei singoli sistemi di input.
\subsubsection{Il connessionismo e la mente situata}
Oggi, più che parlare di mente adattata, è corretto parlare di una mente che sa adattarsi. Questa è una delle intuizioni del \emph{connessionismo}, sorto all'inizio degli anni Ottanta.
\begin{snippetdefinition}{Connessionismo}
    Il \emph{connessionismo} è una prospettiva che mette in relazione l'architettura biologica del cervello con l'architettura funzionale dell'attività cognitiva.
\end{snippetdefinition}
Il connessionismo fa riferimento alle \emph{reti neurali artificiali}, cioè modelli ispirati alla struttura neurale del cervello. Si tratta di simulazioni che riproducono in modo semplificato le proprietà e i processi di funzionamento del sistema nervoso.
Secondo il connessionismo, le informazioni sono elaborate all'interno di reti composte da molte unità che procedono in modo parallelo, anziché seriale come accade in un computer. La conoscenza avviene quindi sulla base di meccanismi di associazione.
Questa impostazione pone le basi per una concezione dinamica e attiva della mente, capace di adattarsi di volta in volta alle condizioni del momento e di autocorreggersi. Si parla perciò di \emph{adapting mind}.
\begin{center}
    \includegraphics[width=0.58\textwidth]{media_psychology/connectionism-neural-networks.jpg}
    \captionof{figure}{Schema del connessionismo e delle reti neurali artificiali.}
\end{center}
\begin{snippetdefinition}{Mente situata}
    La \emph{mente situata} è una mente costantemente immersa in un contesto immediato, dal quale dipendono le sue attività cognitive.
\end{snippetdefinition}
L'importanza attribuita al contesto rimanda a una conoscenza fondata sulle pratiche della vita quotidiana. Di conseguenza, il significato di un oggetto non è una verità conoscibile a priori, ma dipende radicalmente dal suo contesto immediato.
\subsubsection{La mente radicata}
Una mente situata è anche una mente radicata nel corpo, cioè una \emph{embodied mind}. Secondo questa prospettiva, ogni conoscenza è fondata nell'esperienza, cioè nella \emph{grounded cognition}, e deriva dalle informazioni tratte dai diversi sistemi sensoriali, immaginativi, linguistici e affettivi, oltre che dalle azioni compiute dal proprio organismo in una determinata circostanza.
In particolare, le rappresentazioni mentali di un oggetto o di un evento consistono nell'elaborazione di mappe della situazione che consentono di effettuare simulazioni.
\begin{snippetdefinition}{Mente simulativa}
    La \emph{mente simulativa} è una mente capace di elaborare rappresentazioni che permettono di riprodurre, prevedere e immaginare fenomeni, oggetti, eventi e mondi possibili.
\end{snippetdefinition}
Entrano quindi in gioco:
\begin{itemize}
    \item la capacità di riprodurre e prevedere un fenomeno, un oggetto o un evento;
    \item la capacità di creare nuovi mondi possibili a partire dagli elementi conoscitivi disponibili.
\end{itemize}

\subsection{La ricerca in psicologia}
Il \emph{metodo scientifico} in psicologia è fondamentale per affrontare i molteplici aspetti della vita quotidiana, allontanandosi dalle credenze della cosiddetta \emph{psicologia ingenua} o \emph{psicologia del senso comune}.
\begin{snippetdefinition}{Psicologia ingenua}
    La \emph{psicologia ingenua}, o \emph{psicologia del senso comune}, spiega il comportamento attraverso teorie fondate sull'esperienza personale e non sul controllo scientifico. Un esempio è l'idea secondo cui i bambini, dopo aver mangiato molti zuccheri, diventano iperattivi.
\end{snippetdefinition}
Il processo di ricerca in psicologia risponde all'esigenza di determinare quali supposizioni sul comportamento siano corrette, verificarle e approfondirle.
\begin{center}

    \includegraphics[width=0.75\textwidth]{media_psychology/psychology-research-process.jpg}
    \captionof{figure}{Schema generale del processo di ricerca scientifica in psicologia.}
\end{center}
Le domande dei ricercatori nascono da osservazioni, credenze, aspetti culturali e idee formulate in precedenza. Spesso i ricercatori combinano in modo originale conoscenze già disponibili, offrendo una diversa prospettiva di interpretazione.
\begin{snippetdefinition}{Teoria}
    Una \emph{teoria} è un insieme organizzato di proposizioni che spiegano un fenomeno o un insieme di fenomeni. Essa costituisce il contesto fondamentale entro cui formulare le domande che guidano la ricerca.
\end{snippetdefinition}
Quando viene formulata una teoria, ci si aspetta che essa tenga conto di quanto già si conosce sull'argomento e che sia in grado di generare nuove ipotesi.
\begin{snippetdefinition}{Ipotesi}
    Un'\emph{ipotesi} è un'affermazione provvisoria e verificabile sulla relazione tra le cause e le conseguenze di un fenomeno. Per essere verificata, deve essere formulata in modo tale da poter essere falsificata.
\end{snippetdefinition}
Le ipotesi sono spesso espresse nella forma ``se... allora...'': esse prevedono determinati esiti, cioè un \emph{conseguente}, a partire da specifiche condizioni, cioè un \emph{antecedente}.
\begin{snippetexample}{Ipotesi sugli zuccheri}
    Se i bambini consumano molti zuccheri, allora è probabile che manifestino comportamenti iperattivi. Per verificare questa ipotesi, si potrebbero far consumare zuccheri ai bambini e osservare gli effetti sui loro livelli di attivazione. Se non si trovasse alcun effetto, l'ipotesi potrebbe essere respinta e considerata falsificata.
\end{snippetexample}
\begin{snippetdefinition}{Metodo scientifico}
    Il \emph{metodo scientifico} è un insieme generale di procedure per la raccolta e l'interpretazione dei dati, finalizzato a limitare le fonti di errore e a trarre conclusioni verificabili.
\end{snippetdefinition}
La psicologia viene considerata una scienza quando rispetta i vincoli posti dalle regole del metodo scientifico.
A partire dal proprio oggetto d'indagine, il ricercatore stabilisce quale metodo di ricerca utilizzare per verificare la sua ipotesi. Può, ad esempio, ricreare una situazione in laboratorio oppure svolgere una ricerca sul campo.
\begin{snippetdefinition}{Metodo di ricerca}
    Un \emph{metodo di ricerca} è una procedura scientifica che consente di raccogliere e mettere in rapporto dati empirici, in modo da poterli interpretare.
\end{snippetdefinition}
Le ricerche psicologiche si basano su dati che i ricercatori raccolgono secondo modalità diverse in funzione degli obiettivi dell'indagine.
\begin{snippetdefinition}{Dati grezzi}
    I \emph{dati grezzi} sono informazioni raccolte attraverso procedure prestabilite, cioè protocolli sperimentali, e possono essere usati come prove a sostegno delle ipotesi.
\end{snippetdefinition}
Esistono due tipi principali di dati:
\begin{itemize}
    \item \emph{dati comportamentali:} collegati alla rilevazione del comportamento umano o animale;
    \item \emph{dati self-report:} descrizioni più o meno libere che le persone forniscono riguardo a ciò che accade nella loro mente.
\end{itemize}
Una volta raccolti i dati, i ricercatori procedono alla loro analisi e alla successiva accettazione o falsificazione dell'ipotesi.
L'\emph{analisi statistica} contribuisce a verificare se le ipotesi sperimentali sono corrette. Gli psicologi si servono soprattutto di due tipologie di statistiche:
\begin{itemize}
    \item \emph{statistiche descrittive:} procedure matematiche per descrivere e sintetizzare dati numerici;
    \item \emph{statistiche inferenziali:} procedure fondate sulle teorie probabilistiche per formulare previsioni e giudizi.
\end{itemize}
Se i ricercatori ritengono che i dati possano avere una ricaduta scientifica, si passa alla divulgazione dei risultati tramite l'invio di un articolo a una rivista per la pubblicazione. Nell'articolo è necessario riportare sia la procedura utilizzata sia il metodo di analisi dei dati. I dati devono essere conservati e resi accessibili a un'eventuale verifica pubblica.
Dopo la pubblicazione, la comunità scientifica esamina il lavoro e identifica eventuali questioni irrisolte.
\subsection{I metodi della psicologia}
La psicologia si fonda principalmente su tre metodi:
\begin{itemize}
    \item il metodo \emph{osservativo};
    \item il metodo \emph{sperimentale};
    \item il metodo \emph{correlazionale}.
\end{itemize}
\subsubsection{L'osservazione}
L'\emph{osservazione} è una delle modalità principali per studiare le persone. Essa si focalizza sia sul processo sia sull'esito del comportamento.
Esistono due tipi principali di osservazione:
\begin{itemize}
    \item \emph{osservazione naturalistica};
    \item \emph{osservazione clinica}.
\end{itemize}
\begin{snippetdefinition}{Osservazione naturalistica}
    Nell'\emph{osservazione naturalistica}, il ricercatore osserva un comportamento che accade in natura, senza cercare di cambiarlo o influenzarlo. Questo metodo permette di scoprire l'estensione di un fenomeno o di individuare le variabili e le relazioni più significative.
\end{snippetdefinition}
\begin{snippetdefinition}{Osservazione clinica}
    Nell'\emph{osservazione clinica}, l'osservatore interagisce con la persona osservata cercando di ottenere informazioni attraverso domande opportune, evitando che esse siano percepite come disturbanti. Questo metodo viene usato soprattutto a scopi diagnostici e psicoterapeutici.
\end{snippetdefinition}
\subsubsection{Metodo sperimentale}
Nella vita quotidiana le persone possono ipotizzare cause diverse all'origine di uno stesso evento. Gli psicologi affrontano lo stesso problema quando cercano di affermare qualcosa di certo su un particolare fenomeno.
\begin{snippetdefinition}{Ambiguità causale}
    L'\emph{ambiguità causale} si verifica quando esistono più cause possibili per un determinato effetto.
\end{snippetdefinition}
Per ovviare all'ambiguità causale, si utilizza il \emph{metodo sperimentale}, che permette di formulare leggi scientifiche sulle relazioni causali tra variabili. In particolare, viene manipolata la variabile indipendente per verificarne l'effetto sulla variabile dipendente.
\begin{snippetdefinition}{Variabile}
    Una \emph{variabile} è qualsiasi caratteristica che può assumere diversi valori, quantitativi o qualitativi.
\end{snippetdefinition}
Le conclusioni di un esperimento sono accurate solo se non sono influenzate da distorsioni, o \emph{bias}. Questi bias possono riguardare:
\begin{itemize}
    \item il ricercatore;
    \item i soggetti sperimentali;
    \item il disegno di ricerca.
\end{itemize}
\begin{snippetdefinition}{Variabile indipendente}
    La \emph{variabile indipendente} è la variabile manipolata dallo sperimentatore.
\end{snippetdefinition}
\begin{snippetdefinition}{Variabile dipendente}
    La \emph{variabile dipendente} è la variabile osservata e misurata dallo sperimentatore. Il suo valore dovrebbe dipendere dai cambiamenti indotti dalla variabile indipendente.
\end{snippetdefinition}
\begin{snippetexample}{Colore della stanza e creatività}
    Si vuole indagare se il colore di una stanza influenza la creatività dei partecipanti. La \emph{variabile indipendente} è il colore della stanza, ad esempio blu o rosso, mentre la \emph{variabile dipendente} è il punteggio ottenuto in un test di creatività.
\end{snippetexample}
\begin{center}
    \includegraphics[width=0.65\textwidth]{media_psychology/independent-dependent-confounding-variables.jpg}
    \captionof{figure}{Schema della relazione tra variabile indipendente, variabile dipendente e variabili confondenti.}
\end{center}
Per operare in condizioni controllate, costanti e uniformi è necessario ricorrere alla \emph{standardizzazione}.
\begin{snippetdefinition}{Standardizzazione delle procedure}
    La \emph{standardizzazione delle procedure} consiste nell'utilizzo di procedure uniformi e costanti in ogni fase della sperimentazione e per ogni partecipante.
\end{snippetdefinition}
\begin{snippetdefinition}{Standardizzazione delle norme}
    La \emph{standardizzazione delle norme} consiste nell'uniformità nell'attribuzione dei punteggi e nella valutazione dei dati.
\end{snippetdefinition}
La standardizzazione del significato di un concetto implica darne una \emph{definizione operazionale}.
\begin{snippetdefinition}{Definizione operazionale}
    Una \emph{definizione operazionale} formula un concetto nei termini delle procedure utilizzate per determinarlo.
\end{snippetdefinition}
\begin{snippetexample}{Definizione operazionale di bambini}
    In un esperimento sugli effetti degli zuccheri, non è possibile studiare tutti i bambini del mondo. Una definizione operazionale di \emph{bambini} potrebbe allora essere: gli alunni che frequentano la classe terza di una scuola primaria.
\end{snippetexample}
Negli esperimenti, la confusione delle variabili è una delle maggiori minacce alla validità. Quando il comportamento dei partecipanti viene modificato da una o più variabili diverse da quella introdotta dallo sperimentatore, si crea confusione nell'interpretazione dei dati. Questi elementi non previsti sono chiamati \emph{variabili confondenti}.
\begin{snippetdefinition}{Variabile confondente}
    Una \emph{variabile confondente} è una variabile non prevista che modifica il comportamento dei partecipanti e rende ambigua l'interpretazione dei risultati.
\end{snippetdefinition}
Esistono due principali tipi di variabile confondente:
\begin{itemize}
    \item \emph{effetto dell'aspettativa}, o \emph{effetto Rosenthal};
    \item \emph{effetto placebo}.
\end{itemize}
\begin{snippetdefinition}{Effetto dell'aspettativa}
    L'\emph{effetto dell'aspettativa}, o \emph{effetto Rosenthal}, consiste nella distorsione dei risultati provocata dalle aspettative che lo sperimentatore o i soggetti sperimentali hanno riguardo ai risultati stessi.
\end{snippetdefinition}
\begin{snippetdefinition}{Effetto placebo}
    L'\emph{effetto placebo} si verifica quando i partecipanti a un esperimento modificano le proprie risposte in assenza di qualunque manipolazione sperimentale.
\end{snippetdefinition}
I comportamenti degli esseri umani e degli animali sono complessi e spesso dovuti a molteplici cause. Per questa ragione, un buon disegno di ricerca deve anticipare le possibili variabili confondenti e progettare strategie per eliminarle.
\begin{snippetdefinition}{Procedure di controllo}
    Le \emph{procedure di controllo} sono procedure finalizzate a mantenere costanti tutte le variabili e le condizioni non legate all'ipotesi da verificare.
\end{snippetdefinition}
In un esperimento, le istruzioni, la temperatura dell'ambiente, il tempo concesso e molti altri dettagli della situazione sperimentale devono essere uguali per tutti i partecipanti, così che l'esperienza sperimentale sia la stessa per ciascuno.
Per eliminare le distorsioni prodotte dall'effetto dell'aspettativa, i partecipanti dovrebbero essere inconsapevoli delle condizioni sperimentali in cui si trovano.
\begin{snippetdefinition}{Controllo a singolo cieco}
    Il \emph{controllo a singolo cieco} si ha quando i partecipanti non sanno a quale condizione sperimentale sono stati assegnati.
\end{snippetdefinition}
\begin{snippetdefinition}{Controllo a doppio cieco}
    Il \emph{controllo a doppio cieco} si ha quando né i partecipanti né lo sperimentatore sanno a quale condizione sperimentale sono stati assegnati i soggetti esaminati.
\end{snippetdefinition}
Per controllare l'effetto placebo, i ricercatori inseriscono una condizione sperimentale in cui non viene somministrato il trattamento oppure non vi è manipolazione della variabile indipendente. Questa tecnica è chiamata \emph{controllo placebo}.
\begin{snippetdefinition}{Disegno di ricerca}
    Un \emph{disegno di ricerca} è una procedura scientifica che consente un'interpretazione non ambigua dei risultati, escludendo a priori interpretazioni alternative dovute a possibili variabili confondenti.
\end{snippetdefinition}
I principali tipi di disegno di ricerca sono:
\begin{itemize}
    \item il \emph{disegno between-subjects}, o \emph{tra i soggetti};
    \item il \emph{disegno within-subjects}, o \emph{entro i soggetti}.
\end{itemize}
\begin{snippetdefinition}{Disegno tra soggetti}
    Il \emph{disegno tra soggetti} prevede che i partecipanti siano assegnati casualmente alla condizione sperimentale o alla condizione di controllo, cioè all'assenza di trattamento.
\end{snippetdefinition}
\begin{snippetexample}{Disegno tra soggetti}
    In uno studio sull'effetto del colore della stanza sulla creatività, i partecipanti vengono divisi in due gruppi indipendenti. Il primo gruppo completa un test di creatività in una stanza blu, mentre il secondo completa lo stesso test in una stanza rossa. I risultati dei due gruppi vengono poi confrontati per stabilire se il colore della stanza abbia avuto un effetto sui punteggi di creatività.
\end{snippetexample}
\begin{snippetdefinition}{Assegnazione casuale}
    L'\emph{assegnazione casuale} è una strategia che attribuisce a ogni partecipante la stessa probabilità di trovarsi in una o nell'altra condizione sperimentale.
\end{snippetdefinition}
Grazie all'assegnazione casuale si possono ottenere gruppi sperimentali e gruppi di controllo, in base all'esposizione o meno alle condizioni della variabile indipendente.
\begin{snippetdefinition}{Disegno entro i soggetti}
    Il \emph{disegno entro i soggetti} è un disegno sperimentale che utilizza ogni partecipante come controllo di se stesso.
\end{snippetdefinition}
\begin{snippetexample}{Disegno entro i soggetti}
    Il comportamento di un partecipante prima del trattamento può essere confrontato con il suo comportamento dopo il trattamento. In alternativa, ogni partecipante può essere sottoposto a più livelli della variabile indipendente.
\end{snippetexample}
Il metodo sperimentale permette di trarre conclusioni solide sulle relazioni causali tra variabili, ma non è sempre il metodo migliore da utilizzare. Spesso il comportamento è studiato in un ambiente artificiale, in cui i fattori situazionali sono controllati in modo così drastico da poter alterare il corso naturale del comportamento.
Inoltre, i partecipanti alla ricerca di solito sanno di trovarsi in un esperimento e di essere misurati. Possono quindi cercare di compiacere lo sperimentatore o di indovinare lo scopo della ricerca, modificando il proprio comportamento rispetto a quanto accadrebbe se non sapessero di essere osservati.
\begin{snippetdefinition}{Quasi esperimento}
    Si parla di \emph{quasi esperimento} quando esiste un'ipotesi causale da verificare, è possibile manipolare la variabile indipendente e si possono confrontare due o più condizioni sperimentali, ma non sono soddisfatti tutti i requisiti di un vero esperimento.
\end{snippetdefinition}
Le inferenze basate sui dati di un quasi esperimento sono meno affidabili, perché non tengono sotto controllo tutte le variabili confondenti. Le conclusioni risultano quindi più deboli rispetto a quelle di un vero esperimento. Questi studi sono spesso condotti in un contesto naturale o sul campo.
\subsubsection{Metodo correlazionale}
Quando uno psicologo non può manipolare facilmente una variabile, oppure non può farlo nel rispetto delle norme etiche, utilizza il \emph{metodo correlazionale}, che permette di determinare fino a che punto due variabili sono associate.
\begin{snippetdefinition}{Metodo correlazionale}
    Il \emph{metodo correlazionale} è un metodo di ricerca utilizzato per stabilire il grado di associazione tra due variabili, senza manipolarle sperimentalmente.
\end{snippetdefinition}
Per individuare il grado di correlazione lineare tra due variabili, quando esse sono misurate su scala a intervalli o a rapporti, i ricercatori calcolano il \emph{coefficiente di correlazione} \(r\).
\begin{snippetdefinition}{Coefficiente di correlazione}
    Il \emph{coefficiente di correlazione} \(r\) è una misura statistica che indica il grado di correlazione esistente tra due variabili. Il suo valore è compreso tra \(+1{,}0\) e \(-1{,}0\).
\end{snippetdefinition}
Un valore di \(+1{,}0\) indica una perfetta correlazione positiva, un valore di \(-1{,}0\) indica una perfetta correlazione negativa e un valore di \(0{,}0\) indica assenza di correlazione.
\begin{snippetdefinition}{Correlazione positiva}
    Una \emph{correlazione positiva} indica che, all'aumentare del valore di una variabile, aumenta anche il valore dell'altra.
\end{snippetdefinition}
\begin{snippetdefinition}{Correlazione negativa}
    Una \emph{correlazione negativa} indica che, all'aumentare del valore di una variabile, il valore dell'altra diminuisce.
\end{snippetdefinition}
\begin{snippetnote}{Correlazione e causalità}
    La \emph{correlazione} non implica causalità. Un'associazione tra due variabili può dipendere da molteplici nessi causa-effetto e può coinvolgere una terza variabile sottostante.
\end{snippetnote}
\begin{snippetexample}{Sonno e successo scolastico}
    Si vuole indagare la relazione tra bisogno di sonno degli studenti e successo scolastico. Il bisogno di sonno può essere operazionalizzato come quantità media di ore di sonno per notte, mentre il successo scolastico può essere definito come media dei voti pesata in funzione del numero di esami sostenuti. Una correlazione positiva indicherebbe che gli studenti che dormono di più tendono ad avere voti più alti, ma non permetterebbe di concludere che dormire di più causa direttamente voti migliori.
\end{snippetexample}
Un'elevata correlazione indica solo che due serie di dati sono associate in modo sistematico. Non implica che una variabile sia causa dell'altra. Ad esempio, gli studenti potrebbero dormire bene e ottenere buoni voti perché frequentano corsi più facili: in questo caso, la difficoltà dei corsi costituirebbe una terza variabile responsabile della correlazione positiva tra ore di sonno e media dei voti.
\subsection{La misurazione in psicologia}
Data la varietà e la complessità dei fenomeni psicologici, essi non sono sempre semplici da misurare. Alcune azioni o processi possono essere osservati con facilità, mentre altri, come l'ansia o i sogni, non sono direttamente osservabili.
\begin{snippetdefinition}{Misura}
    Per \emph{misura} si intende l'assegnazione di un numero o di un'etichetta alle caratteristiche degli oggetti, in base a determinate regole e alla grandezza assunta come campione.
\end{snippetdefinition}
La misurazione deve essere il risultato del confronto tra un dato osservato e una posizione identificabile su una scala.
Nel 1951, Stevens propose una classificazione di quattro tipi di scale di misurazione, dalla più semplice alla più complessa:
\begin{itemize}
    \item \emph{scala nominale:} classifica sulla base della presenza o dell'assenza della qualità considerata, senza fornire una graduatoria;
    \item \emph{scala ordinale:} permette graduatorie, cioè relazioni di maggiore o minore rispetto a una qualità, ma non informa sulla distanza tra gli eventi;
    \item \emph{scala a intervalli:} permette graduatorie e mantiene costante l'intervallo tra due posizioni successive lungo tutta la scala, come nella misurazione della temperatura;
    \item \emph{scala a rapporti:} possiede le proprietà della scala a intervalli e permette di identificare una posizione corrispondente allo zero assoluto, come nella misurazione del peso.
\end{itemize}
La misurazione psicologica ha come obiettivo la produzione di risultati che siano al tempo stesso \emph{attendibili} e \emph{validi}.
\begin{snippetdefinition}{Attendibilità}
    L'\emph{attendibilità} si riferisce alla coerenza o stabilità dei dati comportamentali ottenuti attraverso test psicologici. Un risultato è attendibile se, a partire da condizioni simili, si ripete in tempi differenti.
\end{snippetdefinition}
L'attendibilità misura quindi la \emph{riproducibilità} del risultato: se un risultato è attendibile, altri ricercatori dovrebbero essere in grado di replicare l'esperimento e ottenere lo stesso esito.
\begin{snippetdefinition}{Validità}
    La \emph{validità} indica il livello di precisione con cui un test o uno strumento di ricerca misura la variabile psicologica che si intende misurare.
\end{snippetdefinition}
Si possono distinguere quattro tipi principali di validità:
\begin{itemize}
    \item \emph{validità interna:} grado in cui le conclusioni di una ricerca sono riferibili al campione di individui studiati e permettono di ritenere esistente una relazione causa-effetto tra le variabili;
    \item \emph{validità esterna:} grado di generalizzazione dei risultati di una ricerca ad altri soggetti e ad altre situazioni;
    \item \emph{validità di costrutto:} accuratezza con cui i risultati di una ricerca sono conformi alla teoria che ne è alla base;
    \item \emph{validità statistica:} grado di probabilità che la relazione tra variabili sia significativa e non ottenuta per puro caso.
\end{itemize}
Un esperimento valido consente al ricercatore di generalizzare i propri risultati a contesti più ampi.
Quando si vogliono ottenere dati relativi a esperienze che non possono essere osservate direttamente, come stati psicologici interni, le ricerche fanno affidamento sui dati di autovalutazione, cioè sulle \emph{misure self-report}.
\begin{snippetdefinition}{Misure self-report}
    Le \emph{misure self-report} sono strumenti di autovalutazione, come questionari e interviste, attraverso cui i partecipanti forniscono informazioni su esperienze, stati interni, comportamenti o abitudini.
\end{snippetdefinition}
\begin{snippetdefinition}{Questionario}
    Un \emph{questionario} consiste in una serie di domande scritte che possono riguardare fatti, abitudini, stati affettivi, comportamenti passati o presenti. Può includere anche domande aperte, che permettono al soggetto di rispondere liberamente con parole proprie.
\end{snippetdefinition}
\begin{snippetdefinition}{Intervista interattiva}
    Un'\emph{intervista interattiva} è un colloquio tra ricercatore e individuo finalizzato alla raccolta di informazioni dettagliate.
\end{snippetdefinition}
\begin{snippetnote}{Limiti degli strumenti self-report}
    L'attendibilità e la validità degli strumenti self-report non sono garantite: i partecipanti possono fraintendere le domande, non ricordare chiaramente le proprie esperienze o rispondere in modo influenzato dalla desiderabilità sociale.
\end{snippetnote}
La \emph{desiderabilità sociale} può portare le persone a fornire risposte false o ingannevoli per creare un'immagine favorevole, o sfavorevole, di sé.
\begin{center}

    \includegraphics[width=0.65\textwidth]{media_psychology/anxiety-self-report-scale.jpg}
    \captionof{figure}{Esempio di scala di autovalutazione dell'ansia.}
\end{center}

\section{Basi biologiche ed evolutive del comportamento}
In questo capitolo si considerano gli aspetti biologici dell'individualità, per comprendere che cosa rende ogni individuo diverso dagli altri esseri umani e in quale misura l'ereditarietà plasma la vita delle persone.
\subsection{Ereditarietà e comportamento}
Un aspetto importante della spiegazione psicologica è costituito dalla contrapposizione tra \emph{natura} e \emph{cultura}, oppure tra \emph{ereditarietà} e \emph{ambiente}.
\begin{snippetexample}{Comportamento aggressivo}
    Consideriamo il comportamento aggressivo. Si potrebbe ipotizzare che alcuni individui siano aggressivi a causa di un aspetto del loro corredo biologico, ad esempio perché hanno ereditato una certa propensione alla violenza da uno dei genitori. In alternativa, si potrebbe pensare che tutti gli esseri umani siano inclini alla violenza nella stessa misura, oppure che il grado di aggressività aumenti in risposta alle caratteristiche dell'ambiente in cui crescono.
\end{snippetexample}
Poiché le caratteristiche ambientali sono più facilmente osservabili, è spesso più semplice comprendere il modo in cui il comportamento delle persone ne è influenzato. Le forze biologiche che incidono sul comportamento, invece, non sono sempre direttamente evidenti.
\emph{Charles Darwin}, ne \emph{L'origine delle specie} del 1859, espone una delle teorie più importanti della scienza: l'evoluzione della vita. Riflettendo sulle specie animali incontrate durante il suo viaggio, Darwin suggerisce che la varietà delle specie sia il risultato di un processo chiamato \emph{selezione naturale}.
\begin{snippetdefinition}{Selezione naturale}
    La \emph{selezione naturale} è il processo secondo cui l'adattamento alle caratteristiche dell'ambiente permette ad alcuni membri di una specie di riprodursi con più successo rispetto ad altri, generando così individui con tratti favorevoli alla sopravvivenza.
\end{snippetdefinition}
Per comprendere la selezione naturale è necessario introdurre alcuni concetti propri della teoria dell'evoluzione.
\begin{snippetdefinition}{Genotipo}
    Il \emph{genotipo} è la struttura genetica ereditata da un individuo al momento del concepimento.
\end{snippetdefinition}
\begin{snippetdefinition}{Fenotipo}
    Il \emph{fenotipo} è l'insieme delle caratteristiche osservabili di un individuo, cioè il suo aspetto esteriore e il repertorio dei suoi comportamenti.
\end{snippetdefinition}
\begin{snippetexample}{Giraffa e vantaggio selettivo}
    Una giraffa eredita dai genitori un certo genotipo. All'interno di un particolare ambiente, questo genotipo contribuisce allo sviluppo del fenotipo, ad esempio un collo lungo e la capacità di strappare foglie poste molto in alto. Se le foglie basse non sono sufficienti per nutrire l'intera popolazione, una giraffa dal collo lungo possiede un \emph{vantaggio selettivo}.
\end{snippetexample}
\begin{snippetdefinition}{Vantaggio selettivo}
    Un \emph{vantaggio selettivo} è una caratteristica che aumenta la probabilità di sopravvivenza e riproduzione di un individuo in un determinato ambiente.
\end{snippetdefinition}
\begin{center}
    \includegraphics[width=0.58\textwidth]{media_psychology/natural-selection-process.jpg}
    \captionof{figure}{Schema della selezione naturale.}
\end{center}
Nell'evoluzione della specie umana, la selezione naturale ha favorito due principali tipi di adattamento: il \emph{bipedismo} e l'\emph{encefalizzazione}. Insieme, questi adattamenti hanno reso possibile l'ascesa della civiltà umana.
Con l'aumentare delle dimensioni del cervello, i nostri progenitori divennero più intelligenti e svilupparono capacità complesse di pensiero, ragionamento, ricordo e pianificazione.
Dopo il bipedismo e l'encefalizzazione, un'altra tappa evolutiva fondamentale per la specie umana è stata la comparsa del \emph{linguaggio}, che sta alla base dell'evoluzione culturale.
\subsection{Ereditarietà e comportamento}
Le condizioni in cui avviene l'evoluzione degli esseri umani favoriscono lo sviluppo di un importante potenziale biologico condiviso, al cui interno rimangono tuttavia considerevoli variazioni individuali.
\begin{snippetdefinition}{Genetica}
    La \emph{genetica} è lo studio dei geni e dei meccanismi dell'ereditarietà, cioè della trasmissione biologica dei caratteri dai genitori ai figli.
\end{snippetdefinition}
\begin{snippetdefinition}{Genetica del comportamento umano}
    La \emph{genetica del comportamento umano} è un'area di studio che unisce la componente genetica delle differenze individuali e la psicologia, con l'obiettivo di esaminare comportamenti e caratteristiche della personalità.
\end{snippetdefinition}
La genetica del comportamento umano si concentra soprattutto sulle origini delle differenze individuali, chiedendosi quali fattori dell'ereditarietà genetica contribuiscano a spiegare il modo in cui una persona pensa e si comporta.
\subsubsection{Basi di genetica}
Nel nucleo di ogni cellula si trova il materiale genetico, cioè il \emph{DNA}, organizzato in minuscole unità chiamate \emph{geni}.
\begin{snippetdefinition}{Gene}
    Un \emph{gene} è un segmento di DNA responsabile della trasmissione dei caratteri ereditari e della produzione di proteine.
\end{snippetdefinition}
Quando si considerano aspetti complessi dell'esperienza umana, è importante osservare che alla manifestazione di un particolare attributo contribuiscono più geni. Questi tratti sono chiamati \emph{caratteri poligenici}.
\begin{snippetdefinition}{Carattere poligenico}
    Un \emph{carattere poligenico} è un tratto alla cui manifestazione contribuiscono più geni.
\end{snippetdefinition}
Per comprendere le funzioni dei geni, la ricerca sulla genetica del comportamento umano si concentra spesso sulla valutazione dell'\emph{ereditabilità} di particolari caratteri e comportamenti.
\begin{snippetdefinition}{Ereditabilità}
    L'\emph{ereditabilità} indica l'influenza relativa della componente genetica, contrapposta a quella ambientale, sui pattern comportamentali. Si misura su una scala da \(0\) a \(1\), dove una stima vicina a \(1\) suggerisce che l'attributo considerato dipende in larga misura da influenze genetiche.
\end{snippetdefinition}
Per distinguere le influenze ambientali da quelle genetiche, i ricercatori utilizzano spesso gli \emph{studi sulle adozioni} e gli \emph{studi sui gemelli}.
\begin{snippetdefinition}{Studi sulle adozioni}
    Gli \emph{studi sulle adozioni} confrontano le caratteristiche di bambini cresciuti in famiglie adottive con quelle dei genitori naturali e dei genitori adottivi, per valutare il peso relativo di ereditarietà e ambiente.
\end{snippetdefinition}
\begin{snippetdefinition}{Studi sui gemelli}
    Gli \emph{studi sui gemelli} esaminano il livello di somiglianza in particolari tratti comportamentali tra gemelli monozigoti e gemelli dizigoti. Poiché i gemelli monozigoti condividono una quota genetica maggiore rispetto ai dizigoti, il confronto permette di stimare l'ereditabilità di un attributo.
\end{snippetdefinition}
I ricercatori hanno documentato che geni e ambiente interagiscono tra loro nel determinare il fenotipo.
\subsection{Il sistema nervoso in azione}
\begin{snippetdefinition}{Neuroscienza}
    La \emph{neuroscienza} è lo studio scientifico del cervello e dei collegamenti tra attività cerebrale e comportamento.
\end{snippetdefinition}
\begin{snippetdefinition}{Neurone}
    Il \emph{neurone} è una cellula specializzata nel ricevere, elaborare e trasmettere informazioni ad altre cellule del corpo.
\end{snippetdefinition}
I neuroni variano per forma, dimensione, composizione chimica e funzione svolta, ma presentano tutti una struttura fondamentale comune.
Esistono tre classi fondamentali di neuroni:
\begin{itemize}
    \item \emph{neuroni sensoriali:} trasmettono i messaggi dalle cellule dei recettori sensoriali verso il sistema nervoso centrale;
    \item \emph{neuroni motori:} portano i messaggi dal sistema nervoso centrale verso muscoli e ghiandole;
    \item \emph{interneuroni:} mettono in comunicazione tra loro altri neuroni, permettendo un'elaborazione più fine dei segnali.
\end{itemize}
\begin{snippetexample}{Riflesso di allontanamento dal dolore}
    Quando i recettori del dolore situati vicino alla pelle sono stimolati da un oggetto pungente, inviano messaggi a un interneurone presente nel midollo spinale attraverso i neuroni sensoriali. L'interneurone stimola i neuroni motori, che a loro volta attivano i muscoli della zona interessata, permettendo al corpo di allontanarsi dall'oggetto doloroso.
\end{snippetexample}
\begin{center}

    \includegraphics[width=0.65\textwidth]{media_psychology/pain-withdrawal-reflex.jpg}
    \captionof{figure}{Schema del riflesso di allontanamento dal dolore.}
\end{center}
Nel riflesso di allontanamento dal dolore, l'informazione raggiunge il cervello solo dopo che la sequenza riflessa ha avuto luogo e il corpo si è già allontanato dallo stimolo. In questi casi, la percezione del dolore tende a verificarsi dopo la reazione fisica al dolore stesso.
A metà degli anni Novanta, grazie a un esperimento condotto su macachi, venne scoperto casualmente un nuovo tipo di neurone: il \emph{neurone specchio}.
\begin{snippetdefinition}{Neurone specchio}
    Un \emph{neurone specchio} è un neurone che si attiva sia quando un individuo esegue un'azione motoria sia quando osserva un altro individuo eseguire la stessa azione.
\end{snippetdefinition}
Per ciascun neurone, la domanda fondamentale è se attivarsi oppure no, cioè se produrre o meno una risposta in un determinato momento.
I neuroni prendono questa decisione combinando le informazioni in arrivo e determinando se gli impulsi segnalano di attivarsi oppure di non attivarsi. Ogni neurone riceve infatti una combinazione di \emph{input eccitatori} e \emph{input inibitori}.
\begin{snippetdefinition}{Input eccitatorio}
    Un \emph{input eccitatorio} è un segnale che aumenta la probabilità che il neurone si attivi.
\end{snippetdefinition}
\begin{snippetdefinition}{Input inibitorio}
    Un \emph{input inibitorio} è un segnale che riduce la probabilità che il neurone si attivi.
\end{snippetdefinition}
Nei neuroni, una corretta sequenza di input eccitatori nel tempo o nello spazio può portare alla produzione di un \emph{potenziale d'azione}.
\begin{snippetdefinition}{Potenziale d'azione}
    Il \emph{potenziale d'azione} è un impulso nervoso che si attiva in un neurone, percorre l'assone e porta al rilascio di neurotrasmettitori nella sinapsi.
\end{snippetdefinition}
\begin{snippetdefinition}{Legge del tutto o del nulla}
    La \emph{legge del tutto o del nulla} è il principio secondo cui la dimensione del potenziale d'azione non è influenzata dall'aumento dell'intensità della stimolazione.
\end{snippetdefinition}
Due neuroni non entrano mai direttamente in contatto, ma sono connessi da una \emph{sinapsi}. Quando le informazioni viaggiano da un neurone all'altro attraverso lo spazio sinaptico, si parla di \emph{trasmissione sinaptica}.
\begin{snippetdefinition}{Sinapsi}
    La \emph{sinapsi} è lo spazio di comunicazione tra due neuroni, attraverso cui avviene la trasmissione delle informazioni.
\end{snippetdefinition}
\begin{snippetdefinition}{Neurotrasmettitore}
    Un \emph{neurotrasmettitore} è un messaggero chimico rilasciato da un neurone, che attraversa la sinapsi e stimola il neurone postsinaptico. Può avere effetto eccitatorio o inibitorio.
\end{snippetdefinition}
\begin{snippetexample}{Catecolamine}
    Le \emph{catecolamine} sono una classe di sostanze chimiche che comprende due importanti neurotrasmettitori: \emph{norepinefrina} e \emph{dopamina}. Entrambe svolgono un ruolo rilevante in disturbi psicologici come i disturbi d'ansia e dell'umore.
\end{snippetexample}
\subsection{Biologia e comportamento}
I neuroscienziati cercano di comprendere il funzionamento del cervello a molteplici livelli. Diversi metodi di ricerca comportano un intervento diretto sulle strutture cerebrali.
Questi metodi affondano le loro radici nel caso di \emph{Phineas Gage}, caposquadra ferroviario che nel 1848 subì un grave incidente: in seguito a un'esplosione, una barra di metallo gli attraversò la testa.
Il danno fisico fu relativamente limitato: perse la vista da un occhio e il lato sinistro del viso rimase in parte paralizzato, ma postura, movimenti e capacità di linguaggio non furono compromessi. Tuttavia, dal punto di vista psicologico, Gage divenne un altro uomo.
I cambiamenti comportamentali osservati dopo la lesione spinsero il suo medico a ipotizzare l'esistenza di basi cerebrali collegate ad aspetti della personalità e del comportamento razionale.
Nello stesso periodo, \emph{Paul Broca} studiava il ruolo del cervello in relazione al linguaggio. Egli osservò che persone con disturbi del linguaggio presentavano lesioni nella stessa area cerebrale, che prese poi il nome di \emph{area di Broca}.
\begin{snippetdefinition}{Area di Broca}
    L'\emph{area di Broca} è una regione cerebrale frontale sinistra cruciale per la produzione del linguaggio verbale.
\end{snippetdefinition}
Nello studio dei cervelli danneggiati da incidenti, i ricercatori incontrano però una difficoltà: spesso non è possibile circoscrivere con precisione la sede e l'estensione del danno.
\begin{snippetdefinition}{Stimolazione magnetica transcranica}
    La \emph{stimolazione magnetica transcranica}, anche ripetitiva, è una tecnica che utilizza impulsi di stimolazione magnetica per produrre effetti simili a lesioni temporanee e reversibili nei soggetti umani.
\end{snippetdefinition}
\begin{snippetdefinition}{Elettroencefalogramma}
    L'\emph{elettroencefalogramma} registra l'attività elettrica cerebrale tramite elettrodi posizionati sulla superficie del cuoio capelluto.
\end{snippetdefinition}
\begin{snippetdefinition}{Tomografia a emissione di positroni}
    La \emph{tomografia a emissione di positroni}, o \emph{PET}, permette di ottenere immagini dettagliate dell'attività cerebrale registrando la radioattività delle cellule durante differenti attività cognitive e comportamentali.
\end{snippetdefinition}
\begin{snippetdefinition}{Risonanza magnetica}
    L'\emph{imaging con risonanza magnetica}, o \emph{MRI}, è una tecnica di neuroimaging che sottopone il cervello a scansione attraverso campi magnetici e onde radio. È utile soprattutto per ottenere immagini chiare dei dettagli anatomici del cervello.
\end{snippetdefinition}
Il sistema nervoso è suddiviso in due parti principali:
\begin{itemize}
    \item \emph{sistema nervoso centrale}, o \emph{SNC}, formato dall'encefalo e dal midollo spinale;
    \item \emph{sistema nervoso periferico}, o \emph{SNP}, formato da tutti i neuroni che mettono in comunicazione il sistema nervoso centrale con le diverse parti del corpo.
\end{itemize}
\begin{center}
    \includegraphics[width=0.60\textwidth]{media_psychology/nervous-system-overview.jpg}
    \captionof{figure}{Schema del sistema nervoso centrale e periferico.}
\end{center}
Il compito del sistema nervoso centrale è integrare e coordinare tutte le funzioni corporee, elaborare i messaggi neuronali in entrata e inviare comandi alle diverse parti del corpo. Il sistema nervoso centrale invia e riceve messaggi neuronali attraverso il sistema nervoso periferico.
Nonostante la sua posizione di comando, il sistema nervoso centrale è isolato da ogni contatto diretto con il mondo esterno. È il sistema nervoso periferico a fornire al sistema nervoso centrale le informazioni provenienti dai recettori sensoriali e a trasmettere i comandi dal cervello agli organi e ai muscoli del corpo.
Il sistema nervoso periferico si divide in:
\begin{itemize}
    \item \emph{sistema nervoso somatico:} regola le azioni dei muscoli del corpo;
    \item \emph{sistema nervoso autonomo:} regola le funzioni del corpo che non sono controllate in modo conscio e che devono funzionare anche durante il sonno.
\end{itemize}
Il sistema nervoso autonomo si divide a sua volta in:
\begin{itemize}
    \item \emph{sistema simpatico:} regola le risposte alle situazioni di emergenza;
    \item \emph{sistema parasimpatico:} monitora le procedure di ordinaria amministrazione dell'organismo.
\end{itemize}
La componente principale del sistema nervoso centrale è l'\emph{encefalo}, costituito da tre strutture interconnesse:
\begin{itemize}
    \item \emph{tronco encefalico:} regola i processi vitali e fondamentali del corpo, come il battito cardiaco;
    \item \emph{sistema limbico:} regola il comportamento emotivo e gli impulsi motivazionali;
    \item \emph{telencefalo:} regola le funzioni cognitive ed emotive di livello superiore.
\end{itemize}
\begin{center}
    \includegraphics[width=0.60\textwidth]{media_psychology/major-brain-structures.jpg}
    \captionof{figure}{Strutture principali dell'encefalo.}
\end{center}
Alcune informazioni hanno portato i ricercatori a supporre che le funzioni dei due emisferi cerebrali differiscano. Ad esempio, Paul Broca osservò che un deficit nel linguaggio era associato a lesioni del lato sinistro del cervello, mentre danni alle stesse aree nel lato destro non producevano gli stessi effetti.
La possibilità di esaminare le differenze tra i due emisferi si manifestò per la prima volta nell'ambito dei trattamenti chirurgici dell'epilessia grave, che consistevano nel recidere il \emph{corpo calloso}.
\begin{snippetdefinition}{Corpo calloso}
    Il \emph{corpo calloso} è un fascio di fibre nervose che trasferisce le informazioni tra i due emisferi cerebrali in entrambe le direzioni.
\end{snippetdefinition}
\begin{snippetdefinition}{Pazienti split-brain}
    I \emph{pazienti split-brain} sono pazienti sottoposti a un intervento chirurgico di recisione del corpo calloso, solitamente nel trattamento di forme gravi di epilessia.
\end{snippetdefinition}
\begin{center}

    \includegraphics[width=0.65\textwidth]{media_psychology/split-brain-visual-pathways.jpg}
    \captionof{figure}{Organizzazione delle vie visive e comunicazione tra emisferi nei pazienti split-brain.}
\end{center}
I percorsi neuronali delle informazioni visive provenienti dalla parte interna di ciascun occhio si incrociano dirigendosi da un emisfero cerebrale all'altro. I percorsi che trasmettono le informazioni dalle parti esterne di ciascun occhio, invece, non si incrociano.
L'asportazione del corpo calloso impedisce alle informazioni selezionate nel campo visivo destro, che hanno raggiunto l'emisfero sinistro, di accedere all'emisfero destro. Allo stesso modo, impedisce alle informazioni selezionate nel campo visivo sinistro di accedere all'emisfero sinistro.
Poiché nella maggior parte delle persone il linguaggio è controllato dall'emisfero sinistro, soltanto questo emisfero era in grado di fornire risposte verbali ai ricercatori. L'emisfero destro, invece, non poteva rispondere verbalmente. La comunicazione con l'emisfero destro era possibile attraverso compiti manuali di riconoscimento, abbinamento o assemblaggio di oggetti.

\begin{snippetexample}{Paziente split-brain}
    Quando un paziente usa la mano sinistra per afferrare un oggetto apparso brevemente nel campo visivo sinistro, l'azione riesce perché entrambe le informazioni, visiva e tattile, sono registrate nell'emisfero destro. Tuttavia, il paziente non sa nominare l'oggetto, poiché il linguaggio è prevalentemente una funzione dell'emisfero sinistro. Se allo stesso paziente viene richiesto di svolgere il medesimo compito con la mano destra, egli non riesce ad afferrare l'oggetto corretto, perché l'informazione visiva e quella tattile sono elaborate da emisferi differenti. Tuttavia, è in grado di nominare correttamente l'oggetto che ha in mano.
\end{snippetexample}
Grazie a molteplici metodi di ricerca, oggi sappiamo che molte funzioni del linguaggio sono \emph{lateralizzate} nell'emisfero sinistro.
\begin{snippetdefinition}{Lateralizzazione}
    Una funzione è \emph{lateralizzata} quando, per la sua realizzazione, un emisfero cerebrale svolge un ruolo primario.
\end{snippetdefinition}
Non bisogna però assumere una visione rigidamente dicotomica: è vero che l'emisfero sinistro svolge un ruolo predominante nel linguaggio, ma l'emisfero destro svolge un ruolo importante per altre funzioni. È l'azione combinata dei due emisferi a dare pienezza alle esperienze.
Ciò che rende il cervello particolarmente interessante è una conseguenza della comunicazione neuronale: il suo cambiare nel corso del tempo.
\begin{snippetdefinition}{Plasticità cerebrale}
    La \emph{plasticità cerebrale} indica i cambiamenti delle prestazioni cerebrali nel corso del tempo. Può comportare la formazione di nuove sinapsi o modifiche nelle funzioni delle sinapsi già esistenti.
\end{snippetdefinition}
La plasticità dipende anche dalle esperienze di vita: sul cervello pesa infatti l'impatto di differenti ambienti e attività.
\begin{snippetdefinition}{Neurogenesi}
    La \emph{neurogenesi} è la produzione di nuove cellule cerebrali, cioè di nuovi neuroni.
\end{snippetdefinition}
\section{Sensazione e percezione}
Per noi non è possibile immaginare un organismo privo di sensi. La \emph{sensazione} e la \emph{percezione} forniscono i materiali grezzi per i processi cognitivi, ma le sensazioni non si limitano a registrare gli stimoli sensoriali.
Il materiale sensoriale viene elaborato in tempo reale da processi cognitivi complessi, che generano interpretazioni del mondo a partire dalle conoscenze precedenti e dagli stimoli esterni.
Il mondo fornisce continuamente un numero elevato di stimoli sensoriali, ma tali stimoli non contengono informazioni sufficienti a spiegare completamente le nostre percezioni.
L'esplorazione e la ricerca sono modi per fronteggiare l'eccesso di stimoli a cui siamo sottoposti. Il processo di \emph{attenzione selettiva} consente di selezionare una parte dello stimolo sensoriale per sottoporla a un'elaborazione accurata, trascurando altri aspetti dello stimolo.
\subsection{Cosa sono sensazione e percezione}
Il compito di questo capitolo è spiegare come corpo e mente attribuiscano un senso alla molteplicità di stimoli visivi, acustici e sensoriali costantemente presenti intorno a noi.
\begin{snippetdefinition}{Percezione}
    La \emph{percezione} è il processo generale di individuazione di oggetti ed eventi nell'ambiente esterno e intorno al soggetto, volto ad attribuire loro un senso, comprenderli, riconoscerli e prepararsi ad agire rispetto a essi.
\end{snippetdefinition}
\begin{snippetdefinition}{Percetto}
    Un \emph{percetto} è ciò che viene percepito, cioè il risultato fenomenico di cui si fa esperienza tramite il processo percettivo.
\end{snippetdefinition}
Il processo della percezione può essere suddiviso in tre fasi:
\begin{enumerate}
    \item \emph{sensazione};
    \item \emph{organizzazione percettiva};
    \item \emph{identificazione e riconoscimento degli oggetti}.
\end{enumerate}
\begin{center}
    \includegraphics[width=0.62\textwidth]{media_psychology/perception-process-stages.jpg}
    \captionof{figure}{Processo percettivo: sensazione, organizzazione percettiva, identificazione e riconoscimento.}
\end{center}
\begin{snippetdefinition}{Sensazione}
    La \emph{sensazione} è l'impressione soggettiva, immediata e semplice che corrisponde a una data intensità dello stimolo fisico. È il processo attraverso cui la stimolazione dei recettori sensoriali produce impulsi neuronali che riproducono le esperienze vissute fuori e dentro il corpo.
\end{snippetdefinition}
\begin{snippetdefinition}{Organizzazione percettiva}
    L'\emph{organizzazione percettiva} è la fase in cui il cervello integra i segnali raccolti dagli organi recettori, insieme alla conoscenza pregressa del mondo, per formare una rappresentazione interna di uno stimolo esterno.
\end{snippetdefinition}
Nel caso della visione, l'organizzazione percettiva permette di valutare proprietà come \emph{dimensione}, \emph{forma}, \emph{movimento} e \emph{distanza} di un oggetto.
I processi di \emph{identificazione} e \emph{riconoscimento} attribuiscono significato ai percetti. In relazione alla visione, le prime fasi della percezione rispondono alla domanda: ``che aspetto ha l'oggetto?''. Le fasi successive rispondono invece a domande come: ``che oggetto è?'' e ``qual è la funzione dell'oggetto?''.
\subsubsection{Stimoli prossimali e distali}
Le differenze tra un oggetto fisico collocato nell'ambiente e la sua immagine sulla retina sono profonde. Una differenza importante è che l'immagine retinica è bidimensionale, mentre l'ambiente è tridimensionale.
Gli psicologi distinguono quindi due differenti stimoli percettivi: lo \emph{stimolo distale} e lo \emph{stimolo prossimale}.
\begin{snippetdefinition}{Stimolo distale}
    Lo \emph{stimolo distale} è l'oggetto fisico reale collocato nell'ambiente, distante dall'osservatore.
\end{snippetdefinition}
\begin{snippetdefinition}{Stimolo prossimale}
    Lo \emph{stimolo prossimale} è l'immagine dello stimolo distale sulla retina, cioè lo stimolo vicino all'osservatore da cui il sistema percettivo ricava le informazioni.
\end{snippetdefinition}
Ciò che aspiriamo a percepire è lo stimolo distale, cioè l'oggetto reale presente nell'ambiente. Tuttavia, lo stimolo da cui otteniamo direttamente le informazioni è lo stimolo prossimale, cioè l'immagine sulla retina.
A livello percettivo, i processi di elaborazione sono volti a determinare lo stimolo distale a partire dalle informazioni contenute nello stimolo prossimale.
\subsubsection{Processi dal basso verso l'alto e dall'alto verso il basso}
I sistemi percettivi sono caratterizzati da una profonda complessità funzionale e strutturale. Sul piano fisiologico, molte aree visive del cervello inviano informazioni in uscita ad altre aree cerebrali e ricevono informazioni in entrata dalle stesse aree.
Di conseguenza, la percezione è l'esito di processi che agiscono sia \emph{dal basso verso l'alto}, cioè \emph{bottom-up}, sia \emph{dall'alto verso il basso}, cioè \emph{top-down}.
\begin{snippetdefinition}{Processi bottom-up}
    I \emph{processi bottom-up}, o processi dal basso verso l'alto, sono guidati dalle informazioni sensoriali provenienti dal mondo fisico. Essi consistono nel recepire messaggi sensoriali dall'ambiente e inviarli al cervello affinché siano elaborati.
\end{snippetdefinition}
I processi bottom-up sono processi \emph{basati sui dati}, perché il punto di partenza dell'identificazione è l'evidenza sensoriale rilevata nel contesto.
\begin{snippetdefinition}{Processi top-down}
    I \emph{processi top-down}, o processi dall'alto verso il basso, sono guidati da conoscenze, credenze, aspettative, motivazioni, obiettivi ed esperienze passate.
\end{snippetdefinition}
Quando le aspettative influenzano la percezione, entrano in gioco i processi top-down. Essi coinvolgono nella percezione del mondo le esperienze passate, la conoscenza, le motivazioni e il background culturale.
Attraverso i processi top-down, le funzioni mentali più complesse influenzano il modo in cui riconosciamo gli oggetti e comprendiamo gli eventi. In questo caso, il processo di identificazione è \emph{guidato dai concetti}, perché le rappresentazioni mentali immagazzinate in memoria influenzano l'interpretazione dei dati sensoriali.
\begin{center}
    \includegraphics[width=0.58\textwidth]{media_psychology/bottom-up-top-down-processing.jpg}
    \captionof{figure}{Schema dei processi bottom-up e top-down.}
\end{center}
\subsubsection{La psicofisica}
Le sensazioni sono eventi personali e soggettivi, dei quali ciascuno ha un'esperienza diretta. Pur essendo soggette a variazioni interindividuali e intraindividuali in funzione delle condizioni del momento, le sensazioni possono essere comunicate agli altri e comprese.
Per rispondere a domande come ``quanto zucchero bisogna mettere in una tazza di caffè prima che inizi a essere dolce?'' è necessario misurare l'intensità delle esperienze sensoriali. Questo è il compito della \emph{psicofisica}.
\begin{snippetdefinition}{Psicofisica}
    La \emph{psicofisica} è lo studio della relazione tra gli stimoli fisici e il comportamento o le esperienze mentali evocate da tali stimoli.
\end{snippetdefinition}
La figura più significativa nella storia della psicofisica è \emph{Gustav Fechner}, che coniò il termine stesso di psicofisica.
Una domanda centrale della psicofisica è: qual è l'intensità minima che uno stimolo deve avere per essere rilevato?
\begin{snippetdefinition}{Soglia assoluta}
    La \emph{soglia assoluta} della stimolazione è la quantità minima di energia fisica necessaria per produrre un'esperienza sensoriale rilevata nel \(50\%\) dei casi.
\end{snippetdefinition}
\begin{snippetexample}{Soglia assoluta}
    I ricercatori potrebbero chiedere ai partecipanti di ascoltare un suono a volume molto basso in una stanza insonorizzata. Nel corso di una serie di prove vengono presentati stimoli di diversa intensità e, in ciascuna prova, i partecipanti indicano se sono riusciti a rilevarli.
\end{snippetexample}
\begin{center}

    \includegraphics[width=0.62\textwidth]{media_psychology/psychophysical-detection-function.jpg}
    \captionof{figure}{Esempio di funzione psicofisica per la rilevazione di uno stimolo.}
\end{center}
Non esiste una singola soglia assoluta reale, perché la sensibilità dei partecipanti cambia ogni volta che provano a rilevare uno stimolo e talvolta essi rispondono anche in assenza di stimolo, producendo un \emph{falso allarme}.
Sebbene sia possibile identificare una soglia assoluta, i sistemi sensoriali sono più sensibili alla presenza di cambiamenti in un dato ambiente sensoriale che a contesti immutati.
\begin{snippetdefinition}{Abituazione sensoriale}
    L'\emph{abituazione sensoriale} è la responsività diminuita di un sistema sensoriale sottoposto a una stimolazione prolungata.
\end{snippetdefinition}
I sistemi sensoriali si sono evoluti in modo tale da favorire la percezione di stimoli ambientali nuovi rispetto a stimoli già presenti. Per questo, ad esempio, il sole può sembrare meno abbagliante dopo un po' di tempo trascorso all'aperto. L'esperienza sensoriale più soggetta ad abituazione è quella olfattiva.
Finora si è assunto che tutte le persone esprimano giudizi in modo simile. In realtà, le misurazioni della soglia possono essere influenzate da \emph{distorsioni di giudizio}.
\begin{snippetdefinition}{Distorsioni di giudizio}
    Le \emph{distorsioni di giudizio}, o \emph{bias}, sono tendenze sistematiche dell'osservatore a rispondere in un modo particolare per cause non inerenti alle caratteristiche sensoriali dello stimolo.
\end{snippetdefinition}
\begin{snippetdefinition}{Teoria della detezione del segnale}
    La \emph{teoria della detezione del segnale}, o \emph{TDS}, è un approccio sistematico al problema delle distorsioni di giudizio. Essa permette di distinguere il ruolo degli stimoli sensoriali e quello del criterio individuale nei processi di giudizio circa la presenza o l'assenza di uno stimolo.
\end{snippetdefinition}
Nella struttura base della teoria della detezione del segnale, in una parte delle prove viene presentato uno stimolo di debole intensità, mentre in un'altra parte non viene presentato nessuno stimolo. In ogni prova, gli osservatori devono rispondere \emph{sì} se pensano che lo stimolo sia presente e \emph{no} se pensano che non lo sia.
\begin{center}

    \includegraphics[width=0.62\textwidth]{media_psychology/signal-detection-theory-outcomes.jpg}
    \captionof{figure}{Schema delle possibili risposte nella teoria della detezione del segnale.}
\end{center}
Le risposte possibili sono quattro:
\begin{itemize}
    \item \emph{hit:} il soggetto indica che il segnale è presente quando esso è effettivamente comparso;
    \item \emph{omissione:} il soggetto indica che il segnale non è presente quando esso è effettivamente comparso;
    \item \emph{falso allarme:} il soggetto indica che il segnale è presente quando esso non è comparso;
    \item \emph{rifiuto corretto:} il soggetto indica che il segnale è assente quando esso non è effettivamente comparso.
\end{itemize}
\begin{snippetdefinition}{Soglia differenziale}
    La \emph{soglia differenziale} è la più piccola differenza fisica tra due stimoli che può essere riconosciuta come differenza minima rilevabile nel \(50\%\) dei casi.
\end{snippetdefinition}
\begin{snippetexample}{Soglia differenziale}
    Una compagnia di bibite vuole produrre una cola percepita come più dolce rispetto a quelle già presenti sul mercato, ma riducendo al minimo la quantità di zucchero utilizzata. In questo caso è necessario misurare la soglia differenziale, cioè stabilire quale sia la minima differenza nella quantità di zucchero che può essere percepita.
\end{snippetexample}
Per misurare una soglia differenziale si utilizzano coppie di stimoli e si chiede ai soggetti se i due stimoli sono percepiti come differenti.
\begin{snippetdefinition}{Differenza minima rilevabile}
    La \emph{differenza minima rilevabile} è un'unità di grandezza quantitativa che misura la grandezza della differenza psicologica tra due sensazioni.
\end{snippetdefinition}
\emph{Ernst Weber} realizzò alcuni studi sulla soglia differenziale e giunse a una relazione oggi chiamata \emph{legge di Weber}.
\begin{snippetdefinition}{Legge di Weber}
    La \emph{legge di Weber} afferma che la differenza minima rilevabile tra due stimoli è una frazione costante dell'intensità dello stimolo standard.
\end{snippetdefinition}
\[
    \frac{\Delta I}{I} = K
\]
dove \(K\) indica la costante di Weber, \(I\) l'intensità dello stimolo standard, mentre \(\Delta I\) è la misura dell'incrmeento che produce una differenza minima rilevabile.
\begin{snippetexample}{Legge di Weber}
    In un esperimento si può chiedere ai partecipanti di indicare, tra una coppia di linee, se le due linee hanno la stessa lunghezza. Weber ha osservato che ogni modalità sensoriale presenta un valore specifico di \(K\); ad esempio, per la lunghezza la costante \(K\) è pari a \(0{,}1\).
\end{snippetexample}
\begin{center}

    \includegraphics[width=0.62\textwidth]{media_psychology/webers-law-perceptual-comparison.jpg}
    \captionof{figure}{Esempio di confronto percettivo tra stimoli nella legge di Weber.}
\end{center}
La legge di Weber fornisce una buona approssimazione, ma non perfetta, del modo in cui la differenza minima rilevabile aumenta con l'intensità dello stimolo standard. I maggiori problemi di generalizzazione della legge si riscontrano quando l'intensità dello stimolo diventa molto elevata.
Nel 1860 Gustav Fechner osservò che la sensazione è direttamente proporzionale al logaritmo dell'intensità dello stimolo. In altri termini, mentre lo stimolo aumenta secondo una progressione geometrica, la sensazione corrispondente aumenta secondo una progressione aritmetica.
\[
    S = K \log I + C
\]

Nella formula di Fechner:
\begin{itemize}
    \item \(S\) indica l'\emph{intensità della sensazione};
    \item \(K\) indica la \emph{costante di Weber};
    \item \(\log I\) indica il \emph{logaritmo dell'intensità dello stimolo};
    \item \(C\) indica la \emph{costante di integrazione}.
\end{itemize}
\begin{center}

    \includegraphics[width=0.62\textwidth]{media_psychology/fechner-logarithmic-law.jpg}
    \captionof{figure}{Relazione logaritmica tra intensità dello stimolo e ampiezza della sensazione secondo Fechner.}
\end{center}
\subsection{Processi dal basso verso l'alto: l'organizzazione percettiva}
Sarebbe difficile orientarsi nel mondo se non fossimo in grado di organizzare e integrare le informazioni disponibili prodotte dall'attività di milioni di recettori.
\begin{snippetdefinition}{Organizzazione percettiva}
    I \emph{processi di organizzazione percettiva} sono processi di integrazione delle informazioni sensoriali che permettono di costruire una percezione coerente.
\end{snippetdefinition}
Questi processi sono stati descritti per la prima volta dalla teoria della Gestalt. Secondo tale teoria, la percezione è governata da leggi che ne regolano l'organizzazione. La distinzione tra figura e sfondo, la percezione di profondità e movimento e l'integrazione di informazioni spaziali e temporali nei percetti dipendono da alcuni principi percettivi di base.
\subsubsection{Articolazione figura-sfondo}
Uno dei compiti principali dei processi percettivi è decidere, all'interno di una scena, quale elemento debba essere considerato \emph{figura} e quale \emph{sfondo}.
\begin{snippetdefinition}{Articolazione figura-sfondo}
    L'\emph{articolazione figura-sfondo} è un processo percettivo universale e costante, poiché non esiste figura senza sfondo. La figura ha forma ed estensione definita, mentre lo sfondo è amorfo, indifferenziato e continua dietro la figura in modo indeterminato.
\end{snippetdefinition}
I fattori alla base dell'articolazione figura-sfondo sono:
\begin{itemize}
    \item \emph{inclusione};
    \item \emph{convessità};
    \item \emph{area relativa};
    \item \emph{orientamento}.
\end{itemize}
Se questi fattori non intervengono, si creano le condizioni per ottenere le \emph{figure reversibili}, nelle quali si ha un'inversione tra figura e sfondo. Si tratta di configurazioni instabili e ambigue, in cui figura e sfondo si alternano regolarmente.
\begin{snippetdefinition}{Figure reversibili}
    Le \emph{figure reversibili} sono configurazioni percettive instabili e ambigue, nelle quali figura e sfondo possono invertirsi e alternarsi.
\end{snippetdefinition}
\subsubsection{Principi di raggruppamento percettivo}
I principi di raggruppamento percettivo sono stati studiati in modo approfondito dagli psicologi della Gestalt. Secondo questo approccio, i fenomeni psicologici possono essere compresi solo se percepiti all'interno di un'organizzazione, come un tutto unitario, e non se vengono scomposti in singoli elementi percettivi di base.
\emph{Max Wertheimer} formulò i seguenti principi:
\begin{itemize}
    \item \emph{legge della vicinanza:} a parità di condizioni, si unificano gli elementi vicini;
    \item \emph{legge della somiglianza:} a parità di condizioni, si unificano gli elementi simili;
    \item \emph{legge della buona direzione:} a parità di condizioni, si unificano gli elementi che presentano continuità di direzione;
    \item \emph{legge della chiusura:} a parità di condizioni, si unificano gli elementi che tendono a chiudersi tra loro;
    \item \emph{legge del destino comune:} a parità di condizioni, si unificano gli elementi che condividono la medesima direzione di movimento.
\end{itemize}
\subsubsection{Integrazione spaziale e temporale}
I principi di raggruppamento percettivo suggeriscono che molti processi percettivi si basano sulla capacità di mettere insieme le informazioni nel modo corretto. Tuttavia, spesso non siamo in grado di cogliere un'intera scena attraverso una sola occhiata o una sola fissazione oculare.
Per avere un'idea esaustiva di ciò che ci circonda è necessario combinare le informazioni relative a fissazioni oculari collocate in punti spaziali differenti, attraverso l'\emph{integrazione spaziale}, e in momenti temporali differenti, attraverso l'\emph{integrazione temporale}.
\begin{snippetdefinition}{Integrazione spaziale}
    L'\emph{integrazione spaziale} è il processo attraverso cui il sistema percettivo combina informazioni provenienti da differenti collocazioni spaziali.
\end{snippetdefinition}
\begin{snippetdefinition}{Integrazione temporale}
    L'\emph{integrazione temporale} è il processo attraverso cui il sistema percettivo combina informazioni raccolte in differenti momenti nel tempo.
\end{snippetdefinition}
\begin{center}

    \includegraphics[width=0.70\textwidth]{media_psychology/boundary-extension-example.jpg}
    \captionof{figure}{Esempio di estensione dei confini: la stessa scena è rappresentata con due inquadrature differenti.}
\end{center}
\begin{snippetdefinition}{Estensione dei confini}
    L'\emph{estensione dei confini} è un errore percettivo e mnestico per cui, dopo aver osservato una scena, le persone tendono a ricordare di aver visto una porzione più ampia rispetto a quella effettivamente presentata.
\end{snippetdefinition}
Quando osserviamo una scena, utilizziamo i processi mnestici per estendere i confini dell'inquadratura. Per questo motivo possiamo ricordare di aver visto prospettive in grandangolo anche quando, in realtà, abbiamo osservato figure in primo piano.
La nostra conoscenza del mondo viene usata per estendere la visuale oltre ciò che riusciamo a vedere direttamente, come accade quando guardiamo fuori da una finestra. Per questo possiamo avere difficoltà a notare cambiamenti tra prospettive diverse di una stessa scena.
\begin{snippetdefinition}{Cecità al cambiamento}
    La \emph{cecità al cambiamento}, o \emph{change blindness}, è la difficoltà a notare cambiamenti tra diverse prospettive o versioni di una stessa scena.
\end{snippetdefinition}
\subsubsection{Percezione del movimento}
Un tipo di percezione che richiede confronti tra differenti osservazioni visive è la \emph{percezione del movimento}.
\begin{snippetdefinition}{Percezione del movimento}
    La \emph{percezione del movimento} è la capacità di combinare informazioni provenienti da differenti osservazioni visive per percepire lo spostamento di uno stimolo nello spazio.
\end{snippetdefinition}
Le conseguenze dei processi di integrazione visiva possono essere comprese osservando il \emph{fenomeno phi}. Questo fenomeno si verifica quando due punti luce statici, collocati in punti differenti del campo visivo, vengono accesi e spenti in modo alternato a un ritmo di circa quattro o cinque volte al secondo. Gli osservatori umani non percepiscono due punti separati, ma un unico percorso semplice, cioè una linea retta.
\subsubsection{Percezione della profondità}
Nei paragrafi precedenti sono state considerate soprattutto immagini bidimensionali poste su una superficie piatta. Nel mondo reale, però, abbiamo a che fare con oggetti collocati in uno spazio tridimensionale.
Percepire le dimensioni dello spazio è fondamentale per interagire con persone e oggetti, raggiungere ciò che è desiderato ed evitare ciò che è pericoloso. Questa percezione richiede informazioni accurate sulla \emph{profondità dell'oggetto}, ottenute a partire dalla distanza dell'oggetto percepito e dalla sua direzione rispetto alla posizione dell'osservatore.
\begin{snippetdefinition}{Percezione della profondità}
    La \emph{percezione della profondità} è la capacità di percepire la distanza e la posizione degli oggetti nello spazio tridimensionale.
\end{snippetdefinition}
L'interpretazione della profondità si basa su numerose fonti informative relative alla distanza, tra cui gli \emph{indizi binoculari} e gli \emph{indizi monoculari}.
\begin{snippetdefinition}{Indizi binoculari di profondità}
    Gli \emph{indizi binoculari di profondità} sono informazioni sulla profondità che derivano dal confronto tra le informazioni visive fornite dai due occhi.
\end{snippetdefinition}
Le due principali sorgenti di informazione binoculare sono la \emph{disparità retinica} e la \emph{convergenza}.
\begin{snippetdefinition}{Disparità retinica}
    La \emph{disparità retinica} è lo scarto tra le posizioni orizzontali retiniche di immagini corrispondenti nei due occhi. Fornisce informazioni sulla profondità perché il grado di disparità dipende dalla distanza relativa tra l'oggetto e l'individuo.
\end{snippetdefinition}
I due occhi sono separati da una distanza di circa \(5\)-\(6\) centimetri e osservano quindi la realtà da punti di vista leggermente diversi.
\begin{snippetdefinition}{Convergenza}
    La \emph{convergenza} è un'informazione binoculare relativa alla profondità, basata sul movimento degli occhi verso l'interno quando si osserva uno stimolo vicino.
\end{snippetdefinition}
\begin{snippetdefinition}{Movimento di parallasse}
    Il \emph{movimento di parallasse} è un'informazione di profondità determinata dalla velocità e dalla direzione del movimento relativo degli oggetti nell'immagine retinica.
\end{snippetdefinition}
\begin{snippetexample}{Movimento di parallasse}
    Quando ci si trova sul sedile del passeggero in automobile, gli oggetti collocati a una certa distanza sembrano più statici rispetto agli oggetti più vicini, che sembrano muoversi più rapidamente.
\end{snippetexample}
Anche osservando un oggetto con un solo occhio è possibile percepirne la profondità. Le fonti informative disponibili nella visione monoculare sono chiamate \emph{indizi monoculari di profondità}.
\begin{snippetdefinition}{Indizi monoculari di profondità}
    Gli \emph{indizi monoculari di profondità} sono informazioni sulla profondità che possono essere ricavate anche usando un solo occhio.
\end{snippetdefinition}
I principali indizi monoculari sono:
\begin{itemize}
    \item \emph{occlusione:} quando un oggetto opaco si sovrappone parzialmente a un altro, l'oggetto nascosto viene percepito come più lontano;
    \item \emph{dimensione relativa:} oggetti della stessa dimensione posti a distanze differenti proiettano immagini di dimensioni differenti sulla retina;
    \item \emph{prospettiva lineare:} indizio di profondità che dipende dalla relazione tra dimensione e distanza;
    \item \emph{gradienti di trama:} la densità della trama aumenta quanto più la superficie si allontana.
\end{itemize}
\subsubsection{Le costanze percettive}
In generale, il mondo è percepito come \emph{invariante}, \emph{costante} e \emph{stabile} al di là dei cambiamenti nella stimolazione dei recettori sensoriali.
\begin{snippetdefinition}{Costanza percettiva}
    La \emph{costanza percettiva} è il fenomeno per cui le proprietà degli oggetti vengono percepite come stabili e costanti nonostante le variazioni della stimolazione sensoriale.
\end{snippetdefinition}
Per ragioni di sopravvivenza è fondamentale percepire come stabili le proprietà degli oggetti al di là delle grandi variazioni delle stimolazioni luminose che colpiscono gli occhi. Il compito critico della percezione è scoprire proprietà invarianti dell'ambiente oltre le variazioni delle impressioni retiniche.
Le principali costanze percettive sono:
\begin{itemize}
    \item \emph{costanza di grandezza};
    \item \emph{costanza di forma};
    \item \emph{costanza di luminosità}.
\end{itemize}
\begin{snippetdefinition}{Costanza di grandezza}
    La \emph{costanza di grandezza} è la capacità di percepire le dimensioni effettive di un oggetto al di là delle variazioni di grandezza della sua immagine retinica.
\end{snippetdefinition}
La costanza di grandezza dipende sia dalle dimensioni reali di un oggetto sia dalla sua distanza dall'occhio. Il sistema percettivo può valutare le dimensioni di un oggetto anche utilizzando la conoscenza precedente circa la grandezza di oggetti simili.
\begin{snippetdefinition}{Costanza di forma}
    La \emph{costanza di forma} è la capacità di percepire la forma di un oggetto anche quando esso è inclinato verso l'osservatore e produce un'immagine retinica differente dalla forma dell'oggetto fisico.
\end{snippetdefinition}
\begin{snippetdefinition}{Costanza di luminosità}
    La \emph{costanza di luminosità} è la tendenza a percepire il bianco, il grigio o il nero degli oggetti come costanti nonostante le variazioni di illuminazione.
\end{snippetdefinition}
\subsection{Processi dall'alto verso il basso: l'importanza dei contesti e delle aspettative}
L'obiettivo principale della percezione è fornire una visione accurata del mondo. La sopravvivenza dipende anche dalla capacità di compiere valutazioni affidabili di oggetti ed eventi collocati in un determinato contesto.
Tuttavia, in molte occasioni i processi \emph{bottom-up} non permettono di costruire un'identificazione univoca di uno stimolo percepito e, in alcuni casi, non consentono di formulare alcuna ipotesi interpretativa. In queste circostanze, i processi \emph{top-down} utilizzano il contesto e le aspettative per attribuire un significato chiaro a ciò che viene percepito.
\begin{snippetdefinition}{Illusione}
    Un'\emph{illusione} si verifica quando uno stimolo viene percepito in modo non corretto.
\end{snippetdefinition}
Gli studi sulle illusioni visive hanno evidenziato che il contesto influenza le assunzioni relative alle caratteristiche visive dello stimolo.
\begin{center}

    \includegraphics[width=0.45\textwidth]{media_psychology/duck-rabbit-ambiguous-figure.jpg}
    \captionof{figure}{Figura ambigua interpretabile come anatra o coniglio.}
\end{center}
È possibile anche controllare le illusioni per ottenere determinati effetti. Ad esempio, gli psicologi della NASA hanno studiato gli effetti dell'ambiente sulla percezione per creare capsule spaziali dotate di qualità sensoriali gradevoli.
\subsubsection{Effetti contestuali sul raggruppamento percettivo}
Il processo di raggruppamento percettivo avviene in modo automatico. Più elementi sono presenti nel contesto, più diventa difficile percepire un elemento separatamente dagli altri.
\begin{snippetexample}{Illusione di Ebbinghaus}
    Nell'\emph{illusione di Ebbinghaus}, i cerchi centrali nei due gruppi hanno la stessa dimensione, ma quello circondato da cerchi più piccoli sembra più grande dell'altro. L'illusione è tanto più forte quanto più le forme degli oggetti sono simili e vengono percepite come appartenenti allo stesso gruppo.
\end{snippetexample}
\begin{center}

    \includegraphics[width=0.68\textwidth]{media_psychology/ebbinghaus-illusion.jpg}
    \captionof{figure}{Illusione di Ebbinghaus.}
\end{center}
In certi casi l'analisi del contesto è necessaria per raggiungere risultati ancora più difficili. Il contesto e le aspettative svolgono un ruolo fondamentale nella vita quotidiana: riconoscere una persona in un luogo inatteso richiede più tempo e può produrre incertezza.
\subsection{Considerazioni conclusive}
Se i processi percettivi fossero completamente dal basso verso l'alto, saremmo legati alla stessa realtà contingente e concreta del \emph{qui e ora}. Saremmo in grado di registrare l'esperienza, ma non potremmo utilizzarla in occasioni future né percepire il mondo in modo differente a seconda delle circostanze.
Se i processi percettivi fossero invece completamente dall'alto verso il basso, potremmo perderci nel mondo di \emph{fantasie personali} relative a ciò che ci aspettiamo di ricevere o che speriamo di percepire.
\begin{snippetnote}{Equilibrio tra bottom-up e top-down}
    Il giusto equilibrio tra processi \emph{bottom-up} e processi \emph{top-down} permette di realizzare il principale obiettivo della percezione: sperimentare la realtà nel modo più efficace per rispondere alle necessità biologiche e sociali dell'individuo.
\end{snippetnote}
\begin{center}

    \includegraphics[width=0.48\textwidth]{media_psychology/bottom-up-top-down-perception-balance.jpg}
    \captionof{figure}{Equilibrio tra processi bottom-up e processi top-down nella percezione.}
\end{center}
Grazie a questo equilibrio, possiamo muoverci e adattarci all'ambiente fisico e sociale.

\section{Attenzione e coscienza}
Provate a pensare a un evento del vostro passato che ricordate con piacere. Ora pensate a che cosa vi piacerebbe che accadesse domani o tra due giorni. Sebbene nel cervello sia immagazzinata una grande quantità di informazioni, è molto probabile che questi pensieri non fossero già presenti nella vostra coscienza prima che vi fossero richiesti. Di conseguenza, pensieri di cui prima non eravate consapevoli sono giunti alla coscienza proprio in quel momento.
\emph{Coscienza} e \emph{attenzione} sono due dimensioni psicologiche fondamentali per l'uso di numerose capacità psichiche, dalla percezione all'apprendimento. Si ritiene che la coscienza aiuti la sopravvivenza e fornisca un senso di chi si è e del posto che si occupa nel mondo.
\subsection{L'attenzione}
\begin{snippetexample}{Attenzione in una festa}
    Durante una festa rumorosa state cercando un'amica in mezzo alla folla. Tra molti colori individuate il vestito rosa che indossa; nonostante la musica e le conversazioni, riuscite a sentire la sua voce che vi chiama. Poco dopo, un rumore di vetri rotti vi induce a fermarvi e a girarvi. Tutti questi processi, cioè individuare l'amica, sentire il proprio nome, orientarsi verso un rumore improvviso e poi riprendere la ricerca, richiedono il coinvolgimento dell'\emph{attenzione}.
\end{snippetexample}
\begin{snippetdefinition}{Attenzione}
    Dal punto di vista cognitivo, l'\emph{attenzione} è l'insieme dei meccanismi che consentono di concentrare e focalizzare le proprie risorse mentali su alcune informazioni, definendo ciò di cui siamo consapevoli in un dato momento.
\end{snippetdefinition}
La \emph{valorizzazione} permette di selezionare determinate informazioni da sottoporre a ulteriori elaborazioni, mentre l'\emph{inibizione} permette di trascurare altre informazioni considerate irrilevanti.
Durante la giornata siamo continuamente esposti a una quantità molto ampia di informazioni. Le nostre capacità cognitive non sono in grado di attribuire un significato a tutti gli stimoli presenti; per questo è necessario selezionare, momento dopo momento, le informazioni rilevanti e ignorare quelle trascurabili.
Una possibile soluzione consiste nel focalizzarsi su alcune parti delle informazioni disponibili, sottoponendole a elaborazione ulteriore in base all'importanza attribuita loro in un dato momento. Tuttavia, anche quando conosciamo bene le nostre intenzioni e i nostri obiettivi, alcuni aspetti dello stimolo, se abbastanza salienti, possono attirare l'attenzione e distrarci.
\subsubsection{Attenzione selettiva}
\begin{snippetdefinition}{Attenzione selettiva}
    L'\emph{attenzione selettiva} è un processo che seleziona le informazioni sulla base della loro salienza e le elabora con particolare efficienza.
\end{snippetdefinition}
L'attenzione selettiva non riguarda soltanto la posizione di uno stimolo nello spazio, ma tutte le proprietà degli oggetti e degli eventi, come forma, colore, dimensione, suono e movimento.
L'attenzione selettiva può riferirsi a:
\begin{itemize}
    \item una specifica modalità sensoriale, come l'attenzione visiva o uditiva;
    \item una specifica zona dello spazio circostante;
    \item uno specifico stimolo e alle sue caratteristiche, come colore o forma;
    \item uno specifico movimento.
\end{itemize}
In questo senso, l'attenzione è \emph{basata sugli oggetti}, poiché procede selezionando gli oggetti piuttosto che semplici coordinate spaziali.
\begin{snippetdefinition}{Attenzione basata sugli oggetti}
    L'\emph{attenzione basata sugli oggetti} è una forma di attenzione per cui, quando l'attenzione è diretta verso un oggetto, tutte le parti di quell'oggetto vengono selezionate simultaneamente per essere elaborate.
\end{snippetdefinition}
\begin{snippetexample}{Oggetto e attenzione}
    Se vi focalizzate su un'amica vestita di rosa, avete una maggiore probabilità di notare l'orologio al suo polso, poiché è parte dello stesso oggetto percettivo, rispetto all'orologio indossato dalla persona accanto, anche se i due orologi si trovano alla stessa distanza da voi.
\end{snippetexample}
\begin{snippetdefinition}{Fuoco dell'attenzione}
    Il \emph{fuoco dell'attenzione} consente di concentrare le risorse attentive su uno specifico stimolo dell'ambiente.
\end{snippetdefinition}
Il fuoco dell'attenzione può variare per dimensione: talvolta è concentrato sui particolari, altre volte è più esteso e disperso. Esso presenta una relazione inversa con l'efficienza di elaborazione delle informazioni: quanto più è ristretta l'area dell'attenzione, tanto maggiore è l'efficienza cognitiva, e viceversa.
\begin{center}
    \includegraphics[width=0.62\textwidth]{media_psychology/selective-attention-example.jpg}
    \captionof{figure}{Esempio di attenzione selettiva.}
\end{center}
In ogni momento siamo circondati da molti stimoli. Le cose su cui ci focalizziamo, insieme alle memorie che esse evocano, determinano in larga misura che cosa si trova nella coscienza.
Secondo \emph{Donald Broadbent}, la mente opera soltanto su alcuni stimoli in entrata grazie all'intervento di un sistema di filtraggio. In questa prospettiva, l'elaborazione degli stimoli a cui non si presta attenzione fallisce prima che le persone riescano a identificarli.
\begin{snippetdefinition}{Effetto cocktail party}
    L'\emph{effetto cocktail party} è il fenomeno per cui una persona riesce a percepire informazioni rilevanti, come il proprio nome, in un ambiente rumoroso, anche quando è impegnata in un'altra conversazione.
\end{snippetdefinition}
L'effetto cocktail party sembra contraddire l'idea di un filtro attentivo completamente precoce, perché mostra che le persone sono sensibili a percepire il proprio nome anche tra informazioni a cui non stanno prestando attenzione.
\subsubsection{Attenzione divisa}
In molte situazioni, l'attenzione è divisa e distribuita su diversi stimoli, come quando si eseguono simultaneamente due compiti. In questi casi la selezione dell'informazione è meno accurata.
\begin{snippetdefinition}{Attenzione divisa}
    L'\emph{attenzione divisa} è un processo in cui l'attenzione viene distribuita tra diversi stimoli o compiti, a partire da una quantità limitata di risorse cognitive.
\end{snippetdefinition}
\begin{snippetexample}{Esperimento di Duncan}
    In un esperimento di \emph{Duncan}, i partecipanti vedevano al centro di un monitor un rettangolo con una piccola apertura su un lato e una linea che lo attraversava. Quando dovevano rispondere a due domande relative allo stesso oggetto, ad esempio se il rettangolo fosse grande o piccolo e se l'apertura fosse a sinistra o a destra, l'accuratezza era elevata. Quando invece dovevano rispondere a domande relative sia al rettangolo sia alla linea, l'accuratezza diminuiva in modo significativo.
\end{snippetexample}
\begin{center}

    \includegraphics[width=0.55\textwidth]{media_psychology/duncan-divided-attention-stimuli.jpg}
    \captionof{figure}{Esempio di stimoli utilizzati nello studio di Duncan sull'attenzione divisa.}
\end{center}
È esperienza comune che, di norma, sia più difficile eseguire due compiti nello stesso momento che eseguirli uno dopo l'altro.
Secondo Duncan, l'attenzione è l'esito di una competizione tra diversi stimoli che può verificarsi in ogni fase dell'elaborazione delle informazioni. Lo stimolo su cui si concentrano maggiori risorse viene analizzato in modo più esaustivo rispetto agli altri.
\subsection{Che cos'è la coscienza?}
\begin{snippetdefinition}{Coscienza}
    La \emph{coscienza} è la consapevolezza degli stimoli esterni e interni da parte del soggetto. Essa è uno stato più che un processo e di norma accompagna un'esperienza che riflette un determinato processo di elaborazione delle informazioni.
\end{snippetdefinition}
Per essere consapevoli di una particolare informazione è necessario essere coscienti.
Negli ultimi anni si è affermata l'idea secondo cui le operazioni mentali si svolgono abitualmente in modo inconscio. Solo in casi particolari, grazie all'intervento dei processi attentivi, le rappresentazioni prodotte dall'elaborazione delle informazioni possono diventare consce.
\subsubsection{I contenuti della coscienza}
Nell'ambito dell'elaborazione delle informazioni, la coscienza sarebbe presente nelle \emph{conoscenze dichiarative}, o \emph{proposizionali}, cioè nelle proposizioni che stabiliscono una relazione tra due o più idee e riguardano i contenuti della vita quotidiana.
\begin{snippetdefinition}{Conoscenze dichiarative}
    Le \emph{conoscenze dichiarative}, o \emph{proposizionali}, sono conoscenze espresse sotto forma di proposizioni che stabiliscono relazioni tra idee e riguardano i contenuti della vita quotidiana.
\end{snippetdefinition}
Le \emph{conoscenze procedurali}, invece, riguardano le procedure con cui vengono svolti i compiti quotidiani e si basano sull'esercizio e sull'esecuzione di un determinato compito.
\begin{snippetdefinition}{Conoscenze procedurali}
    Le \emph{conoscenze procedurali} riguardano le procedure attraverso cui si svolgono i compiti della vita quotidiana e dipendono dall'esercizio e dalla pratica.
\end{snippetdefinition}
La coscienza può essere operante nella prima fase dell'apprendimento di un'abilità, ad esempio quando si impara a nuotare o ad andare in bicicletta. Con l'esercizio, però, queste attività diventano immediate e automatiche.
Per questo motivo si distingue tra:
\begin{itemize}
    \item \emph{elaborazione automatica:} opera molto rapidamente, non richiede risorse attentive e avviene senza l'intervento della coscienza; diversi processi automatici possono svolgersi in parallelo;
    \item \emph{elaborazione controllata:} è lenta, richiede l'intervento delle risorse attentive ed è consapevole; non permette di svolgere un altro compito nello stesso momento, poiché procede in serie.
\end{itemize}
In realtà non esiste un processo totalmente automatico, poiché è sempre richiesta una certa quantità di risorse cognitive. Qualsiasi processo automatico, come guidare l'auto, può tornare a essere controllato in caso di necessità, ad esempio in presenza di ghiaccio sulla strada.
\subsubsection{Consapevolezza e coscienza}
Gran parte della ricerca recente in psicologia riguarda i contenuti della coscienza. Allontanandosi gradualmente dalla filosofia dell'Ottocento, la psicologia è diventata la scienza della mente.
\emph{William James} definisce la psicologia come la descrizione e la spiegazione degli stati di coscienza in quanto tali.
Il nostro stato abituale di coscienza include percezioni, pensieri, sensazioni, immagini e desideri presenti in un dato momento: tutta l'attività mentale su cui stiamo focalizzando l'attenzione. Siamo coscienti di ciò che stiamo facendo e del fatto che lo stiamo facendo. Dall'esperienza privilegiata di osservarsi dall'interno deriva un senso di sé.
La coscienza, pur essendo una funzione autonoma, si fonda e si radica sui processi inconsci. Freud aveva già messo in luce questo aspetto attraverso la metafora dell'\emph{iceberg}: la parte sommersa rappresenta l'inconscio, mentre la punta rappresenta la coscienza.
Alla luce delle teorie più recenti, i contenuti della coscienza sarebbero l'esito funzionale di processi nervosi elaborati in modo inconscio. Di conseguenza, non siamo consapevoli dei processi in se stessi, ma soltanto dei loro risultati finali.
\subsubsection{Studiare i contenuti della coscienza}
Per studiare la coscienza, i ricercatori hanno dovuto sviluppare metodologie capaci di misurare apertamente esperienze profondamente private.
Un metodo deriva dall'evoluzione della pratica di introspezione messa a punto da Wundt: ai partecipanti viene chiesto di parlare ad alta voce mentre svolgono compiti complessi, riportando nel modo più dettagliato possibile la sequenza dei pensieri che attraversano la loro mente.
\begin{snippetdefinition}{Think-aloud protocols}
    I \emph{think-aloud protocols}, o \emph{protocolli di verbalizzazione del pensiero}, sono resoconti prodotti dai partecipanti mentre svolgono un compito e servono a documentare strategie mentali e rappresentazioni della conoscenza impiegate nello svolgimento del compito.
\end{snippetdefinition}
\begin{snippetdefinition}{Experience-sampling method}
    L'\emph{experience-sampling method}, o \emph{ESM}, è un metodo in cui i partecipanti forniscono informazioni sui propri pensieri e sentimenti durante l'arco della giornata.
\end{snippetdefinition}
Queste tecniche permettono ai ricercatori di stabilire, per particolari compiti e momenti, quale sottoinsieme delle informazioni disponibili all'individuo sia presente nella coscienza.
\subsection{Le funzioni della coscienza}
Quando ci si interroga sulle funzioni della coscienza, ci si chiede perché abbiamo bisogno della coscienza e che cosa essa aggiunga all'esperienza umana.
William James descrive l'enorme quantità di informazioni che colpiscono i recettori sensoriali come una ``fiorente e ronzante confusione''. La coscienza aiuta l'adattamento all'ambiente dando senso a questa confusione in tre modi:
\begin{enumerate}
    \item svolge una \emph{funzione restrittiva}, riducendo il flusso di stimoli in ingresso e aiutando a ignorare le informazioni non rilevanti per scopi e obiettivi immediati;
    \item permette l'\emph{immagazzinamento selettivo}, cioè consente di conservare in modo selettivo le informazioni da analizzare;
    \item rende possibile la \emph{pianificazione}, permettendo di fermarsi a pensare, considerare alternative, usare conoscenze passate e immaginare conseguenze future.
\end{enumerate}
\subsubsection{Studiare le funzioni della coscienza}
Il comportamento delle persone è spesso influenzato da processi sia consci sia inconsci. Molte funzioni della coscienza includono rimandi impliciti a ciò che rimane inconscio: frequentemente i processi consci influenzano quelli inconsci e viceversa.
\begin{snippetexample}{Dilemma morale}
    Immaginate che un aereo abbia effettuato un atterraggio di emergenza in pieno oceano. C'è un'unica scialuppa di salvataggio, ma sta affondando per il peso eccessivo. Si potrebbe buttare in acqua alcune persone per salvare il gruppo, oppure non farlo, mettendo però a rischio la vita di tutti. La risposta immediata e ``di pancia'', prodotta da processi inconsci, tende a rifiutare il sacrificio. Una scelta ponderata, invece, può portare a concludere che il sacrificio salverebbe il maggior numero possibile di persone.
\end{snippetexample}
Un modo per studiare le funzioni della coscienza consiste nel mostrare come cambiano le risposte delle persone quando non è consentito il normale funzionamento dei processi consci. Un altro modo consiste nel determinare quali compiti richiedano l'intervento dei processi consci.
\subsection{Sonno e sogni}
Ogni giorno sperimentiamo una profonda alterazione della coscienza: quando decidiamo che la giornata è finita, ci abbandoniamo al \emph{sonno} e, mentre dormiamo, sogniamo.
Circa un terzo della vita è dedicato al sonno, durante il quale i muscoli si trovano in uno stato di paralisi benigna e il cervello mantiene un'intensa attività.
\subsubsection{Ritmi circadiani}
Tutte le creature sono influenzate dai ritmi naturali del giorno e della notte. Il corpo umano è regolato da un ciclo temporale chiamato \emph{ritmo circadiano}.
\begin{snippetdefinition}{Ritmo circadiano}
    Il \emph{ritmo circadiano} è un ciclo temporale interno che regola livelli di attivazione, metabolismo, battito cardiaco e temperatura corporea, sincronizzandoli con l'orologio biologico dell'organismo.
\end{snippetdefinition}
L'informazione relativa alla luce del sole viene raccolta attraverso gli occhi, ma i recettori coinvolti nella regolazione dei ritmi circadiani non sono gli stessi che permettono la visione del mondo.
Le circostanze che producono un allontanamento tra l'orologio interno e il ciclo del sonno influenzano il modo in cui ci si sente e si agisce. Chi svolge il turno di notte può avere difficoltà fisiche e cognitive, perché i suoi ritmi circadiani sono alterati. Sensazioni simili compaiono anche dopo lunghi viaggi in aereo attraverso più fusi orari, cioè nel \emph{jet lag}.
\subsubsection{Il ciclo del sonno}
Circa un terzo del ritmo circadiano è dedicato al sonno, una forma di quiescenza comportamentale. Gran parte delle conoscenze sul sonno riguarda l'attività elettrica del cervello.
La svolta metodologica nello studio del sonno avviene con l'applicazione dell'\emph{elettroencefalogramma}, che fornisce una misura oggettiva e continua del variare dell'attività cerebrale in condizioni di sonno e veglia.
Grazie all'EEG si scopre che le onde cerebrali cambiano forma quando ci si addormenta e mostrano variazioni sistematiche e prevedibili durante tutto il periodo del sonno.
Un'altra scoperta importante è che, a intervalli periodici durante il sonno, si verifica una scarica di movimenti oculari rapidi, chiamati \emph{rapid eye movements}, o \emph{REM}.
\begin{snippetdefinition}{Sonno REM}
    Il \emph{sonno REM} è una fase del sonno caratterizzata da movimenti oculari rapidi, aumento dell'attività elettrica cerebrale e sogni vividi.
\end{snippetdefinition}
Quando ci si addormenta, le onde cerebrali rallentano e si entra nel ciclo del sonno. Ogni fase presenta un andamento EEG distinto. Nelle fasi \(3\) e \(4\) si è in uno stato di sonno profondo, in cui rallentano ulteriormente le onde cerebrali, la respirazione e il battito cardiaco.
Nell'ultima fase l'attività elettrica aumenta: è in questo momento che si verifica il sonno REM e iniziano i sogni vividi.
Le prime quattro fasi del sonno, cioè il sonno non-REM, coprono circa \(90\) minuti, mentre il sonno REM dura circa \(10\) minuti. In una notte, questo ciclo completo di circa \(100\) minuti si ripete dalle quattro alle sei volte. Con il passare dei cicli aumenta il tempo trascorso in fase REM, che può durare anche un'ora nell'ultimo ciclo.
\begin{center}

    \includegraphics[width=0.62\textwidth]{media_psychology/sleep-cycle-nrem-rem.jpg}
    \captionof{figure}{Ciclo del sonno e alternanza tra fasi non-REM e fase REM.}
\end{center}
Non tutti gli individui dormono lo stesso numero di ore. Sebbene esista un bisogno genetico di sonno nella specie umana, la quantità di sonno effettiva dipende anche dalle attività svolte durante la veglia.
Gli schemi del sonno cambiano nel corso della vita. Il neonato dorme circa \(16\) ore al giorno, metà delle quali in fase REM. Verso i \(50\) anni è possibile dormire solo \(6\) ore, mentre i giovani adulti dormono normalmente dalle \(7\) alle \(8\) ore.
\subsubsection{Perché dormiamo?}
Le persone funzionano abbastanza bene quando riescono a dormire \(7\)-\(8\) ore per notte. Le funzioni generali del sonno non-REM sono il \emph{risparmio energetico} e il \emph{riposo}.
Il sonno non-REM potrebbe essersi evoluto per permettere agli animali di risparmiare energia nei momenti in cui non era necessario lavorare o cercare cibo. Inoltre, aiuterebbe il cervello a riparare i danni prodotti dal lavoro svolto dalle cellule neuronali.
Se una notte si viene privati del sonno REM, nella notte successiva se ne avrà di più. Questo suggerisce che anche il sonno REM svolga funzioni fondamentali. Probabilmente esso fornisce al cervello un contesto in cui consolidare nuovi apprendimenti.
\begin{snippetdefinition}{Deprivazione di sonno}
    La \emph{deprivazione di sonno} è una condizione in cui una persona dorme troppo poco, con conseguenze negative sulle prestazioni cognitive e motorie.
\end{snippetdefinition}
La deprivazione di sonno produce effetti negativi sulle prestazioni cognitive, come difficoltà di attenzione e problemi nella memoria di lavoro. Inoltre danneggia la capacità di svolgere compiti di tipo motorio.
\subsection{Stati di coscienza alterati}
Ogni cultura ha sviluppato pratiche che permettono di andare oltre la coscienza ordinaria e sperimentare \emph{stati di coscienza alterati}. Alcune di queste pratiche sono individuali, come la meditazione, mentre altre sono tentativi condivisi di trascendere i confini della normale esperienza consapevole.
\begin{snippetdefinition}{Ipnosi}
    L'\emph{ipnosi} è uno stato di consapevolezza alternativo, caratterizzato dalla particolare abilità di alcune persone di rispondere alle suggestioni con cambiamenti nella percezione, nella memoria, nella motivazione e nel senso di controllo di sé.
\end{snippetdefinition}
Nello stato ipnotico, i partecipanti mostrano una responsività accresciuta ai suggerimenti dell'ipnotizzatore.
\begin{snippetdefinition}{Induzione ipnotica}
    L'\emph{induzione ipnotica} è l'insieme delle attività preliminari che minimizzano le distrazioni esterne e incoraggiano i partecipanti a concentrarsi sugli stimoli suggeriti, credendo di entrare in uno speciale stato di coscienza.
\end{snippetdefinition}
La procedura tipica di induzione utilizza suggestioni che portano a un profondo rilassamento.
\begin{snippetdefinition}{Ipnotizzabilità}
    L'\emph{ipnotizzabilità} è il grado di responsività dell'individuo alla suggestione ipnotica.
\end{snippetdefinition}
Esistono ampie differenze individuali nella suscettibilità all'ipnosi: alcune persone non rispondono quasi per nulla, mentre altre mostrano una responsività molto elevata.
L'ipnosi ha anche un potenziale importante nel controllo del dolore, cioè nell'\emph{analgesia ipnotica}, poiché la mente può amplificare gli stimoli dolorifici attraverso anticipazione e paura.
\begin{snippetdefinition}{Analgesia ipnotica}
    L'\emph{analgesia ipnotica} è il controllo del dolore ottenuto attraverso suggestioni ipnotiche, come immaginare la parte dolente separata dal resto del corpo o priva di sensibilità.
\end{snippetdefinition}
\begin{snippetdefinition}{Meditazione}
    La \emph{meditazione} è una forma di alterazione della coscienza progettata per migliorare la conoscenza di sé e il benessere attraverso il raggiungimento di uno stato di profonda tranquillità.
\end{snippetdefinition}
Nella \emph{meditazione concentrativa}, una persona può focalizzarsi e regolare il respiro, assumere specifiche posizioni corporee, ridurre al minimo le stimolazioni esterne, generare immagini mentali e liberare la mente dai pensieri.
Nella \emph{meditazione mindfulness}, invece, l'individuo impara a lasciare che ricordi e pensieri attraversino la mente senza reagire a essi.
Le tecniche di neuroimaging hanno mostrato che la meditazione influenza gli schemi di attività cerebrale. Risultati recenti suggeriscono che, a lungo termine, la meditazione possa avere un effetto benefico sul cervello stesso.

\section{Apprendimento ed esperienza}
\begin{snippetdefinition}{Apprendimento}
    L'\emph{apprendimento} è un processo continuo, basato sull'esperienza, che si traduce in un cambiamento relativamente stabile e duraturo nel comportamento, nel comportamento potenziale o negli stati rappresentativi.
\end{snippetdefinition}
L'apprendimento può avvenire solo tramite l'\emph{esperienza} e non riguarda i cambiamenti dovuti alla sola maturazione biologica. L'esperienza include una fase di raccolta, valutazione e rielaborazione delle informazioni e una fase di azione volta a produrre risposte capaci di modificare l'ambiente. L'apprendimento consiste dunque in risposte influenzate dalle informazioni contenute nella memoria.
\begin{center}

    \includegraphics[width=0.62\textwidth]{media_psychology/learning-information-memory-action.jpg}
    \captionof{figure}{Schema generale dell'apprendimento come relazione tra raccolta di informazioni, memoria e azione sull'ambiente.}
\end{center}
L'apprendimento è facilmente riscontrabile quando se ne possono osservare i risultati. Tuttavia, può essere analizzato anche in modo più diretto osservando i cambiamenti che avvengono all'interno del cervello: non cambia soltanto il comportamento visibile dall'esterno, ma anche l'organizzazione interna dell'organismo.
L'apprendimento può essere rilevato attraverso un cambiamento a livello delle \emph{prestazioni}, anche se la prestazione da sola non sempre mostra tutto ciò che è stato appreso. Talvolta si acquisiscono attitudini generali, come un interesse verso l'arte moderna, che potrebbero non essere evidenti in azioni concretamente misurabili se non si presenta l'occasione di esibire ciò che è stato appreso.
\begin{snippetdefinition}{Distinzione tra apprendimento e prestazione}
    La \emph{distinzione tra apprendimento e prestazione} riguarda la differenza tra ciò che è stato appreso, cioè le potenzialità disponibili, e ciò che viene effettivamente messo in atto tramite il comportamento osservabile.
\end{snippetdefinition}
Per essere appreso, un cambiamento nel comportamento deve mantenersi relativamente costante nelle diverse circostanze. Ad esempio, una persona che si allena ogni giorno nel tiro con l'arco diventerà molto abile, ma se interrompe la pratica le sue competenze possono regredire fino al livello di partenza.
L'apprendimento è quindi un processo:
\begin{itemize}
    \item che richiede una pratica relativamente stabile e duratura;
    \item continuo, perché si apprende in ogni situazione, in modo deliberato o inconsapevole;
    \item che riguarda sia le conoscenze procedurali sia quelle dichiarative.
\end{itemize}
\subsection{Abituazione e sensibilizzazione}
\begin{snippetexample}{Abituazione}
    Immaginiamo di osservare una fotografia piacevole, come l'immagine di un cucciolo. Alla prima osservazione potremmo provare un'emozione abbastanza intensa. Se però osserviamo la stessa immagine molte volte in un breve intervallo di tempo, la risposta emotiva diventa progressivamente più debole.
\end{snippetexample}
\begin{snippetdefinition}{Abituazione}
    L'\emph{abituazione} è il processo per cui, quando uno stimolo viene presentato ripetutamente, si osserva una diminuzione della risposta comportamentale.
\end{snippetdefinition}
L'abituazione è un meccanismo adattivo perché aiuta a focalizzare l'attenzione sugli eventi nuovi presenti nell'ambiente. Permette inoltre di contenere lo sforzo comportamentale impiegato per rispondere più volte a stimoli già incontrati.
Il processo di abituazione si adatta alla definizione di apprendimento: vi è un cambiamento nel comportamento, perché la risposta emotiva diventa più debole; tale cambiamento è basato sull'esperienza, cioè sull'osservazione ripetuta dell'immagine; inoltre, esso è relativamente stabile, poiché non si ritorna immediatamente al livello originario di risposta emotiva.
\begin{snippetdefinition}{Sensibilizzazione}
    La \emph{sensibilizzazione} è il processo per cui la risposta a uno stimolo presentato ripetutamente diventa più forte invece che più debole.
\end{snippetdefinition}
\begin{snippetexample}{Sensibilizzazione}
    Se una persona viene sottoposta allo stesso stimolo doloroso molte volte in un breve intervallo di tempo, può provare più dolore durante l'ultima somministrazione rispetto alla prima, anche se l'intensità fisica dello stimolo rimane costante.
\end{snippetexample}
In generale, la sensibilizzazione è più probabile quando gli stimoli sono intensi o fastidiosi, mentre negli altri casi è più frequente il meccanismo dell'abituazione.
\subsection{Condizionamento classico}
\begin{snippetexample}{Film horror}
    Guardando un film horror per la seconda volta, il cuore può iniziare a battere più forte quando la colonna sonora suggerisce che sta per accadere qualcosa di pericoloso. Il corpo ha imparato a produrre una risposta fisiologica, come l'aumento del battito cardiaco, quando un evento situazionale, come una musica carica di tensione, è associato a un altro evento, come uno stimolo visivo spaventoso.
\end{snippetexample}
\begin{snippetdefinition}{Condizionamento classico}
    Il \emph{condizionamento classico} è una forma base di apprendimento in cui uno stimolo o evento predice il verificarsi di un secondo stimolo o evento. In questo modo, l'organismo apprende una nuova associazione tra due stimoli.
\end{snippetdefinition}
\subsubsection{L'esperimento di Pavlov}
\emph{Ivan Pavlov} stava conducendo una ricerca sulla digestione nei cani quando si imbatté nel condizionamento classico, scoperta per cui ricevette il premio Nobel nel 1904.
Pavlov aveva sviluppato una tecnica per studiare i processi della digestione nei cani inserendo alcune fistole nelle ghiandole salivari e negli organi digestivi, in modo da deviare le secrezioni corporee in contenitori esterni e poterle misurare e analizzare.
Per attivare queste secrezioni, gli assistenti di Pavlov nutrivano i cani con della carne. Dopo aver ripetuto la procedura diverse volte, Pavlov osservò che i cani iniziavano a salivare prima di essere nutriti. La salivazione compariva alla semplice vista del cibo, alla vista dell'assistente che lo portava o persino al rumore dei suoi passi.
Ogni stimolo che precedeva con regolarità la comparsa di cibo finiva quindi per attivare la produzione di saliva.
Durante gli esperimenti, i cani venivano collocati in imbracature che impedivano loro di muoversi. A intervalli regolari, uno stimolo acustico precedeva la presentazione di cibo. Inizialmente lo stimolo acustico era neutro, perché non aveva alcun significato precedente per il cane in relazione al cibo o alla salivazione. Dopo ripetute associazioni tra stimolo acustico e cibo, il solo stimolo acustico era in grado di produrre salivazione.
\begin{center}

    \includegraphics[width=0.72\textwidth]{media_psychology/pavlov-classical-conditioning-experiment.jpg}
    \captionof{figure}{Schema del condizionamento classico nell'esperimento di Pavlov.}
\end{center}
Il punto di partenza del condizionamento classico è costituito dai \emph{riflessi}, cioè risposte innescate naturalmente da stimoli specifici rilevanti per l'organismo.
\begin{snippetdefinition}{Stimolo incondizionato}
    Lo \emph{stimolo incondizionato}, o \emph{SI}, è uno stimolo che attiva un riflesso sul piano fisiologico. Il comportamento di risposta si verifica in presenza dello stimolo senza necessità di apprendimento.
\end{snippetdefinition}
\begin{snippetdefinition}{Risposta incondizionata}
    La \emph{risposta incondizionata}, o \emph{RI}, è il comportamento attivato naturalmente da uno stimolo incondizionato.
\end{snippetdefinition}
Prima del condizionamento, lo stimolo incondizionato attiva fisiologicamente la risposta incondizionata. Uno stimolo neutro, come un suono, non ha invece alcuna capacità di attivazione in relazione alla salivazione.
Durante il condizionamento, lo stimolo neutro viene abbinato allo stimolo incondizionato. Attraverso questa associazione, lo stimolo neutro diventa uno \emph{stimolo condizionato} e attiva una \emph{risposta condizionata} simile alla risposta incondizionata.
\begin{snippetdefinition}{Stimolo condizionato}
    Lo \emph{stimolo condizionato}, o \emph{SC}, è uno stimolo inizialmente neutro che diventa capace di attivare un comportamento dopo essere stato associato regolarmente a uno stimolo incondizionato.
\end{snippetdefinition}
\begin{snippetdefinition}{Risposta condizionata}
    La \emph{risposta condizionata}, o \emph{RC}, è un comportamento attivato non da uno stimolo fisiologicamente rilevante, ma da uno stimolo inizialmente neutro diventato rilevante tramite associazione con uno stimolo incondizionato.
\end{snippetdefinition}
In sintesi, la biologia fornisce le connessioni \texttt{SI-RI}, mentre l'apprendimento prodotto dal condizionamento classico crea le connessioni \texttt{SC-RC}.
\subsubsection{Processi di condizionamento}
I principali processi implicati nel condizionamento classico sono:
\begin{enumerate}
    \item \emph{acquisizione:} si forma l'associazione tra \texttt{SC} e \texttt{SI}. La risposta condizionata è inizialmente inattiva e aumenta gradualmente in frequenza in seguito a prove ripetute;
    \item \emph{estinzione:} viene presentato lo stimolo condizionato senza lo stimolo incondizionato, e la risposta condizionata diminuisce fino ad azzerarsi;
    \item \emph{riposo:} dopo un breve periodo di pausa, la risposta condizionata può ricomparire anche in assenza dello stimolo incondizionato;
    \item \emph{riacquisizione:} la nuova presentazione dello stimolo condizionato insieme allo stimolo incondizionato riattiva rapidamente la risposta.
\end{enumerate}
\begin{snippetdefinition}{Recupero spontaneo}
    Il \emph{recupero spontaneo} è la ricomparsa di una risposta condizionata dopo un periodo di riposo, anche in assenza dello stimolo incondizionato.
\end{snippetdefinition}
\begin{center}

    \includegraphics[width=0.58\textwidth]{media_psychology/conditioning-acquisition-extinction-reacquisition.jpg}
    \captionof{figure}{Andamento della risposta condizionata nelle fasi di acquisizione, estinzione, riposo e riacquisizione.}
\end{center}
Nel condizionamento classico la \emph{contiguità temporale}, o \emph{timing}, è fondamentale: lo stimolo condizionato e lo stimolo incondizionato devono essere presentati in tempi abbastanza vicini perché l'organismo li percepisca come associati.
I ricercatori hanno individuato quattro modalità principali di presentazione dei due stimoli:
\begin{itemize}
    \item \emph{condizionamento ritardato:} lo stimolo condizionato compare prima dello stimolo incondizionato e rimane fino a quando quest'ultimo viene presentato;
    \item \emph{condizionamento di traccia:} lo stimolo condizionato è presentato in modo discontinuo oppure scompare prima che compaia lo stimolo incondizionato;
    \item \emph{condizionamento simultaneo:} stimolo condizionato e stimolo incondizionato sono presentati nello stesso momento;
    \item \emph{condizionamento retrogrado:} lo stimolo condizionato è presentato dopo lo stimolo incondizionato.
\end{itemize}
\begin{center}

    \includegraphics[width=0.78\textwidth]{media_psychology/conditioning-timing-procedures.jpg}
    \captionof{figure}{Modalità temporali di presentazione di stimolo condizionato e stimolo incondizionato.}
\end{center}
Supponiamo di aver addestrato un cane ad associare uno stimolo acustico alla comparsa di cibo. La risposta del cane sarà specifica solo per quel determinato tipo di stimolo? In generale, una volta che una risposta condizionata è stata associata a un determinato stimolo condizionato, stimoli simili possono attivare la medesima risposta.
\begin{snippetdefinition}{Generalizzazione dello stimolo}
    La \emph{generalizzazione dello stimolo} è l'estensione della risposta condizionata a stimoli che non sono mai stati associati allo stimolo incondizionato originario, ma che risultano simili allo stimolo condizionato iniziale.
\end{snippetdefinition}
Quanto più uno stimolo è simile allo stimolo condizionato iniziale, tanto più forte sarà la risposta attivata.
\begin{center}

    \includegraphics[width=0.45\textwidth]{media_psychology/stimulus-generalization-curves.jpg}
    \captionof{figure}{Curve di generalizzazione dello stimolo.}
\end{center}
L'esistenza dei gradienti di generalizzazione mostra il modo in cui il condizionamento classico svolge la sua funzione nell'esperienza quotidiana. Poiché in natura gli stimoli significativi raramente si presentano ogni volta nello stesso identico modo, la generalizzazione dello stimolo produce un margine di flessibilità che estende la portata dell'apprendimento oltre la specifica esperienza di partenza.
In alcune circostanze, però, è importante che una risposta venga attivata soltanto da una serie limitata di stimoli. Un organismo, ad esempio, non dovrebbe sprecare energie fuggendo da animali solo superficialmente simili ai suoi predatori naturali.
\begin{snippetdefinition}{Discriminazione dello stimolo}
    La \emph{discriminazione dello stimolo} è il processo attraverso cui un organismo impara a rispondere in modo diverso a stimoli che sono distinti dallo stimolo condizionato per qualche dimensione.
\end{snippetdefinition}
Per un organismo che deve agire in modo ottimale nell'ambiente, generalizzazione e discriminazione devono trovare un equilibrio.
\subsubsection{Il processo di acquisizione}
Pavlov riteneva che il condizionamento classico derivasse dalla semplice contiguità temporale tra stimolo condizionato e stimolo incondizionato. Questa teoria dominò il condizionamento classico fino a quando \emph{Robert Rescorla}, nel 1966, dimostrò l'importanza della \emph{contingenza} nell'associazione \texttt{SC-SI}.
Rescorla condusse un esperimento con cani, utilizzando uno stimolo acustico come stimolo condizionato e una scossa elettrica come stimolo incondizionato. Per un gruppo di animali, stimolo condizionato e stimolo incondizionato erano solo temporalmente contigui; per l'altro gruppo, invece, il segnale acustico era un predittore affidabile della scossa.
Nella prima fase dell'esperimento, i cani vennero addestrati a saltare oltre una barriera da una parte all'altra della gabbia per evitare la scossa elettrica prodotta dalla griglia sul pavimento. Se non saltavano ricevevano la scossa; se saltavano, la scossa veniva rimandata. Rescorla utilizzò la frequenza con cui i cani saltavano la barriera come misura della \emph{paura condizionata}.
Nelle prime fasi del comportamentismo, \emph{Watson} e \emph{Rayner} cercarono di dimostrare che molte risposte basate sulla paura potevano essere comprese a partire dalla connessione tra uno stimolo neutro e uno stimolo incondizionato capace di provocare naturalmente paura. Per verificare questa idea condussero un esperimento su un bambino noto come il \emph{piccolo Albert}.
\begin{snippetexample}{Piccolo Albert}
    Watson e Rayner addestrarono Albert ad avere paura di un ratto bianco, che inizialmente era uno stimolo gradito, accoppiando la sua comparsa con uno stimolo incondizionato negativo, cioè un forte rumore. La risposta incondizionata di paura e lo stress emotivo prodotto dal rumore costituirono la base dell'apprendimento di una reazione di paura in presenza del ratto bianco.
\end{snippetexample}
Con questo esperimento Watson voleva dimostrare che risposte emotive come paura e ansia non sono necessariamente il prodotto di processi inconsci o di esperienze molto precoci, come sosteneva Freud. Al contrario, egli riteneva che le paure potessero essere apprese come qualunque altro comportamento.
Un ulteriore requisito necessario affinché uno stimolo possa fungere da base per il condizionamento classico è che esso sia \emph{informativo} rispetto all'ambiente e sufficientemente unico. Questo spiega perché il condizionamento compare più rapidamente quando lo stimolo condizionato risalta rispetto ai molti altri stimoli presenti nell'ambiente.
\begin{snippetdefinition}{Salienza dello stimolo}
    La \emph{salienza dello stimolo} indica il grado con cui uno stimolo viene notato rapidamente, in funzione della sua intensità e del contrasto con altri stimoli presenti nell'ambiente.
\end{snippetdefinition}
\subsection{Condizionamento operante}
\begin{snippetexample}{Film horror e comportamento futuro}
    Dopo aver visto un film horror, una persona potrebbe dire: ``ho imparato che non devo andare al cinema a vedere film horror''. In questo caso l'apprendimento non riguarda solo l'associazione tra due stimoli, ma il modo in cui una conseguenza modifica la probabilità di ripetere un comportamento.
\end{snippetexample}
\subsubsection{La legge dell'effetto}
Nello stesso periodo in cui Pavlov scopriva il condizionamento classico, \emph{Thorndike} osservava il comportamento dei gatti che cercavano di uscire da una gabbia-problema.
I gatti inizialmente si dibattevano contro la situazione di prigionia. Quando alcune azioni impulsive permisero loro di aprire la porta, gli impulsi privi di effetto vennero progressivamente abbandonati, mentre l'azione che aveva portato al successo venne fissata grazie al piacere prodotto dall'esito.
\begin{center}

    \includegraphics[width=0.45\textwidth]{media_psychology/thorndike-problem-box.jpg}
    \captionof{figure}{Gabbia-problema utilizzata negli studi di Thorndike.}
\end{center}
L'apprendimento derivava dall'associazione tra lo stimolo fornito dalla situazione e la risposta che l'animale aveva imparato a dare: una connessione \texttt{S-R}, cioè \emph{stimolo-risposta}. I gatti avevano imparato a produrre una risposta adeguata che, in quella situazione, portava a un esito desiderato.
L'apprendimento avveniva gradualmente e in modo automatico attraverso un meccanismo di \emph{prove ed errori}: le azioni che avevano prodotto un esito soddisfacente aumentavano progressivamente di frequenza fino a diventare la risposta dominante.
\begin{snippetdefinition}{Legge dell'effetto}
    La \emph{legge dell'effetto} afferma che le risposte seguite da conseguenze soddisfacenti tendono a essere ripetute, mentre quelle prive di effetto o seguite da conseguenze sfavorevoli tendono a diminuire.
\end{snippetdefinition}
\subsubsection{Analisi sperimentale del comportamento}
\emph{Skinner} condivideva la prospettiva di Thorndike secondo cui le conseguenze ambientali esercitano un forte effetto sul comportamento.
Egli mise a punto un programma di ricerca finalizzato a scoprire in che modo varie condizioni ambientali influenzano la probabilità che una data risposta si verifichi.
\begin{snippetdefinition}{Condizionamento operante}
    Il \emph{condizionamento operante} è una procedura attraverso cui si manipolano le conseguenze del comportamento di un organismo per valutarne l'effetto sul comportamento successivo.
\end{snippetdefinition}
\begin{snippetdefinition}{Comportamento operante}
    Un \emph{comportamento operante} è qualunque comportamento messo in atto dall'organismo capace di produrre effetti osservabili sull'ambiente. Non è attivato da riflessi innati come quelli del condizionamento classico.
\end{snippetdefinition}
Il condizionamento operante modifica la probabilità del verificarsi di diversi comportamenti operanti in funzione delle conseguenze che essi producono sull'ambiente.
Per portare avanti la sua analisi sperimentale, Skinner inventò la \emph{Skinner box}, o \emph{camera operante}.
\begin{center}

    \includegraphics[width=0.47\textwidth]{media_psychology/skinner-box-operant-chamber.jpg}
    \captionof{figure}{Skinner box, o camera operante, utilizzata per studiare il condizionamento operante.}
\end{center}
Quando un ratto, dopo aver prodotto un comportamento adeguato definito dallo sperimentatore, preme una leva, il meccanismo fornisce una porzione di cibo. Questo strumento permette allo sperimentatore di studiare le variabili che portano i ratti ad apprendere il comportamento desiderato.
Ad esempio, se premere la leva produce una ricompensa in cibo solo dopo che il ratto ha girato in tondo, il ratto imparerà rapidamente a girare in tondo prima di premere la leva. Questa metodologia permise a Skinner di studiare l'effetto del rinforzo sul comportamento animale.
\subsubsection{Il meccanismo del rinforzo}
Il meccanismo del rinforzo si genera grazie a un'associazione ripetuta tra una risposta e il cambiamento che essa produce nell'ambiente.
\begin{snippetexample}{Piccione e cibo}
    Immaginiamo un esperimento in cui la beccata di un piccione su un bottone è generalmente seguita dall'emissione di cibo. Questa relazione costante, cioè il rinforzo, sarà solitamente accompagnata da un aumento della frequenza delle beccate. Affinché il cibo aumenti la probabilità delle beccate, deve essere contingente solo a quella risposta e non ad altri comportamenti, come girare su se stessi o chinarsi.
\end{snippetexample}
\begin{snippetdefinition}{Rinforzo}
    Il \emph{rinforzo} è qualunque stimolo che, somministrato in modo contingente a una risposta, aumenta la probabilità che quella risposta si verifichi.
\end{snippetdefinition}
Esistono tre classi di stimoli:
\begin{itemize}
    \item stimoli \emph{neutri};
    \item stimoli \emph{appetitivi}, cioè piacevoli;
    \item stimoli \emph{aversivi}, cioè spiacevoli e da evitare.
\end{itemize}
\begin{snippetdefinition}{Rinforzo positivo}
    Il \emph{rinforzo positivo} si verifica quando un comportamento è seguito dall'erogazione di uno stimolo piacevole, aumentando la probabilità che quel comportamento si ripresenti.
\end{snippetdefinition}
\begin{snippetdefinition}{Rinforzo negativo}
    Il \emph{rinforzo negativo} si verifica quando un comportamento è seguito dalla rimozione, riduzione o evitamento di uno stimolo avversivo, aumentando la probabilità che quel comportamento si ripresenti.
\end{snippetdefinition}
Nel rinforzo negativo si distinguono due circostanze di apprendimento:
\begin{itemize}
    \item \emph{condizionamento operante di fuga:} l'organismo impara una risposta che permette di scappare da uno stimolo avversivo, come usare un ombrello durante un acquazzone per evitare di bagnarsi;
    \item \emph{condizionamento operante di evitamento:} l'organismo impara risposte che permettono di evitare uno stimolo spiacevole prima che compaia, come allacciare la cintura di sicurezza prima che l'auto emetta il segnale acustico.
\end{itemize}
\begin{snippetnote}{Rinforzo positivo e rinforzo negativo}
    Sia il rinforzo positivo sia il rinforzo negativo aumentano la probabilità della risposta che li precede. Il rinforzo positivo lo fa aggiungendo uno stimolo piacevole, mentre il rinforzo negativo lo fa rimuovendo, riducendo o prevenendo uno stimolo spiacevole.
\end{snippetnote}
Un'altra tecnica che modifica il comportamento è la \emph{punizione}, che porta alla diminuzione della probabilità che una risposta si verifichi.
\begin{snippetdefinition}{Stimolo punitivo}
    Uno \emph{stimolo punitivo} è uno stimolo che, somministrato in modo contingente a una risposta, diminuisce la probabilità che quella risposta si verifichi.
\end{snippetdefinition}
\begin{snippetdefinition}{Punizione positiva}
    La \emph{punizione positiva} si verifica quando un comportamento è seguito dalla somministrazione di uno stimolo spiacevole. È detta positiva perché aggiunge qualcosa alla situazione.
\end{snippetdefinition}
\begin{snippetexample}{Punizione positiva}
    Bere un caffè bollente produce dolore sulla punta della lingua; in futuro, quindi, sarà più probabile aspettare prima di berlo.
\end{snippetexample}
\begin{snippetdefinition}{Punizione negativa}
    La \emph{punizione negativa} si verifica quando un comportamento è seguito dalla rimozione di uno stimolo piacevole. È detta negativa perché sottrae qualcosa alla situazione.
\end{snippetdefinition}
\begin{snippetexample}{Punizione negativa}
    Se un genitore toglie la paghetta al figlio dopo che ha litigato con il fratellino, il bambino può imparare a non comportarsi più in quel modo.
\end{snippetexample}
Punizioni e rinforzi sono operazioni strettamente legate, ma devono essere distinti sulla base degli effetti che hanno sul comportamento: la punizione riduce la probabilità che una risposta si verifichi nuovamente, mentre il rinforzo la aumenta.
In genere non si vuole modificare la probabilità di un comportamento in tutte le situazioni, ma solo all'interno di un particolare contesto. Ad esempio, un insegnante può voler aumentare la probabilità che un bambino stia seduto tranquillo durante la lezione, senza modificare la probabilità che sia attivo e vivace durante l'intervallo.
\begin{snippetdefinition}{Stimoli discriminativi}
    Gli \emph{stimoli discriminativi} sono stimoli che precedono una particolare risposta e creano il contesto in cui quel comportamento produrrà probabilmente un certo effetto sull'ambiente.
\end{snippetdefinition}
Gli organismi imparano che, in presenza di alcuni stimoli ma non di altri, il loro comportamento produrrà probabilmente una particolare conseguenza.
\begin{snippetdefinition}{Contingenza a tre termini}
    La \emph{contingenza a tre termini}, secondo Skinner, è la sequenza composta da \emph{stimolo discriminativo}, \emph{comportamento} e \emph{conseguenza}. Essa rappresenta un concetto fondamentale delle teorie comportamentiste.
\end{snippetdefinition}
\begin{center}

    \includegraphics[width=0.65\textwidth]{media_psychology/three-term-contingency.jpg}
    \captionof{figure}{Esempi di contingenza a tre termini in diversi tipi di comportamento.}
\end{center}
\subsubsection{Proprietà dei rinforzi}
I rinforzi sono i mediatori più potenti del condizionamento operante: possono cambiare o mantenere invariato il comportamento. Essi possiedono proprietà interessanti e complesse.
Si possono distinguere due tipi di rinforzi:
\begin{itemize}
    \item \emph{rinforzi primari:} rinforzi come cibo e acqua, le cui proprietà rinforzanti sono biologicamente determinate;
    \item \emph{rinforzi condizionati:} stimoli inizialmente neutri che, associandosi ai rinforzi primari, diventano capaci di fungere da rinforzi per le risposte operanti.
\end{itemize}
In teoria, ogni stimolo può diventare un rinforzo condizionato se viene associato a un rinforzo primario.
Gli sperimentatori considerano i rinforzi condizionati più efficaci e semplici da usare rispetto ai rinforzi primari per diversi motivi:
\begin{itemize}
    \item i rinforzi primari immediatamente disponibili in una situazione sperimentale sono pochi;
    \item i rinforzi condizionati possono essere dispensati con rapidità;
    \item i rinforzi condizionati sono facilmente replicabili;
    \item l'effetto rinforzante dei rinforzi condizionati è più immediato, perché dipende dalla percezione di averli ricevuti e non da processi biologici.
\end{itemize}
\subsubsection{Schemi di rinforzo}
In alcuni casi non è possibile o non è utile rinforzare un comportamento ogni volta che viene emesso.
\begin{snippetexample}{Skinner e il rinforzo parziale}
    Il giovane Skinner rimase isolato in laboratorio senza cibo sufficiente per i suoi animali. Per risparmiare, iniziò a fornire le ricompense ai ratti solo dopo un certo intervallo di tempo: non importava quante volte premessero la leva, non ottenevano altro cibo prima del momento opportuno. Con questo schema di rinforzo parziale, i ratti mantennero un tasso di risposta simile a quello ottenuto con il rinforzo continuo e, durante l'estinzione, continuarono a rispondere più a lungo e intensamente rispetto ai ratti rinforzati continuamente.
\end{snippetexample}
La scoperta di Skinner portò allo studio sistematico degli effetti dei diversi schemi di rinforzo sul comportamento umano. Nella vita quotidiana si sperimentano spesso schemi di rinforzo: ad esempio, durante una lezione, alzando la mano il professore a volte dà la parola e altre volte no.
Nella vita reale i rinforzi possono essere somministrati:
\begin{itemize}
    \item secondo uno \emph{schema a rapporto}, cioè dopo un certo numero di risposte;
    \item secondo uno \emph{schema a intervallo}, cioè in seguito alla prima risposta dopo un intervallo di tempo specifico.
\end{itemize}
In entrambi i casi, lo schema può essere \emph{fisso} oppure \emph{variabile}.
\begin{snippetdefinition}{Schema a rapporto fisso}
    Nello \emph{schema a rapporto fisso}, o \emph{RF}, il rinforzo viene somministrato dopo che l'organismo ha prodotto un numero prestabilito di risposte.
\end{snippetdefinition}
Quando il rinforzo segue ogni singola risposta si parla di \texttt{RF-1}, cioè di rinforzo continuo. Quando, ad esempio, viene erogato ogni venticinque risposte, lo schema è definito \texttt{RF-25}. Questo tipo di schema produce tassi di risposta molto alti perché esiste una correlazione diretta tra risposta e rinforzo. Gli schemi a rapporto fisso generano una pausa dopo ogni rinforzo: più alto è il rapporto, più lunga sarà la pausa.
\begin{snippetexample}{Rapporto fisso}
    Un venditore deve vendere un certo numero di prodotti prima di ricevere il compenso. Il rinforzo dipende quindi dal raggiungimento di un numero fisso di risposte.
\end{snippetexample}
\begin{snippetdefinition}{Schema a rapporto variabile}
    Nello \emph{schema a rapporto variabile}, o \emph{RV}, il numero medio di risposte tra un rinforzo e il successivo è predeterminato, ma il rinforzo può arrivare dopo un numero variabile di risposte.
\end{snippetdefinition}
Uno schema \texttt{RV-10} significa che, in media, il rinforzo viene somministrato ogni \(10\) risposte, ma potrebbe arrivare dopo la prima o dopo la ventesima. Questo tipo di schema produce il tasso di risposta più alto e la maggiore resistenza all'estinzione, soprattutto quando il rapporto è alto.
\begin{snippetexample}{Rapporto variabile}
    Il gioco d'azzardo è governato da schemi a rapporto variabile. La risposta di inserire gettoni nella slot-machine viene mantenuta alta e costante dalla vincita, che viene erogata dopo un numero sconosciuto e variabile di gettoni inseriti.
\end{snippetexample}
\begin{snippetdefinition}{Schema a intervallo fisso}
    Nello \emph{schema a intervallo fisso}, o \emph{IF}, il rinforzo viene erogato in seguito alla prima risposta fornita dopo un intervallo di tempo costante e prestabilito.
\end{snippetdefinition}
In uno schema \texttt{IF-10}, dopo aver ricevuto il rinforzo, il soggetto deve attendere \(10\) secondi prima che un'altra risposta possa essere rinforzata, indipendentemente dal numero di risposte prodotte. Il tasso di risposta segue un andamento a merlo: subito dopo ogni risposta rinforzata l'animale risponde meno, mentre all'avvicinarsi del rinforzo successivo le risposte aumentano progressivamente.
\begin{snippetdefinition}{Schema a intervallo variabile}
    Nello \emph{schema a intervallo variabile}, o \emph{IV}, la media dell'intervallo tra i rinforzi è determinata, ma il rinforzo viene somministrato dopo intervalli variabili.
\end{snippetdefinition}
In uno schema \texttt{IV-20}, ad esempio, i rinforzi sono forniti a un tasso medio di uno ogni \(20\) secondi. Questo schema genera un tasso di risposta moderato ma costante.
\begin{center}

    \includegraphics[width=0.74\textwidth]{media_psychology/reinforcement-schedules-ratio.jpg}
    \captionof{figure}{Schemi di rinforzo a rapporto fisso e variabile.}
\end{center}
\begin{center}

    \includegraphics[width=0.58\textwidth]{media_psychology/reinforcement-schedules-interval.jpg}
    \captionof{figure}{Schemi di rinforzo a intervallo fisso e variabile.}
\end{center}
\subsubsection{Modellamento}
Nei paragrafi precedenti si è parlato di ratti che premono una leva per ottenere cibo. Tuttavia, anche premere una leva è un comportamento appreso. Quando un ratto viene inserito nella camera operante, è molto improbabile che prema spontaneamente la leva. Per addestrarlo a produrre un comportamento nuovo e complesso si utilizza il \emph{modellamento per approssimazioni successive}.
\begin{snippetdefinition}{Modellamento}
    Il \emph{modellamento}, o \emph{shaping}, è un metodo di apprendimento in cui si rinforza ogni risposta che si avvicina progressivamente al comportamento desiderato, fino a ottenere la risposta finale.
\end{snippetdefinition}
Il processo di modellamento consiste nel definire che cosa costituisce un progresso verso il risultato desiderato e nell'utilizzare un \emph{rinforzo differenziale} per perfezionare ogni passaggio lungo il percorso.
\begin{snippetdefinition}{Rinforzo differenziale}
    Il \emph{rinforzo differenziale} consiste nel rinforzare selettivamente le risposte che si avvicinano al comportamento-obiettivo, trascurando quelle meno adeguate.
\end{snippetdefinition}
Il condizionamento classico e il condizionamento operante sono le modalità fondamentali attraverso cui si è espresso l'approccio dell'\emph{apprendimento associativo}.
\begin{snippetdefinition}{Apprendimento associativo}
    L'\emph{apprendimento associativo} è una teoria meccanicistica secondo cui l'apprendimento viene misurato come variazione delle risposte comportamentali dopo una situazione di stimolazione.
\end{snippetdefinition}
In questo approccio, l'organismo è rappresentato come essenzialmente \emph{passivo} nel processo di apprendimento.

\subsection{Imprinting e apprendimento}
L'apprendimento è spesso un \emph{processo sociale} fondato sull'interazione fra due o più organismi. Fra i diversi fenomeni osservati in ambito etologico, quello probabilmente più noto è l'\emph{imprinting}.
\begin{snippetdefinition}{Imprinting}
    L'\emph{imprinting} è una forma di apprendimento precoce osservabile in animali appena nati che, come i pulcini, mostrano una reazione di inseguimento verso il primo oggetto mobile che vedono o sentono.
\end{snippetdefinition}
Attraverso l'imprinting, i piccoli stabiliscono un forte legame sociale, cioè un legame di \emph{attaccamento}, con il primo oggetto mobile incontrato nelle prime ore di vita, di norma la madre.
L'oca di \emph{Konrad Lorenz} è passata alla storia per aver riconosciuto come madre e seguito il famoso etologo austriaco. Con i suoi esperimenti sugli uccelli, Lorenz dimostrò il fenomeno dell'imprinting, interpretandolo come una forma di apprendimento innato che funziona soltanto entro un \emph{periodo critico}.
\begin{center}

    \includegraphics[width=0.55\textwidth]{media_psychology/lorenz-imprinting-geese.jpg}
    \captionof{figure}{Lorenz seguito dalle oche durante gli studi sull'imprinting.}
\end{center}
\begin{snippetdefinition}{Periodo critico}
    Il \emph{periodo critico} è un periodo breve e circoscritto che corrisponde alle prime fasi dello sviluppo del neonato, durante il quale alcune forme di apprendimento avverrebbero in modo particolarmente efficace.
\end{snippetdefinition}
Attraverso studi condotti in laboratorio, \emph{Vallortigara} e colleghi hanno criticato la tesi di Lorenz, sostenendo che questo tipo di apprendimento non avviene né rapidamente né in modo irreversibile. Inoltre, sembra che l'imprinting non sia circoscrivibile a un vero e proprio periodo critico.
Per questo motivo, il concetto di periodo critico è stato rinominato \emph{periodo sensibile}.
\begin{snippetdefinition}{Periodo sensibile}
    Il \emph{periodo sensibile} è il periodo in cui le influenze ambientali sono più efficaci per l'apprendimento di conoscenze e abilità.
\end{snippetdefinition}
Questo concetto è stato esteso, almeno in parte, anche alla specie umana da \emph{Bowlby}, il quale ha osservato che il secondo semestre di vita è un periodo sensibile per lo sviluppo del legame di attaccamento del bambino nei confronti della madre.
\subsection{Apprendimento cognitivo}
Nell'opera \emph{La doppia elica}, del 1968, Watson racconta come lui e Crick riuscirono a scoprire la forma del DNA grazie a un'improvvisa intuizione, cioè un \emph{insight}, dopo anni di tentativi ed errori.
Fenomeni di questo tipo costituiscono una difficoltà per il comportamentismo, perché sono manifestazioni difficili da spiegare nei termini del condizionamento pavloviano o skinneriano. Molti psicologi ritengono che queste intuizioni dipendano da un processo differente, chiamato \emph{apprendimento cognitivo}.
\begin{snippetdefinition}{Apprendimento cognitivo}
    L'\emph{apprendimento cognitivo} è una forma di apprendimento spiegata in termini di cambiamenti nei processi mentali, piuttosto che soltanto in termini di cambiamenti comportamentali osservabili.
\end{snippetdefinition}
Secondo la psicologia cognitiva, alcune forme di apprendimento devono essere interpretate considerando i cambiamenti che avvengono nelle rappresentazioni e nei processi mentali dell'individuo.
\subsubsection{Apprendimento per insight}
\emph{Wolfgang Köhler} propose una nuova prospettiva per spiegare l'apprendimento. Egli riteneva che la psicologia dovesse riconoscere i processi mentali come una componente essenziale dell'apprendimento, al contrario di quanto sostenevano i comportamentisti.
In una serie di studi, Köhler mostrò che gli scimpanzé potevano imparare a risolvere problemi complessi non solo per prove ed errori, ma anche attraverso improvvise intuizioni capaci di combinare risposte più semplici precedentemente apprese.
\begin{snippetexample}{Sultan}
    Uno degli esperimenti di Köhler coinvolse \emph{Sultan}, uno scimpanzé che aveva imparato sia ad ammucchiare scatole e arrampicarsi su di esse per raggiungere frutta appesa nella gabbia, sia a usare un bastone per prendere frutta fuori dalla sua portata. Quando venne posto in una nuova situazione che richiedeva di combinare i due apprendimenti, Sultan agì inizialmente per prove ed errori, fallendo. Dopo un momento di apparente riflessione, trascinò una cassa e un bastone sotto la frutta, salì sulla cassa e raggiunse il cibo con il bastone.
\end{snippetexample}
\begin{center}

    \includegraphics[width=0.58\textwidth]{media_psychology/kohler-sultan-insight-learning.jpg}
    \captionof{figure}{Esperimento di Köhler sull'apprendimento per insight nello scimpanzé Sultan.}
\end{center}
\begin{snippetdefinition}{Apprendimento per insight}
    L'\emph{apprendimento per insight} è una forma di apprendimento che avviene attraverso una riorganizzazione improvvisa della percezione del problema in una data situazione.
\end{snippetdefinition}
Il comportamento di Sultan mostrava che gli animali non si limitano a utilizzare meccanicamente risposte condizionate, ma possono apprendere attraverso una ristrutturazione cognitiva del problema.
\subsubsection{Le mappe cognitive}
Anche i ratti nel laboratorio di \emph{Tolman} mostrarono comportamenti difficili da spiegare secondo la dottrina comportamentista. Essi riuscivano a esplorare i labirinti come se seguissero una sorta di \emph{mappa mentale}, invece di eseguire meccanicamente una serie di comportamenti appresi.
\begin{snippetdefinition}{Mappa cognitiva}
    Una \emph{mappa cognitiva} è un'immagine mentale che un organismo utilizza per muoversi all'interno di un ambiente familiare.
\end{snippetdefinition}
Chiunque abbia camminato al buio nella propria casa può intuire cosa Tolman intendesse con l'espressione \emph{mappa cognitiva}: ci si orienta nello spazio non solo grazie a risposte meccaniche, ma grazie a una rappresentazione mentale dell'ambiente.
Secondo Tolman, la mappa cognitiva è l'unico modo per spiegare come un ratto riesca a selezionare rapidamente un percorso alternativo in un labirinto quando la strada abituale verso l'obiettivo è bloccata. I ratti sceglievano spesso la deviazione più breve intorno a una barriera, anche se quel percorso non era stato precedentemente rinforzato.
\begin{center}

    \includegraphics[width=0.48\textwidth]{media_psychology/tolman-cognitive-map-maze.jpg}
    \captionof{figure}{Schema del labirinto utilizzato per studiare le mappe cognitive nei ratti.}
\end{center}
I soggetti preferivano la via diretta, cioè il percorso \(1\), quando era accessibile. Se veniva introdotta la barriera \(A\), preferivano il percorso \(2\). Quando invece veniva introdotta la barriera \(B\), i ratti sceglievano di solito il percorso \(3\).
In uno studio ulteriore, Tolman permise ai ratti di vagare liberamente all'interno di un labirinto per molte ore. Durante questo periodo, i ratti non ricevettero rinforzi di alcun genere, ma si limitarono a esplorare il labirinto. In seguito, nonostante l'assenza di rinforzi precedenti, questi ratti impararono a muoversi nel labirinto per cercare cibo più rapidamente rispetto ad altri ratti che non lo avevano mai esplorato.
\begin{snippetdefinition}{Apprendimento latente}
    L'\emph{apprendimento latente} è una forma di apprendimento che avviene senza manifestarsi immediatamente nel comportamento e senza rinforzi evidenti, ma che può emergere in seguito quando diventa utile.
\end{snippetdefinition}
Anche se durante la prima esperienza non era possibile individuare un indice osservabile di apprendimento, i ratti avevano comunque imparato qualcosa.
\subsection{Apprendimento situato, simulazione ed esperienza}
L'apprendimento avviene sempre in un contesto fisico e relazionale in cui \emph{esperto} e \emph{novizio} sono entrambi protagonisti nella costruzione degli apprendimenti. Essi sono necessariamente interdipendenti, perché gli esperti insegnano in funzione di quanto i novizi sono in grado di assimilare, e viceversa.
\begin{snippetdefinition}{Apprendimento situato}
    L'\emph{apprendimento situato} è un apprendimento legato a una specifica situazione e immerso in un determinato contesto immediato.
\end{snippetdefinition}
Rispetto all'insegnamento-apprendimento tradizionale, come la lezione frontale e il rapporto ``da uno a molti'', l'apprendimento situato garantisce un'esperienza interattiva e aperta. Esso permette inoltre un'esperienza fortemente individualizzata, perché ogni fruitore partecipa secondo i propri atteggiamenti e le proprie aspettative.
L'apprendimento situato è anche \emph{contingente}, cioè legato a determinate circostanze e contesti che pongono inevitabilmente vincoli e opportunità. Inoltre, implica l'acquisizione di \emph{conoscenze tacite}.
\begin{snippetdefinition}{Conoscenza tacita}
    La \emph{conoscenza tacita} è una conoscenza individuale implicita, difficile da codificare, trasmettere e condividere con altri. È in larga misura inconsapevole ed è immersa nell'esperienza pratica.
\end{snippetdefinition}
La conoscenza tacita si fonda sul \emph{training} e tende a trasformarsi in abitudini automatiche o semiautomatiche. Durante un periodo di osservazione attenta, l'imitazione accurata di modelli esperti e una profonda interazione con loro consentono al novizio di appropriarsi di procedure, strategie e accorgimenti.
L'apprendimento situato e partecipativo facilita quindi l'\emph{apprendimento riflessivo}, il \emph{pensiero critico} e la \emph{partecipazione guidata} basata sulla condivisione dei significati.
\subsection{Nuove frontiere dell'apprendimento}
Con l'avvento delle nuove tecnologie, l'apprendimento ha notevolmente allargato i propri orizzonti, assumendo forme inimmaginabili fino a pochi decenni fa. In particolare, oggi si parla di \emph{apprendimento elettronico}, o \emph{e-learning}.
\begin{snippetdefinition}{E-learning}
    L'\emph{e-learning} è una forma di apprendimento che sceglie la tecnologia come strumento privilegiato.
\end{snippetdefinition}
\subsubsection{Serious Games}
La \emph{mente simulativa} è una disposizione generale che consente agli individui di riprodurre e anticipare fenomeni dell'esperienza, al fine di intervenire in modo opportuno secondo quanto richiesto dalla situazione.
\begin{snippetdefinition}{Serious Games}
    I \emph{Serious Games} sono attività digitali interattive che, attraverso la simulazione virtuale, consentono ai partecipanti di fare esperienze precise e accurate, promuovendo attraverso il gioco percorsi attivi, partecipati e coinvolgenti di apprendimento.
\end{snippetdefinition}
I Serious Games rappresentano una manifestazione innovativa e fondamentale della mente simulativa.
Nel corso dell'esistenza si manifesta un legame profondo tra attività di gioco e serietà dell'apprendimento. Gioco e simulazione mentale costituiscono infatti un intreccio indissolubile con l'apprendimento fin dall'infanzia, quando ci si appropria di molte competenze e abilità attraverso varie forme di gioco.
Gli elementi principali dei Serious Games sono:
\begin{itemize}
    \item la presenza di un \emph{obiettivo sfidante};
    \item il fatto di essere \emph{coinvolgenti} e divertenti da usare;
    \item l'attribuzione di un \emph{punteggio};
    \item la possibilità di fornire al partecipante abilità, atteggiamenti o conoscenze applicabili al mondo reale.
\end{itemize}
\begin{center}
    \includegraphics[width=0.58\textwidth]{media_psychology/serious-games-learning.jpg}
    \captionof{figure}{Esempio di Serious Games applicati all'apprendimento.}
\end{center}
\subsection{Valutazione dell'apprendimento}
La \emph{valutazione} è un momento essenziale e critico per ogni forma di apprendimento. È necessario verificare e misurare quanto il discente ha imparato e in che modo. Senza questo passaggio non è possibile accertare i progressi realizzati dalle persone o fare confronti tra gruppi diversi di allievi.
Nei sistemi informativi tradizionali, la valutazione viene generalmente svolta attraverso strumenti come l'interrogazione, il colloquio e diversi tipi di test.
Tali strumenti devono rispettare criteri di \emph{attendibilità}, \emph{validità}, \emph{accuratezza} e \emph{finezza discriminativa}. Quando questi criteri sono rispettati, si ottiene una misura standardizzata delle competenze prese in considerazione.
Nonostante i vantaggi, i metodi tradizionali di valutazione presentano il rischio di distorsioni dovute ai limiti dell'attendibilità, della validità, degli stati emotivi e degli atteggiamenti degli individui.
I Serious Games rappresentano potenzialmente una rivoluzione anche nei processi di valutazione. La valutazione può infatti diventare \emph{dinamica}, fornendo al fruitore un feedback in tempo reale mentre sta imparando.
Oltre a poter integrare tecniche tradizionali, i Serious Games sono in grado di offrire forme più complesse e complete di misurazione, come la \emph{valutazione in itinere}.
\section{La memoria}
Provate a pensare al vostro primo ricordo. A quando risale? Quanto è vivida la scena che riuscite a ricordare? Questa immagine è stata influenzata dal ricordo che altre persone conservano dello stesso evento?
Immaginate cosa accadrebbe se non aveste alcun ricordo del passato: se non ricordaste il vostro decimo compleanno o l'ultima serata passata con gli amici. Senza questi ancoraggi temporali, sarebbe difficile attribuire un senso a ciò che vi circonda e alla vostra stessa identità.
Supponete poi di aver perso la capacità di acquisire nuovi ricordi. Le esperienze più recenti svanirebbero come se non fossero mai accadute e diventerebbe difficile seguire il filo di una conversazione o la trama di una serie televisiva.
\subsection{Che cos'è la memoria?}
\begin{snippetdefinition}{Memoria}
    La \emph{memoria} è la capacità di codificare, immagazzinare e recuperare le informazioni.
\end{snippetdefinition}
La memoria può essere descritta come una particolare modalità di elaborazione delle informazioni. Poiché normalmente svolge il proprio compito senza che ce ne accorgiamo, tendiamo a darla per scontata, come accade per processi corporei quali la digestione o la respirazione.
Tuttavia, prestiamo attenzione alla memoria quando qualcosa non funziona: ad esempio quando dimentichiamo le chiavi della macchina, una ricorrenza importante o la risposta a una domanda d'esame. Il cervello umano è in grado di memorizzare mediamente l'equivalente di \(100\) miliardi di bit di informazione, mostrando una capacità di gestione delle informazioni estremamente ampia.
\subsubsection{Funzioni della memoria}
Quando si pensa alla memoria, vengono in mente soprattutto le situazioni in cui la si usa per ricordare eventi specifici o informazioni, come il proprio film preferito o il numero di matricola.
Una delle funzioni principali della memoria è consentire un accesso consapevole agli eventi passati, personali e collettivi. Tuttavia, la memoria svolge anche altre funzioni: garantisce senza sforzo una percezione di continuità rispetto alle proprie esperienze. Ad esempio, quando si guida nel proprio quartiere, la memoria rende familiari i negozi situati ai lati della strada.
\paragraph{Memoria implicita ed esplicita}
Osservate una scena in cui un coniglio si trova in cucina. Il coniglio appare immediatamente fuori luogo, ma questa sensazione non deriva da un controllo consapevole di ogni singolo oggetto presente nell'immagine. Piuttosto, i processi mnestici confrontano automaticamente la scena attuale con le conoscenze passate relative a ciò che dovrebbe trovarsi in una cucina.
\begin{center}

    \includegraphics[width=0.48\textwidth]{media_psychology/implicit-memory-out-of-context.jpg}
    \captionof{figure}{Esempio di uso implicito della memoria nel riconoscimento di un elemento fuori contesto.}
\end{center}
Questo esempio permette di comprendere la differenza tra uso implicito e uso esplicito della memoria.
\begin{snippetdefinition}{Uso esplicito della memoria}
    L'\emph{uso esplicito della memoria} comporta uno sforzo consapevole per codificare e recuperare le informazioni.
\end{snippetdefinition}
\begin{snippetdefinition}{Uso implicito della memoria}
    L'\emph{uso implicito della memoria} si verifica quando la codifica o il recupero delle informazioni avvengono senza che ce ne rendiamo conto.
\end{snippetdefinition}
La scoperta del coniglio è un esempio di memoria implicita, perché i processi mnestici confrontano la conoscenza passata relativa alla cucina con l'interpretazione dell'immagine senza uno sforzo consapevole. Se invece venisse chiesto che cosa manca nella cucina, sarebbe necessario usare la memoria esplicita, recuperando consapevolmente ciò che di solito è presente in quell'ambiente.
La stessa distinzione può essere applicata alla fase iniziale di immagazzinamento dei ricordi. Alcune conoscenze su ciò che dovrebbe trovarsi in una cucina vengono acquisite senza sforzo consapevole; altre, come l'associazione tra la parola \emph{frigorifero} e l'oggetto corrispondente, richiedono processi mnestici espliciti.
\paragraph{Memoria dichiarativa e procedurale}
Gli esempi relativi alla memoria implicita ed esplicita riguardano spesso il recupero di fatti o eventi, cioè la \emph{memoria dichiarativa}. Esiste però anche una memoria relativa al modo in cui si fanno le cose, chiamata \emph{memoria procedurale}.
\begin{snippetdefinition}{Memoria dichiarativa}
    La \emph{memoria dichiarativa} riguarda il recupero di fatti o eventi che possono essere dichiarati verbalmente o rappresentati consapevolmente.
\end{snippetdefinition}
\begin{snippetdefinition}{Memoria procedurale}
    La \emph{memoria procedurale} riguarda l'acquisizione, il mantenimento e l'uso delle abilità e delle procedure con cui si fanno le cose. È una forma implicita di memoria.
\end{snippetdefinition}
Gli studi sulla memoria procedurale riguardano spesso l'intensità e la durata del periodo di apprendimento: si cerca di capire come si passi da un insieme consapevole di azioni, messe in atto durante una certa attività, allo svolgimento inconsapevole e automatico della stessa attività.
Questi fenomeni possono essere osservati anche in un'attività semplice, come digitare un numero di telefono che nel tempo è diventato familiare.
\subsubsection{I processi mnestici}
La capacità di fare ricorso alle esperienze pregresse richiede l'intervento di tre processi differenti:
\begin{itemize}
    \item la \emph{codifica};
    \item l'\emph{immagazzinamento};
    \item il \emph{recupero}.
\end{itemize}
\begin{snippetdefinition}{Codifica}
    La \emph{codifica} è la fase iniziale dell'elaborazione delle informazioni, che porta alla costruzione di una rappresentazione della memoria.
\end{snippetdefinition}
\begin{snippetdefinition}{Immagazzinamento}
    L'\emph{immagazzinamento} consiste nella conservazione nel tempo del materiale codificato.
\end{snippetdefinition}
\begin{snippetdefinition}{Recupero}
    Il \emph{recupero} è il reperimento delle informazioni immagazzinate in un momento successivo.
\end{snippetdefinition}
In altre parole, la codifica permette l'ingresso delle informazioni, l'immagazzinamento le conserva e il recupero le estrae quando sono necessarie.
Nello specifico, la codifica richiede la costruzione di \emph{rappresentazioni mentali} di informazioni relative al mondo esterno.
\begin{snippetexample}{Rappresentare un oggetto}
    Se dovete descrivere a un amico il regalo più bello ricevuto al vostro compleanno, ma non lo avete con voi, potete descriverne le proprietà, fare un disegno oppure mimarne l'uso. In tutti questi casi costruite rappresentazioni interne dell'oggetto originario, che sostituiscono la realtà fornendo alcune informazioni su di essa e tralasciandone altre.
\end{snippetexample}
\begin{snippetdefinition}{Rappresentazione mentale}
    Una \emph{rappresentazione mentale} è una rappresentazione interna che sta al posto di un oggetto o evento e trasmette informazioni congruenti con esso.
\end{snippetdefinition}
Le rappresentazioni mentali conservano le caratteristiche salienti delle esperienze passate, permettendo di rievocarne il ricordo.

\subsection{Utilizzo della memoria nel breve periodo}
La maggior parte delle informazioni che incontriamo non rimane conservata stabilmente nella memoria. Al contrario, spesso le informazioni sono disponibili solo per un breve periodo.
In questa sezione si considerano tre forme di utilizzo non permanente della memoria:
\begin{itemize}
    \item la \emph{memoria iconica};
    \item la \emph{memoria a breve termine};
    \item la \emph{memoria di lavoro}.
\end{itemize}
\subsubsection{Memoria iconica}
\begin{snippetdefinition}{Memoria iconica}
    La \emph{memoria iconica}, o \emph{memoria sensoriale visiva}, è il sistema di memoria relativo al dominio visivo che permette di immagazzinare una grande quantità di informazioni per un intervallo di tempo molto breve.
\end{snippetdefinition}
La memoria iconica è un esempio di \emph{memoria sensoriale}. I ricercatori hanno ipotizzato che ogni sistema sensoriale possieda un deposito di memoria capace di conservare per pochi secondi le rappresentazioni delle caratteristiche fisiche degli stimoli ambientali.
Un ricordo visivo, o \emph{icona}, dura circa mezzo secondo.
\begin{snippetexample}{Esperimento di Sperling}
    George Sperling presentò ai partecipanti una serie di tre righe di lettere per un tempo brevissimo, attraverso una presentazione tachistoscopica. Nella procedura di \emph{recupero totale}, i partecipanti dovevano elencare tutti gli elementi che riuscivano a ricordare. Nella procedura di \emph{recupero parziale}, invece, dovevano riportare soltanto il contenuto di una riga, indicata da un segnale acustico alto, medio o basso.
\end{snippetexample}
\begin{center}

    \includegraphics[width=0.45\textwidth]{media_psychology/sperling-iconic-memory-matrix.jpg}
    \captionof{figure}{Esempio di matrice di lettere utilizzata negli esperimenti di Sperling sulla memoria iconica.}
\end{center}
Poiché i partecipanti riuscivano a riportare accuratamente qualunque delle tre righe fosse indicata dal segnale acustico, Sperling concluse che tutte le informazioni presentate erano state acquisite dalla memoria iconica.
Da un lato, questi risultati dimostrano che la memoria iconica può immagazzinare una grande quantità di informazioni; dall'altro, la differenza tra recupero totale e recupero parziale mostra che l'informazione svanisce molto rapidamente. I partecipanti assegnati al recupero totale, infatti, non riuscivano a ricordare tutte le informazioni presenti nell'icona.
Questa conclusione venne rafforzata usando un segnale acustico leggermente ritardato: in questo caso il numero di elementi ricordati diminuiva progressivamente.
\begin{snippetnote}{Memoria iconica e immagine eidetica}
    La memoria iconica non coincide con la cosiddetta memoria fotografica. Il termine tecnico per indicare la memoria fotografica è \emph{immagine eidetica}: chi possiede immagini eidetiche riesce a ricordare i dettagli di un'immagine per un tempo molto più lungo della memoria iconica, come se stesse ancora osservando la fotografia.
\end{snippetnote}
\subsubsection{Memoria a breve termine}
\begin{snippetdefinition}{Memoria a breve termine}
    La \emph{memoria a breve termine}, o \emph{MBT}, consente di focalizzare le proprie risorse cognitive su un numero limitato di rappresentazioni mentali per un periodo temporaneo.
\end{snippetdefinition}
La memoria a breve termine permette, ad esempio, di ricordare un numero telefonico abbastanza a lungo da poterlo digitare dopo averlo cercato in rubrica. Le sue risorse sono instabili e la capacità di mantenere più informazioni attive è limitata.
Per verificare la capacità massima della memoria a breve termine, i ricercatori hanno elaborato il \emph{test dello span di memoria}, o \emph{memory span}, spesso applicato alle cifre.
\begin{snippetdefinition}{Span di memoria}
    Lo \emph{span di memoria} è la quantità di elementi che una persona riesce a mantenere temporaneamente in memoria e a riprodurre nell'ordine corretto.
\end{snippetdefinition}
\begin{snippetexample}{Span di cifre}
    Si legge una lista di numeri casuali una sola volta, poi si copre la lista e si prova a scrivere le cifre ricordate nello stesso ordine in cui sono state presentate. Ad esempio: \texttt{8 1 7 3 4 9 4 2 8 5}.
\end{snippetexample}
La maggior parte delle persone riesce a ricordare un numero di elementi compreso tra \(5\) e \(9\). \emph{George Miller} affermò che \(7\), più o meno \(2\), è il ``numero magico'' che caratterizza le prestazioni di memoria quando le persone devono ricordare liste di lettere, parole, numeri o altri elementi familiari e dotati di significato.
Nonostante la capacità limitata della MBT, la codifica delle informazioni può essere rafforzata attraverso due processi:
\begin{itemize}
    \item la \emph{ripetizione}, o \emph{rehearsal};
    \item la creazione di \emph{unità di informazione}, o \emph{chunk}.
\end{itemize}
\begin{snippetdefinition}{Ripetizione}
    La \emph{ripetizione} è un processo che impedisce all'informazione di essere eliminata dalla memoria a breve termine, mantenendola attiva attraverso la reiterazione.
\end{snippetdefinition}
\begin{snippetdefinition}{Chunk}
    Un \emph{chunk}, o \emph{unità di informazione}, è un'organizzazione di più bit di informazione in una singola unità dotata di significato.
\end{snippetdefinition}
\begin{snippetdefinition}{Chunking}
    Il \emph{chunking} è il processo di riconfigurazione degli elementi attraverso il raggruppamento basato su somiglianza, significato o altri principi di organizzazione, a partire da informazioni già presenti nella memoria a lungo termine.
\end{snippetdefinition}
\begin{snippetexample}{Chunking}
    La sequenza \texttt{1-9-8-4} è composta da quattro cifre che potrebbero occupare molto spazio nella memoria a breve termine. Se però viene interpretata come l'anno \texttt{1984}, essa forma una singola unità di informazione, lasciando più spazio disponibile per altri elementi.
\end{snippetexample}
È quindi possibile strutturare l'informazione secondo significati personali oppure combinare nuovi stimoli con codici già memorizzati nella memoria a lungo termine.
\begin{snippetexample}{Date come chunk}
    Nella sequenza \texttt{2001200620142021} si possono individuare \(16\) cifre indipendenti oppure \(4\) unità, se la sequenza viene spezzata in anni: \texttt{2001}, \texttt{2006}, \texttt{2014}, \texttt{2021}. Nel secondo caso è più facile memorizzare la sequenza corretta dopo una rapida occhiata.
\end{snippetexample}
\subsubsection{Memoria di lavoro}
\begin{snippetdefinition}{Memoria di lavoro}
    La \emph{memoria di lavoro} è una risorsa mnestica che permette di mantenere e manipolare temporaneamente le informazioni necessarie per svolgere compiti come ragionare, comprendere il linguaggio e risolvere problemi.
\end{snippetdefinition}
Immaginiamo di cercare di ricordare un numero di telefono mentre si cerca un foglio e una matita per scriverlo. La memoria a breve termine permette di tenere il numero in mente, mentre la memoria di lavoro permette di compiere le operazioni mentali necessarie a cercare gli oggetti.
La memoria di lavoro concerne il presente: connette il passato con il presente, integra il presente con i ricordi ed è responsabile sia della memorizzazione sia dell'elaborazione dell'informazione.
Secondo \emph{Alan Baddeley}, la memoria di lavoro è composta da quattro componenti:
\begin{itemize}
    \item \emph{circuito fonologico:} consente di ricordare e manipolare informazioni acustiche basate sulla produzione verbale, come quando si ripete mentalmente un numero telefonico;
    \item \emph{taccuino visuospaziale:} svolge funzioni simili al circuito fonologico, ma relative all'informazione visiva e spaziale;
    \item \emph{esecutivo centrale:} controlla l'attenzione e coordina le informazioni provenienti dal circuito fonologico e dal taccuino visuospaziale;
    \item \emph{buffer episodico:} è un sistema di immagazzinamento a capacità limitata, controllato dall'esecutivo centrale, che permette di recuperare informazioni dalla memoria a lungo termine e combinarle con le informazioni della situazione attuale.
\end{itemize}
\subsection{Memoria a lungo termine: codifica e recupero}
Quando gli psicologi parlano di \emph{memoria a lungo termine}, si riferiscono ai ricordi che accompagnano una persona per molto tempo, talvolta per tutta la vita.
\begin{snippetdefinition}{Memoria a lungo termine}
    La \emph{memoria a lungo termine}, o \emph{MLT}, è il magazzino di tutte le esperienze, gli eventi, le informazioni e le capacità che una persona possiede su di sé e sul mondo, acquisite attraverso la memoria sensoriale e la memoria di lavoro.
\end{snippetdefinition}
\begin{center}

    \includegraphics[width=0.58\textwidth]{media_psychology/memory-systems-model.jpg}
    \captionof{figure}{Schema dei sistemi di memoria e delle relazioni tra memoria sensoriale, memoria di lavoro e memoria a lungo termine.}
\end{center}
\subsubsection{Indizi di recupero}
Per recuperare un ricordo si utilizzano \emph{indizi di recupero}. Essi possono provenire dall'ambiente esterno, come la domanda di un test, oppure essere generati internamente, come quando ci si chiede: ``dove ho già visto questa persona?''.
\begin{snippetdefinition}{Indizio di recupero}
    Un \emph{indizio di recupero} è un elemento contestuale che consente il recupero di ricordi immagazzinati nella memoria a lungo termine.
\end{snippetdefinition}
\paragraph{Recupero e riconoscimento}
\begin{snippetdefinition}{Recupero}
    Il \emph{recupero} è la tecnica di richiamo della memoria in cui si deve riprodurre un'informazione a cui si è stati precedentemente esposti.
\end{snippetdefinition}
\begin{snippetdefinition}{Riconoscimento}
    Il \emph{riconoscimento} è la realizzazione che un dato evento o stimolo è già stato esperito in passato.
\end{snippetdefinition}
\begin{snippetexample}{Testimonianza}
    Se la polizia chiede a una vittima di ricordare alcune caratteristiche distintive dell'aggressore, sta usando un metodo basato sul recupero. Se invece mostra alla vittima fotografie di vari sospetti, sta usando un metodo basato sul riconoscimento.
\end{snippetexample}
\paragraph{Memoria episodica e semantica}
La memoria dichiarativa assume forme diverse in funzione degli indizi necessari per il recupero. Lo psicologo canadese \emph{Endel Tulving} propose per primo la distinzione tra \emph{memoria episodica} e \emph{memoria semantica}.
\begin{snippetdefinition}{Memoria episodica}
    La \emph{memoria episodica} conserva eventi specifici di cui si fa esperienza personalmente. Per riattivare questi ricordi sono necessari indizi di contesto legati al momento in cui l'evento è accaduto.
\end{snippetdefinition}
\begin{snippetdefinition}{Memoria semantica}
    La \emph{memoria semantica} è una memoria generica e categoriale, relativa al significato delle parole, dei concetti, dei simboli e alle relazioni tra essi.
\end{snippetdefinition}
Il ricordo del primo bacio è un esempio di memoria episodica, perché è legato al contesto di apprendimento. Al contrario, conoscenze come la capitale della Francia o una costante fisica fanno parte della memoria semantica e, per la maggior parte delle persone, non richiedono indizi di contesto specifici per essere richiamate alla mente.
\subsubsection{Contesto e codifica}
\begin{snippetdefinition}{Specificità del contesto di codifica}
    Il principio della \emph{specificità del contesto di codifica} afferma che i ricordi si attivano più rapidamente quando il contesto di recupero è coerente con il contesto di codifica.
\end{snippetdefinition}
\begin{snippetexample}{Riconoscimento fuori contesto}
    Può capitare di vedere in una stanza affollata una persona che si sa di aver già incontrato, senza riuscire a ricordare dove. Dopo averla osservata più a lungo, la si riconosce: la difficoltà nasceva dal fatto che non era mai stata vista in quel particolare contesto.
\end{snippetexample}
I ricercatori hanno fornito diverse dimostrazioni della \emph{memoria contesto-dipendente}. In un esperimento, alcuni subacquei dovevano memorizzare una lista di parole sulla spiaggia oppure sott'acqua. La prestazione era migliore di circa il \(50\%\) quando il contesto di codifica e quello di recupero coincidevano.
\begin{snippetdefinition}{Memoria contesto-dipendente}
    La \emph{memoria contesto-dipendente} è il fenomeno per cui il recupero di un'informazione è facilitato quando il contesto in cui viene recuperata è simile a quello in cui è stata codificata.
\end{snippetdefinition}
I cambiamenti di contesto aiutano a spiegare anche l'\emph{effetto di posizione seriale}.
\begin{snippetdefinition}{Effetto di posizione seriale}
    L'\emph{effetto di posizione seriale} è il fenomeno per cui, in una lista, vengono ricordati meglio i primi elementi e gli ultimi elementi, mentre gli elementi intermedi sono ricordati con maggiore difficoltà.
\end{snippetdefinition}
L'effetto di posizione seriale comprende:
\begin{itemize}
    \item l'\emph{effetto priorità}, cioè il miglior ricordo dei primi item;
    \item l'\emph{effetto recenza}, cioè il miglior ricordo degli ultimi item.
\end{itemize}
\begin{center}

    \includegraphics[width=0.62\textwidth]{media_psychology/serial-position-curve.jpg}
    \captionof{figure}{Curva di posizione seriale: gli item iniziali e finali sono ricordati meglio degli item intermedi.}
\end{center}
Il ruolo del contesto nella curva di posizione seriale dipende dalla \emph{differenziazione contestuale}, cioè dalla diversità del contesto in cui le informazioni sono state apprese.
\begin{snippetdefinition}{Differenziazione contestuale}
    La \emph{differenziazione contestuale} è la diversità del contesto in cui sono state apprese le informazioni, che può renderle più o meno distinguibili e quindi più o meno facili da recuperare.
\end{snippetdefinition}
\begin{center}

    \includegraphics[width=0.58\textwidth]{media_psychology/contextual-differentiation.jpg}
    \captionof{figure}{Rappresentazione visiva della differenziazione contestuale.}
\end{center}
L'effetto recenza si verifica perché gli ultimi elementi vengono differenziati quasi automaticamente. Lo stesso principio può spiegare l'effetto priorità: ogni volta che si inizia qualcosa di nuovo, i processi mnestici creano un nuovo contesto, nel quale le prime esperienze risultano particolarmente differenziate.
\subsubsection{I processi di codifica e di recupero}
La memoria funziona in modo più efficace quando le fasi di codifica e di recupero sono ben abbinate. Il modo in cui un'informazione viene elaborata nel momento della codifica influenza la qualità del suo ricordo.
\begin{snippetdefinition}{Teoria dei livelli di elaborazione}
    La \emph{teoria dei livelli di elaborazione} sostiene che quanto più è profondo il livello a cui un'informazione viene elaborata, tanto più è probabile che essa venga recuperata dalla memoria.
\end{snippetdefinition}
La profondità dell'elaborazione è spesso definita dal tipo di giudizio richiesto ai partecipanti rispetto a un materiale sperimentale.
\begin{snippetexample}{Livelli di elaborazione}
    Consideriamo la parola \texttt{CANE}. Si può osservare che è scritta in maiuscolo, che termina con \texttt{-ne}, oppure che indica un animale a quattro zampe. Ognuna di queste osservazioni richiede un livello di elaborazione via via più profondo.
\end{snippetexample}
\begin{snippetdefinition}{Priming}
    Il \emph{priming} è un fenomeno per cui uno stimolo precedente influenza la risposta a uno stimolo successivo, anche in assenza di una correlazione evidente tra i due.
\end{snippetdefinition}
\subsubsection{Perché dimentichiamo}
Nella maggior parte dei casi la memoria funziona in modo adeguato, ma talvolta possono emergere dimenticanze. \emph{Ebbinghaus}, nel 1885, fu il primo ricercatore a documentare l'interferenza attraverso esperimenti rigorosi.
\begin{snippetdefinition}{Interferenza}
    L'\emph{interferenza} si verifica quando ricordi diversi competono tra loro e rendono più difficile l'acquisizione o il recupero di informazioni.
\end{snippetdefinition}
\begin{snippetdefinition}{Interferenza proattiva}
    L'\emph{interferenza proattiva} si verifica quando le informazioni acquisite in passato rendono difficile l'acquisizione di nuove informazioni.
\end{snippetdefinition}
\begin{snippetdefinition}{Interferenza retroattiva}
    L'\emph{interferenza retroattiva} si verifica quando l'acquisizione di nuove informazioni rende difficile il ricordo di informazioni memorizzate in precedenza.
\end{snippetdefinition}
\begin{snippetexample}{Cambio di numero di telefono}
    Quando si cambia numero di telefono, all'inizio il vecchio numero può rendere difficile memorizzare quello nuovo: è interferenza proattiva. Dopo aver ripetuto più volte il numero nuovo, può invece diventare difficile ricordare quello vecchio: è interferenza retroattiva.
\end{snippetexample}
\subsubsection{Come memorizzare informazioni non strutturate}
Una strategia per migliorare la codifica delle informazioni è il \emph{ripasso elaborativo}.
\begin{snippetdefinition}{Ripasso elaborativo}
    Il \emph{ripasso elaborativo} è una tecnica di memorizzazione che consiste nell'arricchire la codifica di una nuova informazione collegandola ad altro materiale significativo.
\end{snippetdefinition}
\begin{snippetexample}{Cane e tetto}
    Per ricordare la coppia \texttt{Cane-Tetto}, si può costruire un'immagine mentale di un cane finito sul tetto per recuperare una pallina. In questo modo il recupero della coppia viene agevolato.
\end{snippetexample}
\begin{snippetdefinition}{Mnemotecniche}
    Le \emph{mnemotecniche} sono strumenti che permettono di codificare una lunga serie di fatti associandoli a informazioni familiari già codificate in precedenza.
\end{snippetdefinition}
\begin{snippetdefinition}{Metodo dei loci}
    Il \emph{metodo dei loci} è una mnemotecnica che consente di ricordare l'ordine di una lista di nomi o oggetti associandoli a una sequenza di luoghi familiari.
\end{snippetdefinition}
\begin{snippetexample}{Metodo dei loci}
    Per ricordare la lista della spesa, si può posizionare mentalmente ogni prodotto lungo il percorso abituale da casa all'università. Il recupero avviene poi ripercorrendo mentalmente gli stessi luoghi nello stesso ordine.
\end{snippetexample}
\subsubsection{La metamemoria}
\begin{snippetdefinition}{Metamemoria}
    La \emph{metamemoria} è la capacità di riflettere sul funzionamento dei propri processi mnestici e sulle informazioni che si è certi di possedere.
\end{snippetdefinition}
Le domande di metamemoria riguardano, ad esempio, la sicurezza di possedere una certa informazione, l'attendibilità dei propri ricordi e la valutazione delle strategie di recupero utilizzate.
\subsubsection{Il processo di ricostruzione della memoria}
In molti casi non si recupera un'informazione direttamente, ma la si ricostruisce basandosi su più forme generali di conoscenza memorizzata.
\begin{snippetdefinition}{Memoria ricostruttiva}
    La \emph{memoria ricostruttiva} è un tipo di memoria che ricostruisce le informazioni basandosi su più forme generali di conoscenza memorizzata.
\end{snippetdefinition}
Poiché la memoria può ricostruire i ricordi, possono verificarsi \emph{distorsioni}, cioè circostanze in cui la memoria differisce dalla realtà.
\begin{snippetdefinition}{Distorsioni della memoria}
    Le \emph{distorsioni della memoria} sono circostanze in cui il ricordo ricostruito differisce dall'evento reale.
\end{snippetdefinition}
\emph{Bartlett} individuò tre processi di ricostruzione:
\begin{itemize}
    \item \emph{livellamento:} semplificazione della storia;
    \item \emph{modellamento:} enfatizzazione di alcuni dettagli;
    \item \emph{assimilazione:} cambiamento dei dettagli.
\end{itemize}

\begin{center}
    \includegraphics[width=0.52\textwidth]{media_psychology/reconstructive-memory-example.jpg}
    \captionof{figure}{Esempio di memoria ricostruttiva.}
\end{center}
\section{Il pensiero}
In questo capitolo ci si focalizza su alcuni \emph{processi cognitivi}. Le capacità di usare il linguaggio e di pensare in astratto sono state spesso considerate l'essenza dell'esperienza umana. Tendiamo a dare per scontata la cognizione perché è qualcosa in cui siamo continuamente impegnati durante le ore di veglia.
\begin{snippetdefinition}{Cognizione}
    La \emph{cognizione} è un termine generale che si riferisce a tutte le forme di conoscenza. Lo studio della cognizione coincide quindi con lo studio della vita mentale.
\end{snippetdefinition}
La cognizione comprende sia i \emph{contenuti} sia i \emph{processi}. I contenuti della cognizione sono ciò che si sa: concetti, fatti, regole e ricordi, ad esempio ``il cane è un mammifero''. I processi cognitivi sono invece il modo in cui questi contenuti mentali vengono elaborati allo scopo di interpretare il mondo circostante e trovare soluzioni creative ai problemi della vita quotidiana.
Gli psicologi cognitivi studiano le funzioni mentali superiori con particolare attenzione al modo in cui le persone acquisiscono conoscenze e le utilizzano per manipolare e comprendere la propria esperienza nel mondo.
\begin{center}

    \includegraphics[width=0.62\textwidth]{media_psychology/cognitive-psychology-domains.jpg}
    \captionof{figure}{Ambiti principali della psicologia cognitiva.}
\end{center}
\subsection{Studiare la cognizione}
\begin{snippetdefinition}{Psicologia cognitiva}
    La \emph{psicologia cognitiva} è il settore della psicologia che si occupa dello studio della cognizione e dei processi mentali.
\end{snippetdefinition}
La psicologia cognitiva è supportata in modo interdisciplinare dalle \emph{scienze cognitive}.
\begin{snippetdefinition}{Scienze cognitive}
    Le \emph{scienze cognitive} raccolgono e integrano il sapere proveniente da diverse discipline accademiche che si occupano delle medesime questioni teoriche relative alla cognizione.
\end{snippetdefinition}
La condivisione di dati e intuizioni costituisce un vantaggio per ciascuna delle scienze coinvolte nello studio della mente.
\begin{center}

    \includegraphics[width=0.48\textwidth]{media_psychology/cognitive-science-interdisciplinarity.jpg}
    \captionof{figure}{Interdisciplinarità delle scienze cognitive.}
\end{center}
Studiare la cognizione presenta una difficoltà evidente: essa si trova nel profondo della mente. Ad esempio, è possibile osservare l'\emph{input}, come un messaggio che dice ``chiamami'', ed è possibile osservare l'\emph{output}, cioè il fatto che una persona effettui una telefonata. Il problema è capire che cosa accade nel mezzo, cioè quali processi cognitivi e quali rappresentazioni mentali permettano di passare dal messaggio alla risposta.
\subsubsection{Scoprire i processi mentali}
Una delle metodologie più importanti per lo studio dei processi mentali fu ideata nel 1868 dallo psicologo tedesco \emph{Donders}. Egli voleva studiare la \emph{velocità dei processi mentali} e mise a punto una serie di compiti sperimentali che si differenziavano in base ai passaggi mentali necessari per la corretta esecuzione.
L'intuizione di Donders era che i \emph{passaggi mentali supplementari} potessero essere dedotti dal maggior tempo richiesto per eseguire un compito.
\begin{center}

    \includegraphics[width=0.78\textwidth]{media_psychology/donders-mental-processes.jpg}
    \captionof{figure}{Esempio di analisi dei processi mentali secondo la logica di Donders.}
\end{center}
\begin{snippetexample}{Compiti di Donders}
    Nel primo compito bisogna determinare se una lettera è maiuscola o minuscola e, se è maiuscola, tracciare una \texttt{C} sopra di essa. Nel secondo compito bisogna determinare se la lettera è maiuscola o minuscola, stabilire se la lettera maiuscola è una vocale o una consonante e poi tracciare una \texttt{V} sopra le vocali e una \texttt{C} sopra le consonanti.
\end{snippetexample}
Il secondo compito richiede almeno due passaggi mentali in più rispetto al primo: la \emph{categorizzazione dello stimolo}, cioè stabilire se la lettera sia vocale o consonante, e la \emph{selezione della risposta}, cioè decidere se tracciare una \texttt{C} o una \texttt{V}. Poiché per risolvere il secondo compito è necessario fare tutto ciò che si fa nel primo, e qualcosa in più, il tempo richiesto è maggiore.
Il terzo compito serve da controllo: è una procedura necessaria per verificare che la differenza di tempo non dipenda dal fatto che sia più veloce tracciare le \texttt{C} rispetto alle \texttt{V}.
\subsubsection{Processi mentali e risorse mentali}
Quando gli psicologi cognitivi scompongono attività di alto livello, come l'uso del linguaggio o il problem solving, nei processi che le compongono, spesso agiscono come se stessero giocando con mattoncini da costruzione. Ogni mattoncino rappresenta una diversa componente che deve essere eseguita. L'obiettivo è determinare la forma e l'ampiezza di ciascun componente e capire come essi si combinino per creare l'attività intera.
\begin{snippetdefinition}{Processi seriali}
    I \emph{processi seriali} sono processi mentali che si svolgono uno dopo l'altro, secondo una sequenza ordinata.
\end{snippetdefinition}
\begin{snippetdefinition}{Processi paralleli}
    I \emph{processi paralleli} sono processi mentali che si sovrappongono nell'unità di tempo e avvengono simultaneamente.
\end{snippetdefinition}
\begin{snippetexample}{Processi seriali e paralleli}
    Al ristorante, quando si decide che cosa ordinare, ci si può concentrare sulle portate una alla volta e giudicare per ciascuna se sia una scelta possibile: in questo caso il processo è seriale. Quando invece il cameriere arriva per prendere l'ordine, i processi linguistici che permettono di comprendere la domanda possono operare contemporaneamente ai processi che permettono di formulare una risposta: in questo caso i processi sono paralleli.
\end{snippetexample}
\begin{center}

    \includegraphics[width=0.68\textwidth]{media_psychology/serial-parallel-processes.jpg}
    \captionof{figure}{Rappresentazione di processi seriali, processi paralleli ed equivalenza temporale.}
\end{center}
Gli psicologi cognitivi usano spesso i tempi di reazione per stabilire se i processi mentali si svolgano in serie o in parallelo. Tuttavia non è sempre semplice determinarlo, perché il tempo totale necessario per completare un compito non permette di stabilire con certezza se due processi siano avvenuti uno dopo l'altro oppure simultaneamente.
In molti casi, gli studiosi valutano se i processi siano seriali o paralleli osservando quanto essi impegnano le \emph{risorse mentali}.
\begin{snippetdefinition}{Risorse di elaborazione}
    Le \emph{risorse di elaborazione} sono risorse cognitive limitate che devono essere distribuite tra i diversi compiti mentali in atto. La loro distribuzione dipende dai processi attentivi.
\end{snippetdefinition}
Non tutti i processi incidono nello stesso modo sulle risorse disponibili. È possibile distinguere una dimensione che va dai \emph{processi controllati} ai \emph{processi automatici}.
\begin{snippetdefinition}{Processi controllati}
    I \emph{processi controllati} richiedono attenzione e molte risorse cognitive. Nella maggior parte dei casi non è possibile svolgere più processi controllati contemporaneamente.
\end{snippetdefinition}
\begin{snippetdefinition}{Processi automatici}
    I \emph{processi automatici} non richiedono attenzione in modo rilevante e spesso possono essere eseguiti contemporaneamente ad altri processi automatici senza produrre interferenze.
\end{snippetdefinition}
\begin{snippetexample}{Camminare e conversare}
    In condizioni normali è semplice camminare lungo un marciapiede e conversare con un amico, perché la navigazione spaziale e il linguaggio possono svolgersi in parallelo. Se però il marciapiede è pieno di pozzanghere, diventa necessario scegliere attentamente il percorso: i processi spaziali richiedono risorse supplementari e i processi linguistici possono essere momentaneamente sospesi.
\end{snippetexample}
Un esempio di processo automatico riguarda il riconoscimento del valore numerico. Se viene chiesto di individuare tra due numeri quello fisicamente più grande, il compito diventa più difficile quando c'è incongruenza tra grandezza fisica e valore numerico. Ad esempio, se il \texttt{5} è scritto più grande del \texttt{7}, ma ha valore numerico minore, il riconoscimento automatico del valore numerico interferisce con il compito.
\begin{center}

    \includegraphics[width=0.52\textwidth]{media_psychology/numeric-stroop-interference.jpg}
    \captionof{figure}{Esempio di interferenza tra grandezza fisica e valore numerico.}
\end{center}
Questo compito mostra come i processi automatici facciano affidamento su un uso efficiente della memoria. Quando si vedono numeri come \texttt{5} o \texttt{7}, i processi mnestici forniscono rapidamente le informazioni relative alla quantità. Inoltre, compiti che inizialmente richiedevano processi controllati possono diventare automatici con la pratica.
\subsection{Le forme del pensiero}
È possibile distinguere tre forme principali di pensiero:
\begin{itemize}
    \item \emph{pensiero verbale:} pensiero strutturato in forma verbale;
    \item \emph{immagine mentale:} rappresentazione analogica, nella mente, del mondo esterno e interno;
    \item \emph{pensiero astratto:} pensiero formato da codici astratti senza l'uso di parole né immagini.
\end{itemize}
Il pensiero verbale e il pensiero per immagini sono forme note e facilmente riconoscibili del pensiero. Pensando con parole e numeri si può risolvere un problema di aritmetica; scegliendo tra immagini mentali di diversi percorsi, invece, si può decidere quale strada permetta di raggiungere una destinazione nel modo più agevole.
Negli ultimi decenni si è fatta strada anche l'ipotesi dell'esistenza di un pensiero formato da \emph{simboli amodali}, cioè un codice che non si identifica con nessuna lingua parlata e che non deriva da nessuna modalità sensoriale. Si tratterebbe di un codice in forma astratta e proposizionale.
\subsubsection{Il pensiero verbale}
Nella vita quotidiana, sia quando si è soli sia quando ci si relaziona con gli altri, una parte rilevante del pensiero è in forma verbale.
Secondo \emph{Watson}, fondatore del comportamentismo, il pensiero non era altro che \emph{linguaggio subvocale}, cioè quel tipo di linguaggio che si attiva quando pensiamo parlando sottovoce a noi stessi. Il pensiero veniva così ridotto a un comportamento osservabile, consistente nei minimi movimenti della muscolatura dell'apparato fonatorio.
\begin{snippetdefinition}{Linguaggio subvocale}
    Il \emph{linguaggio subvocale} è il linguaggio silenzioso che si attiva quando pensiamo come se parlassimo sottovoce a noi stessi.
\end{snippetdefinition}
La convinzione di Watson fu confutata da un esperimento di \emph{Smith} e colleghi. Smith assunse una sostanza capace di paralizzare completamente i muscoli striati e venne tenuto in vita da un respiratore artificiale. Terminato l'effetto della sostanza, riferì che, pur con tutti i muscoli paralizzati, riusciva comunque a pensare senza difficoltà. Da ciò si concluse che linguaggio parlato e pensiero sono dissociabili.
Dei rapporti tra pensiero e linguaggio verbale si sono occupati anche \emph{Vygotskij} e \emph{Piaget}. Per Vygotskij, il linguaggio infantile parte da radici sociali: il linguaggio egocentrico deriva, insieme al linguaggio comunicativo, dalla suddivisione funzionale del primo linguaggio del bambino, chiamato \emph{linguaggio sociale}. Secondo Vygotskij, il linguaggio egocentrico, separandosi dal linguaggio comunicativo, porterebbe nel tempo al \emph{linguaggio interiore}, centrale per l'attività di pensiero.
\subsubsection{Le immagini mentali}
\begin{snippetdefinition}{Immagine mentale}
    Un'\emph{immagine mentale} è una rappresentazione psicologica di percezioni e azioni, generata dalla corrispondente mappa cerebrale momentanea di una data situazione.
\end{snippetdefinition}
\begin{snippetexample}{Il gatto sul tappeto}
    Se bisogna scegliere tra due rappresentazioni visive della frase ``il gatto è sul tappeto'', dal punto di vista verbale entrambe possono sembrare corrette. Tuttavia, di norma si preferisce l'immagine in cui il gatto è orientato nel modo abituale, perché corrisponde alla scena che viene spontaneamente costruita nella mente.
\end{snippetexample}
\begin{center}

    \includegraphics[width=0.48\textwidth]{media_psychology/mental-image-cat-on-mat.jpg}
    \captionof{figure}{Esempio di immagine mentale evocata dalla frase ``il gatto è sul tappeto''.}
\end{center}
La storia è ricca di scoperte riconducibili alle immagini mentali. \emph{Albert Einstein}, ad esempio, sosteneva di aver pensato interamente in termini di immagini visive e di aver tradotto le proprie scoperte in simboli matematici e parole solo dopo la conclusione del lavoro visivo di scoperta.
Nella vita quotidiana si usano regolarmente immagini visive. Un classico esperimento chiedeva ai partecipanti di modificare mentalmente alcune immagini. Venivano presentati stimoli raffiguranti la lettera \texttt{R} oppure la sua immagine speculare, con una rotazione angolare compresa tra \(0^\circ\) e \(180^\circ\). Appena la lettera compariva, il partecipante doveva stabilire se fosse una \texttt{R} normale o speculare.
\begin{center}

    \includegraphics[width=0.50\textwidth]{media_psychology/mental-rotation-letter-r.jpg}
    \captionof{figure}{Compito di rotazione mentale con la lettera \texttt{R}.}
\end{center}
I tempi di reazione mostrarono che i partecipanti vedevano lo stimolo con l'\emph{occhio della mente} e facevano ruotare mentalmente l'immagine fino alla posizione verticale corretta, prima di decidere se si trattasse della \texttt{R} o della sua immagine speculare. Maggiore era l'angolo di rotazione, maggiore era il tempo di reazione.
I dati suggerirono che la rotazione mentale è molto simile alla rotazione fisica dello stimolo.
Le immagini visive vengono usate anche per rispondere a domande sul mondo. Ad esempio, se bisogna stabilire se una palla da golf sia più grande di una palla da ping-pong, si può formare un'immagine mentale delle due palle vicine e confrontarle.
I ricercatori si sono chiesti quanto coincidano le regioni cerebrali impegnate nella percezione visiva e quelle attivate durante l'immaginazione visiva. In un esperimento, i partecipanti apprendevano disegni di oggetti comuni e poi venivano sottoposti a risonanza magnetica funzionale sia mentre osservavano gli stessi disegni su uno schermo sia mentre producevano immagini mentali dei disegni.
\begin{center}

    \includegraphics[width=0.62\textwidth]{media_psychology/visual-perception-imagery-brain-areas.jpg}
    \captionof{figure}{Confronto tra aree cerebrali attivate nella percezione visiva e nell'immaginazione visiva.}
\end{center}
Da questi studi è emerso che aree simili si attivano sia nella percezione sia nella costruzione di immagini mentali.
Spesso ci formiamo immagini visive a partire da descrizioni verbali. Quando leggiamo, possiamo costruire un \emph{modello mentale spaziale} che tiene conto della posizione degli oggetti descritti. È possibile usare l'immaginazione mentale per proiettarsi all'interno di una scena: proprio come se fossimo realmente in una stanza, impiegheremmo meno tempo a dire ciò che si trova davanti a noi rispetto a ciò che si trova dietro.
\begin{center}

    \includegraphics[width=0.46\textwidth]{media_psychology/spatial-mental-model.jpg}
    \captionof{figure}{Esempio di modello mentale spaziale costruito a partire da una descrizione.}
\end{center}
\subsubsection{Il pensiero in forma astratta}
Si ipotizza l'esistenza di un pensiero formato da codici astratti e innati, non provenienti dalle diverse modalità sensoriali e propriocettive. Si tratterebbe di un pensiero in forma astratta proposizionale, costituito da proposizioni non composte né da parole né da immagini.
\emph{Fodor} ha definito questo sistema \emph{linguaggio della mente}.
\begin{snippetdefinition}{Linguaggio della mente}
    Il \emph{linguaggio della mente} è, secondo Fodor, un sistema di rappresentazioni astratte e proposizionali che non coincide con parole o immagini e che costituirebbe la base del pensiero.
\end{snippetdefinition}
Secondo questa prospettiva, le rappresentazioni del linguaggio della mente:
\begin{itemize}
    \item hanno parti costituenti che si combinano secondo una precisa logica, cioè una sintassi;
    \item sono composte da concetti innati che rievocano proprietà del mondo;
    \item sono composizionali, perché le proprietà semantiche di una rappresentazione complessa dipendono dalle proprietà semantiche degli elementi di base;
    \item sono regolate da relazioni di implicazione, ad esempio \emph{russare} implica \emph{dormire}.
\end{itemize}
Nella prospettiva di Fodor, la mente è un sistema chiuso, che non interagisce con l'ambiente esterno sul piano percettivo e motorio.
\begin{snippetnote}{Limite dell'ipotesi amodale}
    A oggi non è emersa un'evidenza neuropsicologica chiara dell'esistenza di simboli amodali nel cervello.
\end{snippetnote}
\subsection{Categorie e concetti}
La nostra esperienza del mondo è un flusso incessante, disomogeneo e casuale di sensazioni, movimenti, pensieri, emozioni e azioni. La mente umana non può governare ogni singolo pezzo dell'esperienza né controllare la sua globalità nel suo insieme. Essa procede invece a un livello intermedio, discernendo, frazionando e raggruppando elementi associabili tra loro per somiglianza, contiguità o altri criteri.
Siamo quindi costretti a suddividere l'interezza dell'esperienza in \emph{unità di informazione}, o \emph{chunk}, più o meno grandi.
\begin{snippetdefinition}{Categoria mentale}
    Una \emph{categoria mentale} è una classe di entità relativamente omogenee al proprio interno ed eterogenee rispetto alle entità di altre classi.
\end{snippetdefinition}
Ogni categoria raggruppa entità che possiedono alcune proprietà simili sulla base di criteri definiti e sufficientemente espliciti. Al tempo stesso, le entità di una categoria presentano differenze discriminanti rispetto alle entità delle altre categorie.
Le categorie danno senso al mondo in cui viviamo e, in quanto elaborazioni della mente, non sono determinate direttamente dall'ambiente: l'ambiente fornisce soltanto stimoli e indizi.
\begin{snippetdefinition}{Concetto}
    Un \emph{concetto} è la conoscenza che abbiamo di una categoria di oggetti o eventi.
\end{snippetdefinition}
Sul modo in cui i concetti si formano esistono varie ipotesi esplicative, nessuna delle quali completamente esaustiva.
Gli studi tradizionali, secondo l'\emph{approccio classico}, distinguono tra \emph{intensione} ed \emph{estensione} di un concetto.
\begin{snippetdefinition}{Intensione}
    L'\emph{intensione} di un concetto è l'insieme degli attributi necessari affinché un oggetto o evento faccia parte di quel concetto.
\end{snippetdefinition}
\begin{snippetdefinition}{Estensione}
    L'\emph{estensione} di un concetto comprende tutti gli oggetti o eventi che fanno parte del concetto stesso.
\end{snippetdefinition}
Secondo l'approccio classico, un concetto è caratterizzato da un insieme di attributi che costituiscono le unità di base con cui esso è costruito. Ciascun attributo è necessario e tutti gli attributi, insieme, sono sufficienti a individuare un membro qualsiasi di una categoria concettuale.
I concetti sono organizzati gerarchicamente: gli attributi che definiscono uno specifico concetto includono tutti gli attributi del concetto a esso sopraordinato.
Nel modello a rete gerarchica, ciascun nodo della rete corrisponde a un concetto e i legami che collegano i nodi rappresentano le relazioni tra concetti.
\begin{center}

    \includegraphics[width=0.56\textwidth]{media_psychology/hierarchical-concept-network.jpg}
    \captionof{figure}{Esempio di rete gerarchica dei concetti.}
\end{center}
L'idea secondo cui tutti gli esemplari di una determinata categoria concettuale la rappresentino in modo ugualmente efficace si scontra con alcune eccezioni. Ad esempio, lo \emph{struzzo} appartiene alla classe degli uccelli ed è quindi subordinato al nodo-concetto di \emph{uccello}, ma non eredita da esso la proprietà di volare, tipica di molti uccelli.
Le ricerche di \emph{Eleanor Rosch} propongono invece un'organizzazione categoriale in ottica tassonomica.
\begin{snippetdefinition}{Tassonomia}
    Una \emph{tassonomia} è un ordinamento sistematico. Nel caso delle categorie, indica una classificazione gerarchica delle categorie stesse.
\end{snippetdefinition}
Rosch ipotizza tre livelli gerarchici:
\begin{enumerate}
    \item \emph{livello sovraordinato}, ad esempio \emph{arredamento};
    \item \emph{livello di base}, ad esempio \emph{sedia};
    \item \emph{livello subordinato}, ad esempio \emph{sedia da cucina} o \emph{sedia a dondolo}.
\end{enumerate}
\begin{center}
    \includegraphics[width=0.60\textwidth]{media_psychology/category-taxonomy-levels.jpg}
    \captionof{figure}{Esempio di tassonomia categoriale.}
\end{center}
Le \emph{categorie di base} sono le più importanti. Gli oggetti che vi appartengono implicano un programma motorio unitario, somiglianze sul piano morfologico e un numero elevato di tratti comuni. Inoltre, le categorie di base sono le più rilevanti sul piano comunicativo, perché le parole che le designano sono le più comunemente utilizzate e vengono apprese prima dai bambini.
A livello orizzontale, le categorie mentali sono organizzate attorno al \emph{prototipo}.
\begin{snippetdefinition}{Prototipo}
    Il \emph{prototipo} è il miglior esemplare di una categoria, cioè quello che la rappresenta meglio perché possiede maggiore salienza.
\end{snippetdefinition}
All'interno di qualsiasi categoria, alcuni componenti possiedono le proprietà di quella categoria in modo più rilevante rispetto ad altri.
Secondo Rosch:
\begin{itemize}
    \item i prototipi sono gli elementi centrali attorno ai quali si organizza la categoria;
    \item l'appartenenza a una categoria è graduale e non dicotomica, in base al grado di somiglianza con i prototipi della categoria;
    \item più un elemento è simile al prototipo, più forte sarà la sua appartenenza alla categoria;
    \item gli elementi di una categoria sono simili tra loro, ma non possiedono proprietà perfettamente identiche.
\end{itemize}
Questa impostazione si basa sui principi di \emph{somiglianza} e \emph{analogia}, perché i vari componenti di una categoria vengono confrontati con il prototipo secondo giudizi di maggiore o minore somiglianza.
Questa concezione ha però sollevato diversi dubbi. Anzitutto, i concetti di \emph{rappresentatività} e \emph{appartenenza} vengono spesso sovrapposti, mentre in realtà indicano due processi distinti.
\begin{snippetdefinition}{Rappresentatività}
    La \emph{rappresentatività} indica il possesso del maggior numero di proprietà tipiche di una categoria.
\end{snippetdefinition}
\begin{snippetdefinition}{Appartenenza categoriale}
    L'\emph{appartenenza categoriale} indica il fatto che un elemento faccia effettivamente parte di una categoria.
\end{snippetdefinition}
I criteri di somiglianza con il prototipo non sono sufficienti per definire l'appartenenza a una categoria, perché possono risultare troppo vaghi. L'appartenenza categoriale richiede criteri più robusti, come il possesso di proprietà essenziali comuni a tutti i componenti.
\begin{snippetexample}{Categoria uccello}
    Per appartenere alla categoria \emph{uccello}, un animale deve essere oviparo e possedere il becco. Senza queste due proprietà essenziali, un animale rimane fuori dalla categoria, anche se somiglia per qualche aspetto a un uccello.
\end{snippetexample}

\subsubsection{Il modello esteso del prototipo}
Verso gli anni Novanta è stato elaborato il \emph{modello esteso del prototipo}: si passa dal prototipo come esemplare concreto al prototipo come \emph{costrutto mentale}, inteso come insieme di proprietà astratte.
Questo passaggio è essenziale per il processo di categorizzazione dell'esperienza e costituisce una condizione fondamentale per la comprensione delle categorie stesse.
\begin{snippetdefinition}{Modello esteso del prototipo}
    Il \emph{modello esteso del prototipo} considera il prototipo non come un esemplare concreto, ma come una configurazione astratta di proprietà salienti che distinguono una categoria da un'altra.
\end{snippetdefinition}
Il prototipo diventa la \emph{configurazione degli effetti prototipici}, cioè l'insieme delle proprietà più salienti che distinguono una categoria da un'altra. Esso assume il valore modale della categoria, poiché ne rappresenta il punto culminante.
Le proprietà di una categoria non costituiscono un insieme omogeneo, ma presentano differenze interne. Vi sono proprietà condivise da tutti i componenti, come nel caso della categoria \emph{uccello}, per cui essere oviparo e avere il becco sono proprietà comuni, e altre proprietà che non tutti i componenti possiedono, come volare o avere le piume.
Occorre quindi distinguere tra \emph{proprietà essenziali} e \emph{proprietà tipiche}.
\begin{snippetdefinition}{Proprietà essenziali}
    Le \emph{proprietà essenziali} sono proprietà che definiscono l'appartenenza categoriale in negativo, poiché permettono di escludere dalla categoria gli elementi che non le possiedono.
\end{snippetdefinition}
\begin{snippetdefinition}{Proprietà tipiche}
    Le \emph{proprietà tipiche} sono proprietà specifiche aggiunte, soggette a eccezioni e cancellabili senza compromettere l'appartenenza alla categoria.
\end{snippetdefinition}
\begin{snippetexample}{Proprietà essenziali e tipiche degli uccelli}
    Per la categoria \emph{uccello}, proprietà essenziali sono l'essere oviparo e il possedere il becco. Se un animale non ha il becco e non è oviparo, non può far parte della categoria \emph{uccello}. Sono invece proprietà tipiche la capacità di volare, poiché alcuni uccelli come gli struzzi non volano, avere le piume, poiché i pinguini non hanno piume nel senso comune del termine, e avere le ali.
\end{snippetexample}
Tra proprietà essenziali e proprietà tipiche esiste una \emph{gerarchia di rilevanza}, perché le prime sono più importanti delle seconde.
\subsection{Problem solving e ragionamento}
Torniamo al messaggio ``Il gatto è sul tappeto''. Se non si tratta di un messaggio in codice, potreste trovarvi nella situazione in cui qualcuno vi segnala di far uscire il gatto rimasto nella stanza in cui tenete il tappeto nuovo. In questo caso usereste processi cognitivi di \emph{problem solving} e \emph{ragionamento}.
Entrambe le attività comprendono la combinazione delle informazioni contingenti con quelle immagazzinate in memoria: lo scopo è giungere a una conclusione o trovare una soluzione.
\subsubsection{Problem solving}
L'attività di \emph{problem solving} è parte fondamentale dell'esistenza di ciascuno. Molte attività quotidiane richiedono soluzioni, come passare un esame o portare a termine un lavoro entro una scadenza.
Alcuni problemi implicano una discrepanza tra ciò che si sa e ciò che sarebbe necessario sapere. Risolvere un problema significa ridurre questa discrepanza, trovando un modo per ottenere le informazioni mancanti.
\begin{center}

    \includegraphics[width=0.45\textwidth]{media_psychology/nine-dot-problem.jpg}
    \captionof{figure}{Problema dei nove punti: unire tutti i punti con quattro linee rette connesse senza staccare la penna dal foglio.}
\end{center}
\begin{center}

    \includegraphics[width=0.45\textwidth]{media_psychology/nine-dot-solution.jpg}
    \captionof{figure}{Soluzione del problema dei nove punti.}
\end{center}
\paragraph{Spazio del problema}
Un problema può essere definito come la percezione di una discrepanza tra lo \emph{stato attuale} e lo \emph{stato desiderato}. Ad esempio, una persona può essere al verde e desiderare di avere più soldi, ipotizzando di ridurre questa distanza trovando un lavoro part-time.
La definizione formale di un problema comprende tre elementi:
\begin{enumerate}
    \item uno \emph{stato iniziale}, cioè l'informazione incompleta o la condizione di insoddisfazione da cui si parte;
    \item uno \emph{stato finale}, cioè l'informazione o la condizione che si vuole ottenere;
    \item un insieme di \emph{operatori}, cioè i passaggi che bisogna compiere per muoversi dallo stato iniziale a quello finale.
\end{enumerate}
\begin{snippetdefinition}{Spazio del problema}
    Lo \emph{spazio del problema} è l'insieme formato da stato iniziale, stato finale e operatori disponibili per passare dall'uno all'altro.
\end{snippetdefinition}
Nella risoluzione di un problema, molte difficoltà iniziali sorgono perché uno di questi elementi non è ben definito.
\begin{snippetdefinition}{Problema ben definito}
    Un \emph{problema ben definito} è un problema in cui lo stato iniziale, lo stato finale e gli operatori sono chiaramente specificati. Il compito consiste nello scoprire come usare le operazioni disponibili per raggiungere la soluzione.
\end{snippetdefinition}
Anche un problema ben definito può essere difficile da risolvere. Nei problemi di algebra, ad esempio, si può usare un \emph{algoritmo}.
\begin{snippetdefinition}{Algoritmo}
    Un \emph{algoritmo} è una procedura passo dopo passo che fornisce sempre la soluzione corretta a un particolare tipo di problema.
\end{snippetdefinition}
Gli algoritmi sono più facilmente disponibili per i problemi ben definiti. Quando non sono disponibili, ci si affida alle \emph{euristiche}.
\begin{snippetdefinition}{Euristica}
    Un'\emph{euristica} è una strategia cognitiva o regola empirica usata come scorciatoia nella soluzione di problemi complessi.
\end{snippetdefinition}
\begin{snippetexample}{Euristica nel romanzo giallo}
    Leggendo un romanzo giallo, si potrebbe escludere il maggiordomo come assassino usando l'euristica secondo cui l'autore non sceglierebbe mai una trama così scontata.
\end{snippetexample}
\paragraph{Migliorare l'abilità di problem solving}
La risoluzione dei problemi può essere difficile perché spesso ci sono troppe cose da tenere in considerazione nello stesso momento. Per risolvere un problema è necessario pianificare la serie di operazioni da mettere in atto. Se la serie è troppo complessa, o se sono complesse le singole operazioni, diventa difficile trovare il modo di passare dallo stato iniziale a quello finale.
Un passaggio importante per migliorare l'abilità di problem solving consiste nel trovare un modo utile di rappresentare il problema, così che ogni operazione diventi possibile. Talvolta, trovare una rappresentazione utile significa pensare al problema in modo completamente diverso.
\begin{snippetdefinition}{Fissità funzionale}
    La \emph{fissità funzionale} è un blocco mentale che ostacola la risoluzione dei problemi, inibendo la percezione di una nuova funzione per un oggetto precedentemente associato a un altro scopo.
\end{snippetdefinition}
Ogni volta che si rimane bloccati su un problema, può essere utile chiedersi come lo si sta rappresentando e se esistano altri modi di pensare al problema o alle componenti della sua soluzione.
\begin{snippetexample}{Problema della candela}
    Immaginate di avere a disposizione una candela, una scatola di puntine da disegno e una bustina di fiammiferi. Il compito è attaccare la candela al muro sopra un tavolo, in modo che la cera sciolta non goccioli sul tavolo. Per risolvere il problema bisogna modificare la prospettiva abituale e vedere la scatola di puntine non come un contenitore, ma come una piattaforma.
\end{snippetexample}
\begin{center}

    \includegraphics[width=0.50\textwidth]{media_psychology/candle-problem-functional-fixedness.jpg}
    \captionof{figure}{Problema della candela come esempio di fissità funzionale.}
\end{center}
Spesso, quando si cerca di risolvere un problema, si mettono in atto forme speciali di pensiero chiamate \emph{ragionamento}: ragionamento deduttivo, ragionamento induttivo e ragionamento abduttivo.
\subsubsection{Ragionamento deduttivo}
Immaginate di dover andare al ristorante. Lungo la strada vi accorgete di non avere contanti e telefonate al locale per chiedere se accettano carte di credito Visa. Il ristorante risponde che vengono accettate tutte le principali carte di credito. A questo punto potete concludere con sicurezza che la vostra Visa sarà accettata.
La conversazione può essere riformulata secondo la struttura del sillogismo:
\begin{itemize}
    \item \emph{Premessa 1:} il ristorante accetta tutte le principali carte di credito;
    \item \emph{Premessa 2:} la Visa è una delle principali carte di credito;
    \item \emph{Conclusione:} il ristorante accetta la Visa.
\end{itemize}
\begin{snippetdefinition}{Ragionamento deduttivo}
    Il \emph{ragionamento deduttivo} consiste nel trarre conclusioni a partire da premesse, basandosi su regole logiche che procedono dal generale al particolare.
\end{snippetdefinition}
Nella vita quotidiana, tuttavia, il ragionamento deduttivo è influenzato sia dalle conoscenze pregresse sul mondo sia dalle risorse disponibili.
Consideriamo il seguente sillogismo:
\begin{itemize}
    \item \emph{Premessa 1:} tutte le cose che hanno un motore hanno bisogno di olio;
    \item \emph{Premessa 2:} le automobili hanno bisogno di olio;
    \item \emph{Conclusione:} le automobili hanno un motore.
\end{itemize}
Secondo le regole della logica, la conclusione non è valida, perché la prima premessa lascia aperta la possibilità che alcune cose senza motore abbiano comunque bisogno di olio. Ciò che non è valido per un problema logico, però, non è necessariamente falso nel mondo reale.
\begin{snippetdefinition}{Bias dovuto alla credenza}
    Il \emph{bias dovuto alla credenza} è la tendenza a giudicare valide le conclusioni che si ritengono credibili e non valide quelle che si giudicano non credibili.
\end{snippetdefinition}
In alcuni casi, l'abilità di utilizzare esperienze passate migliora la prestazione nei compiti di ragionamento.
\begin{snippetexample}{Compito di selezione di Wason}
    Immaginate di avere quattro carte: \texttt{A}, \texttt{D}, \texttt{4}, \texttt{7}. Il compito è determinare quali carte girare per verificare la regola: ``Se una carta ha una vocale da un lato, allora ha un numero pari dall'altro''. La maggior parte delle persone decide di girare \texttt{A}, scelta corretta, e \texttt{4}, scelta errata. Infatti, non importa cosa ci sia dietro il \texttt{4}; la regola non verrà comunque confutata. La mossa corretta è girare anche il \texttt{7}, perché se dietro ci fosse una vocale la regola sarebbe falsa.
\end{snippetexample}
Quando ai partecipanti è permesso applicare allo stesso esercizio conoscenze del mondo reale, le prestazioni migliorano. Se la regola diventa ``Se un cliente sta per bere una bevanda alcolica, allora deve avere almeno sedici anni'', è più facile capire che bisogna controllare chi \emph{beve birra} e chi ha \emph{quindici anni}.
\begin{center}

    \includegraphics[width=0.65\textwidth]{media_psychology/wason-selection-task.jpg}
    \captionof{figure}{Compito di selezione di Wason in forma astratta e in forma legata alla vita reale.}
\end{center}
\subsubsection{Ragionamento induttivo}
Supponiamo di arrivare davanti al ristorante e accorgerci solo allora di non avere contanti. All'esterno del locale non ci sono indicazioni utili per capire se la carta Visa sia accettata. Guardando attraverso la vetrina si nota una clientela elegante, un menu con prezzi alti e un quartiere socialmente benestante. Tutte queste osservazioni portano a credere che il ristorante probabilmente accetterà la carta di credito.
\begin{snippetdefinition}{Ragionamento induttivo}
    Il \emph{ragionamento induttivo} è una forma di ragionamento che utilizza gli indizi disponibili per generare conclusioni probabili ma non certe.
\end{snippetdefinition}
Gran parte della risoluzione dei problemi nella vita reale si affida al ragionamento induttivo. Se ci si chiude fuori casa o fuori dalla macchina, una buona strategia può consistere nel richiamare alla memoria soluzioni che hanno già funzionato in passato.
\begin{snippetdefinition}{Soluzione di problemi per analogia}
    La \emph{soluzione di problemi per analogia} consiste nello stabilire un'analogia tra la situazione corrente e una situazione pregressa, usando una soluzione già sperimentata in passato.
\end{snippetdefinition}
È però necessaria cautela: una soluzione già utilizzata può essere efficace, ma può anche ostacolare la risoluzione se esiste una differenza sostanziale tra la vecchia situazione e quella nuova.
\begin{snippetdefinition}{Disposizione mentale}
    Una \emph{disposizione mentale} è uno stato cognitivo preesistente che può migliorare la qualità e la velocità della percezione e della soluzione dei problemi, ma può anche inibire o distorcere l'attività mentale quando le vecchie modalità di pensiero non sono più adeguate.
\end{snippetdefinition}
La distinzione tra ragionamento deduttivo e induttivo è stata studiata anche attraverso tecniche di neuroimaging. In un esperimento, i partecipanti svolgevano compiti di ragionamento mentre erano sottoposti a PET. I problemi deduttivi richiedevano di valutare se le conclusioni fossero valide, mentre i problemi induttivi prevedevano premesse che portavano a conclusioni incerte.
I risultati mostrarono differenti modelli di attivazione cerebrale: il ragionamento deduttivo produceva una maggiore attivazione nell'emisfero destro, mentre il ragionamento induttivo portava a una maggiore attivazione nell'emisfero sinistro. Questo suggerisce che il ragionamento deduttivo coinvolga un tipo di analisi relativamente indipendente dal linguaggio, mentre il ragionamento induttivo attivi maggiormente la comprensione basata sul linguaggio.
\subsubsection{Ragionamento abduttivo}
Negli scambi comunicativi quotidiani, per capire ciò che l'altro comunica, si fa spesso ricorso al \emph{ragionamento abduttivo}.
\begin{snippetdefinition}{Ragionamento abduttivo}
    Il \emph{ragionamento abduttivo} consiste nel procedere a ritroso dagli effetti alle cause, nel tentativo di spiegare qualcosa che è già accaduto. Si procede per supposizioni cercando di individuare una soluzione.
\end{snippetdefinition}
Le soluzioni individuate tramite ragionamento abduttivo non sono sempre corrette. Questa forma di ragionamento, essendo la più facile, non è esente da rischi e spiegazioni erronee.
Il procedimento abduttivo è influenzato dai processi di \emph{fissazione attentiva}.
\begin{snippetdefinition}{Fissazione attentiva}
    La \emph{fissazione attentiva} è la concentrazione dell'attenzione su aspetti parziali e limitati di ciò che è stato detto o di ciò che è accaduto.
\end{snippetdefinition}
\begin{center}

    \includegraphics[width=0.58\textwidth]{media_psychology/deductive-inductive-abductive-reasoning.jpg}
    \captionof{figure}{Confronto schematico tra ragionamento deduttivo, induttivo e abduttivo.}
\end{center}
\subsection{Giudizio e presa di decisione}
I giudizi sono il risultato dell'applicazione delle risorse cognitive limitate a situazioni che richiedono una presa di decisione. I processi di giudizio e di presa di decisione devono operare in modo da consentire di gestire l'incertezza nel modo migliore.
\begin{snippetdefinition}{Giudizio}
    Il \emph{giudizio} è il processo attraverso cui si formano opinioni, si raggiungono conclusioni e si fanno valutazioni critiche degli eventi e delle persone.
\end{snippetdefinition}
\begin{snippetdefinition}{Presa di decisione}
    La \emph{presa di decisione} è il processo con cui si sceglie tra due o più alternative, accettando o rifiutando le opzioni disponibili.
\end{snippetdefinition}
Giudizio e presa di decisione sono due processi strettamente legati: il giudizio riguarda la formazione di valutazioni, mentre la decisione è più direttamente connessa all'azione comportamentale.
\subsubsection{Euristiche e giudizio}
Un modo formale per formulare un giudizio potrebbe consistere nel fare una lista di pro e contro e vedere quale delle due parti risulti più lunga. Tuttavia, questo procedimento è lungo e complesso, mentre nella vita reale è spesso necessario formulare giudizi rapidamente.
Secondo gli studiosi, le persone utilizzano più spesso le euristiche che i metodi formali di analisi.
\begin{snippetdefinition}{Euristica}
    Un'\emph{euristica} è una regola empirica informale, una scorciatoia cognitiva che riduce la complessità nella formulazione del giudizio.
\end{snippetdefinition}
I ricercatori hanno ipotizzato che gli esseri umani abbiano sviluppato una \emph{cassetta degli attrezzi adattativa}: un repertorio di euristiche veloci ed economiche che portano a giudizi corretti nella maggior parte dei casi.
La ricerca si è però soffermata soprattutto sui casi in cui le euristiche portano a giudizi non corretti. Poiché molti processi euristici operano senza intervento cosciente, è utile conoscere le circostanze in cui l'intervento conscio può essere necessario.
\paragraph{Euristica della disponibilità}
\begin{snippetdefinition}{Euristica della disponibilità}
    L'\emph{euristica della disponibilità} è una scorciatoia di giudizio in cui la valutazione dipende dalla disponibilità delle informazioni in memoria.
\end{snippetdefinition}
Questa euristica ha due componenti:
\begin{itemize}
    \item la \emph{facilità} con cui le informazioni vengono reperite;
    \item i \emph{contenuti della memoria} che si ha la sensazione vengano recuperati più facilmente.
\end{itemize}
Gli stessi indizi di recupero possono essere più o meno efficaci per accedere alla memoria a seconda del contesto in cui vengono usati. Cambiando il contesto di recupero, è possibile modificare il giudizio, perché la facilità con cui le informazioni vengono recuperate dalla memoria varia da contesto a contesto.
Si possono avere difficoltà anche quando l'informazione conservata in memoria contiene un errore.
\paragraph{Euristica della rappresentatività}
\begin{snippetdefinition}{Euristica della rappresentatività}
    L'\emph{euristica della rappresentatività} è una scorciatoia di giudizio per cui si assume che qualcosa appartenga a una categoria se possiede le caratteristiche tipiche dei membri di quella categoria.
\end{snippetdefinition}
Questa euristica si basa sull'uso di informazioni appartenenti al passato per formulare giudizi riferiti al presente. Tuttavia, può portare fuori strada quando si ignorano altre informazioni rilevanti.
\begin{snippetexample}{Avvocato e sport preferito}
    Se viene descritta una persona come un avvocato di successo, competitivo, molto alto, controllato e vanitoso, le persone possono essere portate a scegliere il \emph{tennis} come suo sport preferito. Tuttavia, la risposta più probabile è \emph{un gioco con la palla}, perché include anche il tennis.
\end{snippetexample}
L'implicazione per la vita quotidiana è che si può essere tratti in inganno da un'alternativa rappresentativa prima di aver considerato la struttura di tutte le opzioni disponibili.
Un altro uso della rappresentatività riguarda la valutazione complessiva di un evento. Se si chiede a una persona quanto le sia piaciuto un concerto, la risposta non corrisponde necessariamente alla media di tutti i momenti vissuti. La ricerca suggerisce che i valori rappresentativi spesso riflettono una media tra il momento più intenso e il momento finale dell'evento.
\begin{snippetexample}{Esperimento sull'acqua gelida}
    In un esperimento, i partecipanti affrontavano due prove dolorose immergendo le mani nell'acqua gelida. Nella prova breve tenevano le mani per \(60\) secondi in acqua a \(14^\circ\)C. Nella prova lunga, dopo i primi \(60\) secondi, tenevano le mani immerse per altri \(30\) secondi, mentre l'acqua raggiungeva \(15^\circ\)C. Quando venne chiesto quale prova avrebbero preferito ripetere, la maggioranza scelse quella lunga, perché il momento finale era meno doloroso.
\end{snippetexample}
\begin{center}

    \includegraphics[width=0.55\textwidth]{media_psychology/cold-pressor-peak-end-rule.jpg}
    \captionof{figure}{Esperimento sull'acqua gelida e ruolo del picco e della fase finale nel giudizio complessivo.}
\end{center}
\paragraph{Euristica dell'ancoraggio}
\begin{snippetdefinition}{Euristica dell'ancoraggio}
    L'\emph{euristica dell'ancoraggio} è una regola in base a cui i giudizi sul valore di un evento o di un esito mostrano aggiustamenti insufficienti, verso l'alto o verso il basso, rispetto a un valore di partenza.
\end{snippetdefinition}
\begin{snippetexample}{Prodotto di numeri}
    Tversky e Kahneman chiesero ai partecipanti di stimare rapidamente il risultato di \(1 \times 2 \times 3 \times 4 \times 5 \times 6 \times 7 \times 8\) oppure di \(8 \times 7 \times 6 \times 5 \times 4 \times 3 \times 2 \times 1\). Anche se le due operazioni hanno lo stesso risultato, la media delle stime era \(512\) nella prima condizione e \(2250\) nella seconda. La risposta corretta è \(40320\).
\end{snippetexample}
L'esercizio mostra che il giudizio resta troppo saldamente ancorato alla congettura iniziale. Nel primo caso le prime moltiplicazioni producono un valore basso, mentre nel secondo producono un valore più alto, influenzando la stima finale.
\begin{center}
    \includegraphics[width=0.50\textwidth]{media_psychology/anchoring-heuristic-example.jpg}
    \captionof{figure}{Esempio dell'euristica dell'ancoraggio.}
\end{center}
\subsubsection{La psicologia della decisione}
È necessario comprendere gli aspetti psicologici della presa di decisione per valutare se una scelta regga a un'analisi attenta.
Il modo in cui una domanda è formulata può avere conseguenze importanti sulle decisioni prese, influenzando la risposta stessa.
Uno dei modi più immediati per prendere una decisione è giudicare quale opzione porti al maggior guadagno o quale conduca alla minor perdita. Se venissero offerti \(5\) o \(10\) euro, non ci sarebbero difficoltà nel scegliere \(10\) euro.
\begin{snippetdefinition}{Punti di riferimento}
    I \emph{punti di riferimento} sono criteri rispetto ai quali un esito viene interpretato come guadagno o come perdita. Dipendono in parte dalle aspettative del decisore.
\end{snippetdefinition}
Quando una decisione si rivela sbagliata, spesso si prova \emph{rammarico}. Se il rammarico è prevedibile, e quindi le persone si aspettano di poter rimpiangere la propria scelta, è probabile che diventino più caute nel momento in cui devono decidere.

\section{Emozioni}
Oltre che dalle motivazioni, l'agire umano è largamente influenzato anche dalle \emph{emozioni}. La paura provata di fronte a un serpente comparso all'improvviso può paralizzare oppure indurre a scappare con tutte le proprie forze.
Immaginare una vita senza emozioni permette di comprendere quanto esse siano fondamentali: le emozioni svolgono infatti molteplici funzioni essenziali.
\begin{snippetdefinition}{Emozione}
    Un'\emph{emozione} è un'esperienza complessa, costituita da cambiamenti corporei e mentali che includono attivazione fisiologica, sentimenti, processi cognitivi, espressioni visibili e reazioni comportamentali specifiche, attivate in risposta a situazioni percepite come personalmente significative.
\end{snippetdefinition}
\begin{snippetexample}{Felicità}
    Quando si prova felicità, l'attivazione fisiologica può determinare un leggero aumento del battito cardiaco. Il comportamento manifesto può essere espressivo, come nel caso di un sorriso, oppure orientato all'azione, come quando si abbraccia una persona cara.
\end{snippetexample}
\subsection{Aspetti fisiologici delle emozioni}
Quando si prova una forte emozione, il cuore batte più velocemente, la respirazione aumenta e i muscoli si irrigidiscono. Oltre a questi cambiamenti evidenti, l'emozione coinvolge molti altri processi che agiscono in modo sinergico sull'intero organismo.
Il \emph{sistema nervoso autonomo}, o \emph{SNA}, prepara il corpo alla risposta emotiva attraverso l'azione del sistema simpatico e del sistema parasimpatico. L'equilibrio tra questi due sistemi dipende dalla qualità e dall'intensità della stimolazione.
Con una stimolazione moderata e piacevole si attiva maggiormente il \emph{sistema parasimpatico}. Nel caso di una stimolazione più intensa, invece, vengono coinvolti entrambi i sistemi.
Dal punto di vista fisiologico, emozioni forti come la paura o la rabbia attivano il sistema di reazione di emergenza del corpo, che lo prepara rapidamente a un potenziale pericolo. Il \emph{sistema nervoso simpatico} attiva il rilascio di ormoni, come adrenalina e noradrenalina, che determinano il rilascio di zucchero nel sangue da parte degli organi interni, aumentando pressione sanguigna, sudorazione e salivazione.
La ricerca si è concentrata sulle funzioni svolte dall'\emph{ipotalamo} e dal \emph{sistema limbico} nel controllo dell'integrazione degli aspetti ormonali e neuronali dell'attivazione emotiva.
\begin{snippetdefinition}{Amigdala}
    L'\emph{amigdala} è una struttura cerebrale coinvolta nella decodifica delle informazioni provenienti dagli organi sensoriali e nell'attribuzione di significato a esperienze emotivamente salienti, soprattutto negative.
\end{snippetdefinition}
Le tecniche di scansione cerebrale hanno permesso di costruire mappe delle risposte cerebrali specifiche associate alle diverse emozioni. Le emozioni positive e negative non attivano semplicemente risposte opposte nella stessa area corticale, ma producono una maggiore attività in molteplici regioni del cervello.
Studi recenti indicano l'esistenza di differenze individuali nella rappresentazione cerebrale delle emozioni, in particolare tra uomini e donne. Negli uomini, l'amigdala destra presenta una maggiore connettività funzionale con regioni cerebrali orientate verso l'esterno. Nelle donne, invece, l'amigdala sinistra comunica in modo esteso con regioni come l'ipotalamo, più orientate verso gli stimoli interni. Ciò suggerisce che uomini e donne possano essere biologicamente predisposti a rispondere in modo diverso agli eventi emotivi.
\subsubsection{Teoria periferica}
Le emozioni precedono i comportamenti oppure derivano da essi? Si urla contro qualcuno perché si è arrabbiati, oppure ci si sente arrabbiati perché si sta urlando?
Secondo \emph{William James}, si proverebbe un'emozione in seguito alla reazione corporea.
\begin{snippetdefinition}{Teoria periferica delle emozioni}
    La \emph{teoria periferica delle emozioni}, o teoria di \emph{James-Lange}, sostiene che l'emozione deriva dal feedback corporeo. La percezione di uno stimolo causa un'attivazione del sistema nervoso autonomo che conduce all'esperienza di una specifica emozione.
\end{snippetdefinition}
Questa teoria è detta \emph{periferica} perché assegna un ruolo primario alle reazioni viscerali, cioè alle azioni del sistema nervoso autonomo. Secondo tale prospettiva, non si urla perché si è arrabbiati, ma si è arrabbiati perché si urla.
\subsubsection{Teoria centrale}
Lo psicologo \emph{Walter Cannon} formulò alcune obiezioni alla teoria di James-Lange. Egli sostenne che l'attività viscerale fosse troppo lenta e poco rilevante per dare origine all'esperienza emotiva. Secondo Cannon, l'emozione richiede l'intervento del cervello tra lo stimolo in entrata e la risposta in uscita.
Per \emph{Philip Bard}, lo stimolo emotigeno determina contemporaneamente un'attivazione corporea attraverso il sistema nervoso simpatico e un'esperienza emotiva soggettiva attraverso la corteccia cerebrale.
\begin{snippetdefinition}{Teoria centrale delle emozioni}
    La \emph{teoria centrale delle emozioni}, o teoria di \emph{Cannon-Bard}, sostiene che uno stimolo emotigeno produca due reazioni simultanee: l'attivazione fisiologica e l'esperienza emotiva soggettiva. Le due reazioni sono parallele e non dipendono causalmente l'una dall'altra.
\end{snippetdefinition}
Secondo questa teoria, a differenza della teoria periferica, si urla contro qualcuno perché si è arrabbiati.
La teoria periferica di James-Lange e la teoria centrale di Cannon-Bard, pur essendo contrapposte, si sono dimostrate entrambe fondate ma parziali, perché non riescono da sole a cogliere la complessità delle emozioni.
\subsubsection{Teoria emotiva a due fattori}
Il passaggio a una comprensione più psicologica delle emozioni avviene intorno agli anni Sessanta.
\begin{snippetdefinition}{Teoria emotiva a due fattori}
    La \emph{teoria emotiva a due fattori}, proposta da \emph{Schachter}, sostiene che l'esperienza emotiva sia la conseguenza dell'attivazione fisiologica, o \emph{arousal}, e dell'attribuzione causale.
\end{snippetdefinition}
Entrambe le componenti sono necessarie perché si verifichi l'esperienza emotiva. In questa prospettiva, l'attivazione fisiologica è generica e indifferenziata: costituisce il primo passo della sequenza emotiva. Successivamente il soggetto valuta il proprio livello di attivazione, individua le etichette verbali più adatte per descriverla e attribuisce un significato alle proprie reazioni nel contesto in cui si trova.
La sequenza può essere riassunta così:
\begin{enumerate}
    \item \emph{attivazione fisiologica};
    \item \emph{valutazione dell'attivazione fisiologica}, attraverso etichette verbali;
    \item \emph{attribuzione causale} nel contesto specifico.
\end{enumerate}
\subsubsection{Teoria cognitivo-attivazionale}
Secondo \emph{Lazarus}, l'esperienza emotiva non può essere compresa soltanto in termini di ciò che accade alle persone o nel loro cervello, ma dipende dal modo in cui vengono valutate le interazioni in corso con l'ambiente.
Lazarus sottolinea inoltre che l'attribuzione causale spesso si verifica senza un pensiero consapevole. Quando un'esperienza passata lega certe emozioni a determinate situazioni, non è necessaria una ricerca esplicita nell'ambiente per interpretare la propria attivazione.
\begin{snippetdefinition}{Teoria cognitivo-attivazionale}
    La \emph{teoria cognitivo-attivazionale delle emozioni} sostiene che l'esperienza emotiva dipenda dal modo in cui vengono valutate soggettivamente le interazioni in corso con l'ambiente, anche in modo inconsapevole.
\end{snippetdefinition}
Questa prospettiva evidenzia un'influenza reciproca tra emozione e processo di attribuzione causale, che può portare anche a errori di attribuzione.
\subsubsection{Teorie dell'appraisal}
\begin{snippetdefinition}{Appraisal}
    L'\emph{appraisal} è il processo di valutazione cognitiva attraverso cui una situazione viene interpretata in funzione dei propri interessi, scopi e significati.
\end{snippetdefinition}
Le \emph{teorie dell'appraisal} sostengono che le emozioni scaturiscano in seguito a un atto di conoscenza e valutazione della situazione. Le emozioni sono quindi strettamente legate ai processi cognitivi, poiché la loro attivazione implica l'elaborazione cognitiva della situazione.
Ne consegue che le emozioni cambiano quando cambiano i significati o i valori attribuiti a una determinata situazione. Due individui che valutano in modo diverso la stessa situazione manifesteranno emozioni differenti.
Queste teorie hanno fornito un contributo fondamentale alla comprensione dell'esperienza emotiva, poiché mettono in evidenza la grande \emph{flessibilità} e \emph{versatilità} delle emozioni.
Secondo \emph{Scherer}, l'appraisal è di solito rapido e relativamente semplice e segue una sequenza lineare di \emph{controlli di valutazione dello stimolo}, o \emph{CVS}.
I principali controlli di valutazione dello stimolo sono:
\begin{enumerate}
    \item valutazione della \emph{novità} e della discrepanza dello stimolo rispetto alle aspettative;
    \item valutazione della \emph{qualità edonica}, cioè della piacevolezza o spiacevolezza intrinseca dello stimolo;
    \item valutazione della \emph{pertinenza} dello stimolo rispetto ai bisogni e agli scopi dell'organismo;
    \item valutazione delle capacità di \emph{coping}, cioè delle possibilità di far fronte allo stimolo;
    \item valutazione della \emph{compatibilità} con le norme sociali e con l'immagine di sé.
\end{enumerate}
\begin{snippetdefinition}{Coping}
    Il \emph{coping} è la valutazione delle proprie capacità di far fronte a uno stimolo emotigeno. Può essere primario, se riguarda la capacità di controllare l'evento, oppure secondario, se riguarda la capacità di gestire le proprie reazioni emotive. Può inoltre essere attivo, se prepara all'azione, o passivo, se prepara alla difesa.
\end{snippetdefinition}
\subsection{Manifestazione delle emozioni}
\emph{Charles Darwin}, nel volume \emph{L'espressione delle emozioni nell'uomo e negli animali} del 1872, si chiese se le espressioni facciali delle emozioni fossero culturalmente invarianti. Egli considerava le emozioni come stati mentali ereditati e specializzati, formatisi per relazionarsi con classi ricorrenti di situazioni.
\subsubsection{Teoria dei programmi affettivi}
Se la prospettiva evoluzionistica è corretta, ci si aspetta di trovare le medesime emozioni nei bambini di tutto il mondo.
Secondo \emph{Tomkins}, le emozioni sono strettamente associate alla realizzazione di scopi universali connessi alla sopravvivenza della specie e dell'individuo.
\begin{snippetdefinition}{Programma affettivo}
    Un \emph{programma affettivo} è uno specifico programma nervoso, evolutosi nel tempo, che regola un'emozione di base e consente alla specie di adattarsi efficacemente al proprio ambiente.
\end{snippetdefinition}
Questa prospettiva, sostenuta da \emph{Ekman}, \emph{Izard} e \emph{Panksepp}, afferma che ogni emozione di base sia regolata da uno specifico programma affettivo.
\subsubsection{Ipotesi dell'universalità delle espressioni emotive}
\emph{Ekman} e i suoi collaboratori hanno sostenuto l'esistenza di una gamma di espressioni emotive universali nella specie umana, probabilmente perché si tratta di componenti innate della nostra eredità evolutiva.
L'espressione facciale di ciascuna emozione di base, cioè \emph{collera}, \emph{disgusto}, \emph{paura}, \emph{gioia}, \emph{tristezza} e \emph{sorpresa}, sarebbe unica e universale, presente in tutte le culture e riconosciuta in modo attendibile.
La conclusione generale è che le persone, al di là delle differenze culturali, esprimono le emozioni di base secondo modalità simili e sono in grado di identificare le emozioni provate da altri individui interpretandone le espressioni facciali.
\begin{snippetnote}{Universalità ed emozioni di base}
    L'ipotesi dell'universalità riguarda soltanto le emozioni di base. Ekman non sostiene che tutte le espressioni facciali siano universali o che tutte le culture esprimano ogni emozione nello stesso modo.
\end{snippetnote}
\begin{center}
    \includegraphics[width=0.62\textwidth]{media_psychology/emotional-expressions-context.jpg}
    \captionof{figure}{Espressioni emotive e ruolo del contesto.}
\end{center}
Ekman ha definito la propria posizione \emph{teoria neuroculturale}, per indicare il contributo congiunto del cervello, attraverso l'evoluzione, e della cultura nell'espressione delle emozioni.
Alcune debolezze metodologiche hanno però messo in discussione questi risultati, soprattutto perché non considerano adeguatamente la comparsa delle espressioni facciali in contesti naturali.
L'ipotesi di un legame tra emozioni ed espressioni facciali non nega l'esistenza di significative variazioni culturali. È quindi opportuno parlare di interdipendenza tra componenti biologiche e influenze culturali nella manifestazione delle emozioni.
\subsubsection{La prospettiva contestualista}
La \emph{prospettiva contestualista}, proposta più recentemente da \emph{Masuda}, \emph{Mesquita}, \emph{Russell} e \emph{Fernandez-Dols}, studia la connessione tra le espressioni facciali delle emozioni e il contesto immediato.
\begin{snippetdefinition}{Prospettiva contestualista}
    La \emph{prospettiva contestualista} sostiene che le espressioni facciali assumano un valore emotivo definito soltanto in riferimento a una situazione specifica.
\end{snippetdefinition}
Le espressioni facciali hanno un alto valore di \emph{indessicalità}: fanno riferimento a una realtà specifica attraverso indizi contestuali. Senza il contesto, il loro significato non sarebbe pienamente decifrabile.
\begin{snippetdefinition}{Indessicalità}
    L'\emph{indessicalità} è la proprietà per cui un segno, come un'espressione facciale, rimanda a una realtà specifica attraverso l'uso sistematico di indizi contestuali.
\end{snippetdefinition}
Di norma, le espressioni facciali sono accompagnate da movimenti corporei, inclusi i gesti. Le informazioni provenienti dal corpo non sono quindi un'aggiunta secondaria, ma una componente essenziale per percepire le manifestazioni emotive. Anche la voce contribuisce al riconoscimento delle espressioni facciali, rendendolo più rapido, attendibile e robusto.
La percezione delle manifestazioni emotive è quindi \emph{multimodale} e \emph{corale}, poiché coinvolge l'organismo nelle sue diverse modalità. Anche il contesto fisico e sociale è fondamentale, perché le informazioni contestuali congruenti facilitano il riconoscimento delle emozioni.
\section{La personalità umana}
Provate a paragonare due vostri amici. In che cosa sono simili? In che cosa sono diversi? È probabile che l'analisi arrivi rapidamente a concentrarsi sulle loro personalità.
\begin{snippetdefinition}{Personalità}
    La \emph{personalità} è l'insieme complessivo delle caratteristiche psicologiche che influenzano i pattern comportamentali tipici di un individuo nel corso del tempo e in situazioni differenti.
\end{snippetdefinition}
Le teorie della personalità sono ipotesi sulla struttura e sul funzionamento delle personalità individuali. Ogni teoria ha due obiettivi fondamentali: comprendere l'unicità di ogni individuo e spiegare in che modo ogni personalità produca pattern comportamentali caratteristici.
\subsection{Teorie dei tratti}
Un approccio allo studio della personalità consiste nell'utilizzo di scale di misura di differenti tratti. Le persone sembrano avere una tendenza naturale a classificare il proprio comportamento e quello altrui lungo più dimensioni.
\begin{snippetdefinition}{Tratto}
    Un \emph{tratto} è una caratteristica o attributo duraturo che influenza il comportamento in molte situazioni differenti.
\end{snippetdefinition}
Alcuni teorici considerano i tratti come predisposizioni che causano il comportamento, mentre altri li usano come dimensioni descrittive utili a rendere conto dei pattern comportamentali osservati.
\subsubsection{Approccio di Allport}
\emph{Gordon Allport} considerava i tratti come i mattoni della personalità e la fonte dell'individualità. Secondo Allport, i tratti producono comportamenti coerenti, collegando e uniformando le reazioni di una persona a una molteplicità di stimoli.
I tratti possono agire come \emph{variabili intervenienti}, mettendo in relazione insiemi di stimoli e risposte che, a prima vista, potrebbero sembrare poco collegati tra loro.
\begin{snippetdefinition}{Variabile interveniente}
    Una \emph{variabile interveniente} è una variabile non direttamente osservabile che permette di spiegare il legame tra stimoli e risposte osservabili.
\end{snippetdefinition}
\begin{snippetexample}{Timidezza}
    Il tratto della \emph{timidezza} può spiegare perché una persona risponda in modo simile in situazioni sociali diverse, ad esempio parlando poco, evitando il contatto visivo o mostrando disagio.
\end{snippetexample}
Allport individuò tre tipi di tratti:
\begin{itemize}
    \item \emph{tratti cardinali:} tratti attorno ai quali una persona organizza la propria vita, come il sacrificarsi per gli altri nel caso di Madre Teresa di Calcutta;
    \item \emph{tratti centrali:} caratteristiche fondamentali di una persona, come onestà o pessimismo;
    \item \emph{tratti secondari:} caratteristiche personali specifiche che contribuiscono a prevedere il comportamento, ma non sono decisive per comprendere la personalità, come le preferenze alimentari.
\end{itemize}
Allport era interessato a scoprire la combinazione peculiare di questi tre tipi di tratti che rende ogni persona un'entità unica. Per esaminarli, promosse lo studio di casi.
\begin{center}

    \includegraphics[width=0.58\textwidth]{media_psychology/allport-trait-types.jpg}
    \captionof{figure}{Schema dei principali tipi di tratti secondo Allport.}
\end{center}
\subsubsection{Quanti tratti di personalità esistono?}
Allport scoprì, attraverso una ricerca sul vocabolario, che la lingua inglese possedeva \(17953\) aggettivi per descrivere le differenze individuali. Da allora, gli psicologi hanno cercato di stabilire quante dimensioni esistano e quali permettano una caratterizzazione utile e universale degli individui.
\emph{Raymond Cattell} usò la lista di aggettivi di Allport come punto di partenza per la propria ricerca, finalizzata a individuare un piccolo gruppo di dimensioni fondamentali. Propose l'esistenza di \(16\) fattori soggiacenti alla personalità umana, che chiamò \emph{tratti originari}, perché riteneva rappresentassero l'origine dei comportamenti esteriori considerati come personalità.
I sedici fattori di Cattell includevano contrapposizioni comportamentali come \emph{riservato} ed \emph{espansivo}, oppure \emph{fiducioso} e \emph{sospettoso}.
\emph{Hans Eysenck} identificò invece tre macro-dimensioni nei dati ottenuti dai test di personalità:
\begin{itemize}
    \item \emph{estroversione:} orientamento verso l'interno o verso l'esterno;
    \item \emph{nevroticismo:} stabilità o instabilità emotiva;
    \item \emph{psicoticismo:} tendenza a essere gentili e premurosi oppure aggressivi e antisociali.
\end{itemize}
Se si rappresentano graficamente le dimensioni di \emph{estroversione} e \emph{nevroticismo} su assi ortogonali, ogni individuo può corrispondere a un punto del cerchio: da molto introverso a molto estroverso, da molto instabile a molto stabile. I tratti collocati all'interno del cerchio descrivono le persone attraverso la combinazione di queste due dimensioni.
\subsubsection{Modello a cinque fattori}
In anni recenti è emersa una convergenza sull'idea che la struttura della personalità sia descritta in modo efficace da cinque fattori, sovrapposti solo parzialmente alle tre dimensioni di Eysenck.
\begin{snippetdefinition}{Modello a cinque fattori}
    Il \emph{modello a cinque fattori}, o \emph{Big Five}, è un modello della personalità secondo cui le differenze individuali possono essere descritte attraverso cinque grandi dimensioni.
\end{snippetdefinition}
Le cinque dimensioni sono molto ampie: ciascuna comprende numerosi tratti dotati di connotazioni specifiche, ma organizzati attorno a un'idea comune.
I cinque fattori sono:
\begin{itemize}
    \item \emph{apertura mentale};
    \item \emph{coscienziosità};
    \item \emph{estroversione};
    \item \emph{piacevolezza};
    \item \emph{nevroticismo}.
\end{itemize}
Anche nel modello dei Big Five ogni dimensione presenta due poli: da un lato si trovano i termini simili per significato al nome della dimensione, dall'altro quelli dal significato opposto.
\begin{center}

    \includegraphics[width=0.72\textwidth]{media_psychology/big-five-personality-dimensions.jpg}
    \captionof{figure}{Le cinque dimensioni del modello Big Five e i rispettivi poli.}
\end{center}

Lo spostamento verso il modello a cinque fattori rappresenta il tentativo di dare una struttura al vasto elenco di tratti individuati da Allport e Odbert attraverso l'analisi del vocabolario.
I cinque fattori contribuiscono a definire un sistema di classificazione capace di descrivere le differenze individuali attorno ad alcune dimensioni fondamentali.
\subsubsection{Neuroscienze della personalità}
Un ruolo importante nell'elaborazione degli stimoli emotivi è svolto dall'\emph{amigdala}. I ricercatori hanno ipotizzato che non tutte le amigdale, e quindi non tutte le persone, rispondano agli stimoli nello stesso modo.
Per verificare questa ipotesi, sono stati reclutati partecipanti con differenti livelli di \emph{estroversione}. Poiché l'estroversione è associata ad aspetti importanti della vita emotiva, si prevedeva che essa potesse influenzare l'elaborazione emotiva.
Durante la risonanza magnetica funzionale, ai partecipanti venivano mostrate immagini di visi impauriti, felici e neutri.
\begin{center}

    \includegraphics[width=0.58\textwidth]{media_psychology/extraversion-amygdala-activity.jpg}
    \captionof{figure}{Relazione tra estroversione e attività dell'amigdala nella risposta a volti emotivi.}
\end{center}
Dai risultati non emergeva un legame tra estroversione e risposte cerebrali di fronte a volti impauriti: questi attivavano sia l'amigdala sinistra sia quella destra in modo simile per tutti i livelli di estroversione. Di fronte a volti felici, invece, gli individui più estroversi mostravano una maggiore attività dell'amigdala sinistra.
Questo studio suggerisce che le persone più inclini a relazionarsi con gli altri presentano una maggiore attivazione in regioni cerebrali collegate alle emozioni orientate al rapporto interpersonale.
\subsubsection{Il modello a cinque fattori: prospettive evoluzionistiche}
Le cinque dimensioni di personalità sono state messe in relazione con forme di interazione che gli esseri umani hanno avuto con sé stessi e con il mondo esterno nel corso dell'evoluzione.
Poiché gli esseri umani sono una specie sociale, la variazione nelle cinque dimensioni può essere interpretata come una serie di risposte a domande fondamentali sulla socievolezza:
\begin{itemize}
    \item richiesta di essere di buona compagnia, collegata all'\emph{estroversione};
    \item richiesta di essere gentili e di sostegno, collegata alla \emph{piacevolezza};
    \item richiesta di impegnarsi in sforzi prolungati, collegata alla \emph{coscienziosità};
    \item richiesta di essere emotivamente affidabili o inaffidabili, collegata al \emph{nevroticismo};
    \item richiesta di avere idee e riuscire a svilupparle, collegata all'\emph{apertura mentale}.
\end{itemize}
I ricercatori evoluzionistici si sono chiesti perché esista una variazione così ampia in ciascuna dimensione. Una possibile spiegazione è che, nel corso dell'evoluzione umana, la diversità degli ambienti abbia favorito valori differenti in ognuno dei tratti.
Se i tratti sono stati plasmati dall'evoluzione, ci si può aspettare che siano almeno in parte trasmessi da una generazione all'altra. Gli studi sull'ereditabilità mostrano infatti che quasi tutti i tratti della personalità sono influenzati da fattori genetici.
\subsubsection{Valutazione delle teorie dei tratti}
Le teorie dei tratti consentono ai ricercatori di offrire descrizioni sintetiche della personalità di individui differenti. Tuttavia, sono state criticate perché non spiegano in modo soddisfacente come si generi il comportamento o come si sviluppi la personalità.
Esse individuano e descrivono caratteristiche associate ai comportamenti, ma non spiegano pienamente i processi che li producono. Le teorie psicodinamiche, invece, pongono l'accento sulle forze conflittuali interne all'individuo, considerate responsabili del cambiamento e dello sviluppo.
\subsection{Teorie psicodinamiche}
\begin{snippetdefinition}{Teorie psicodinamiche}
    Le \emph{teorie psicodinamiche} della personalità condividono l'assunto secondo cui potenti forze interiori plasmano la personalità e motivano il comportamento.
\end{snippetdefinition}
Queste teorie prendono avvio da \emph{Sigmund Freud}. La sua teoria della personalità cerca di spiegare le origini e lo sviluppo della personalità, la natura della mente, gli aspetti atipici della personalità e il modo in cui essa può essere modificata attraverso la terapia.
\subsubsection{Psicanalisi freudiana}
Secondo la teoria psicanalitica, il centro della personalità è costituito da eventi interni alla mente dell'individuo, cioè eventi \emph{intrapsichici}, che motivano il comportamento. Per questo motivo, le persone possono agire anche su base inconscia.
Per Freud, tutto il comportamento è motivato: esso non deriva dal caso o da eventi accidentali, ma da motivi determinati.
\paragraph{Pulsioni e sviluppo psicosessuale}
La formazione neurologica di Freud lo portò a postulare l'esistenza di una base biologica comune ai pattern comportamentali. Egli riteneva che il comportamento fosse sempre motivato da un'energia psichica interna, costituita da pulsioni o istinti innati.
Freud individuò due pulsioni fondamentali:
\begin{itemize}
    \item la \emph{pulsione sessuale}, all'origine della componente erotica delle attività mentali;
    \item la \emph{pulsione aggressiva}, alla base della componente distruttiva.
\end{itemize}
\begin{snippetdefinition}{Libido}
    La \emph{libido} è, secondo Freud, la fonte di energia degli impulsi sessuali.
\end{snippetdefinition}
\begin{snippetdefinition}{Destrudo}
    La \emph{destrudo} è la componente energetica associata alla pulsione aggressiva.
\end{snippetdefinition}
Secondo Freud, la pulsione sessuale non compare improvvisamente nella pubertà, ma agisce fin dalla nascita. Essa si manifesta nel piacere che i bambini provano attraverso la stimolazione degli organi genitali e di altre aree sensibili, dette \emph{zone erogene}.
Le fasi dello sviluppo psicosessuale sono:
\begin{itemize}
    \item fase \emph{orale};
    \item fase \emph{anale};
    \item fase \emph{fallica};
    \item fase \emph{di latenza};
    \item fase \emph{genitale}.
\end{itemize}
\begin{snippetdefinition}{Fissazione}
    La \emph{fissazione} è l'incapacità di accedere alla fase successiva dello sviluppo psicosessuale, causata da eccessiva gratificazione o eccessiva frustrazione in una delle prime fasi.
\end{snippetdefinition}
Il concetto di fissazione mostra l'idea freudiana secondo cui i conflitti dei primi anni di vita contribuiscono a determinare i comportamenti successivi.
\begin{snippetdefinition}{Determinismo psichico}
    Il \emph{determinismo psichico} è l'assunto secondo cui tutti gli eventi della vita psichica sono determinati da esperienze precedenti e da episodi significativi, consci o inconsci.
\end{snippetdefinition}
Questo concetto mette in risalto l'\emph{inconscio}, cioè quella parte della psiche i cui contenuti sono rimossi dalla coscienza. Secondo Freud, il comportamento può essere motivato da pulsioni di cui non si è consapevoli.
Nel comportamento si può distinguere:
\begin{itemize}
    \item un \emph{contenuto manifesto}, cioè ciò che si dice, si fa e si percepisce consapevolmente;
    \item un \emph{contenuto latente}, cioè il significato nascosto e non immediatamente consapevole.
\end{itemize}
Freud distingue tre strutture fondamentali della personalità: \emph{Es}, \emph{Io} e \emph{Super-Io}.
\begin{snippetdefinition}{Es}
    L'\emph{Es} è la struttura primitiva e inconscia della personalità. Opera in modo irrazionale, agisce d'impulso ed esige la gratificazione immediata secondo il \emph{principio del piacere}, senza considerare se ciò che desidera sia realistico, socialmente desiderabile o moralmente accettabile.
\end{snippetdefinition}
\begin{snippetdefinition}{Super-Io}
    Il \emph{Super-Io} è l'istanza della psiche che comprende le norme sociali e morali e le aspirazioni ideali. Corrisponde in parte alla nozione comune di coscienza morale ed è spesso in tensione con l'Es.
\end{snippetdefinition}
\begin{snippetdefinition}{Io}
    L'\emph{Io} è l'aspetto della personalità coinvolto nelle attività di autoconservazione e nel dirigere le pulsioni verso sbocchi appropriati. Regola il rapporto tra individuo e ambiente e opera secondo il \emph{principio di realtà}.
\end{snippetdefinition}
L'Io media tra Es e Super-Io, anteponendo alla ricerca immediata del piacere le scelte basate sulla ragione.
Quando si trova in presenza di stati di angoscia, l'Io mette in atto strategie per difendersi dagli impulsi dell'Es e dalle richieste del Super-Io. Queste strategie sono chiamate \emph{meccanismi di difesa dell'Io}.
\begin{snippetdefinition}{Rimozione}
    La \emph{rimozione} è il meccanismo di difesa fondamentale attraverso cui pensieri, sentimenti o ricordi percepiti come minacciosi per la psiche vengono esclusi dalla coscienza.
\end{snippetdefinition}
I desideri ritenuti pericolosi per l'integrità dell'apparato psichico vengono respinti dalla coscienza e rinchiusi nell'inconscio. I meccanismi di difesa, pur essendo utili, sono forme di autoinganno: abusarne crea più problemi di quanti ne risolva.
\paragraph{Valutazioni della teoria freudiana}
Le idee di Freud hanno avuto un enorme impatto sull'interpretazione degli aspetti ordinari e anomali della personalità. Tuttavia, sono state criticate perché molti suoi concetti risultano vaghi e difficili da definire con precisione.
\subsubsection{Post psicanalisi freudiana}
Gli studiosi successivi a Freud ripresero la sua rappresentazione della personalità, apportando modifiche importanti:
\begin{itemize}
    \item attribuirono maggiore importanza alle funzioni dell'Io, incluse le difese, lo sviluppo del sé, i processi di pensiero consci e la padronanza di sé;
    \item riconobbero una maggiore influenza alle variabili sociali, come cultura, famiglia e pari;
    \item ridussero l'enfasi sugli impulsi sessuali generali;
    \item estesero lo sviluppo della personalità oltre l'infanzia.
\end{itemize}
\emph{Carl Jung} rielaborò ampiamente il concetto freudiano di inconscio, distinguendo tra inconscio personale e inconscio collettivo.
\begin{snippetdefinition}{Inconscio personale}
    L'\emph{inconscio personale} è formato da contenuti dimenticati o rimossi dall'individuo.
\end{snippetdefinition}
\begin{snippetdefinition}{Inconscio collettivo}
    L'\emph{inconscio collettivo} è una dimensione inconscia ereditata e condivisa dall'intera specie umana.
\end{snippetdefinition}
\begin{snippetdefinition}{Archetipo}
    Un \emph{archetipo} è un'immagine universale e innata, collocata nell'inconscio collettivo, che orienta il comportamento umano.
\end{snippetdefinition}
Jung considerava sana e integrata la personalità capace di equilibrare forze contrapposte, come aggressività maschile e sensibilità femminile.
\begin{snippetdefinition}{Psicologia analitica}
    La \emph{psicologia analitica} è la prospettiva di Jung secondo cui la personalità è una costellazione di forze interiori che si compensano in equilibrio dinamico, ponendo enfasi sul bisogno di realizzazione del sé.
\end{snippetdefinition}
\subsection{Teorie umanistiche}
Gli approcci umanistici alla personalità sono caratterizzati dall'interesse per la totalità dell'esperienza personale conscia e per il potenziale di crescita dell'individuo.
Per \emph{Carl Rogers}, il \emph{concetto di sé} è un modello mentale dei propri comportamenti tipici e delle proprie qualità peculiari.
\begin{snippetdefinition}{Concetto di sé}
    Il \emph{concetto di sé} è il modello mentale che un individuo possiede dei propri comportamenti tipici e delle proprie qualità personali.
\end{snippetdefinition}
Secondo Rogers, nel corso dell'esistenza ci si sforza di rendere il concetto di sé congruente con le reali esperienze di vita. Questa enfasi sul sé mostra una caratteristica fondamentale delle teorie umanistiche: la centralità del bisogno di \emph{autorealizzazione}.
\begin{snippetdefinition}{Autorealizzazione}
    L'\emph{autorealizzazione} è il costante sforzo di realizzare il proprio potenziale e di sviluppare talenti e capacità. Maslow la colloca al vertice della gerarchia dei bisogni.
\end{snippetdefinition}
Questo bisogno può entrare in conflitto con il bisogno di approvazione di sé e degli altri, soprattutto quando la persona percepisce che, per essere approvata, deve soddisfare obblighi o condizioni.
\begin{snippetdefinition}{Accettazione incondizionata}
    L'\emph{accettazione incondizionata} è l'accettazione totale di un individuo da parte di un altro, senza imporre condizioni.
\end{snippetdefinition}
Per Rogers, l'accettazione incondizionata è importante sia nell'infanzia sia nell'età adulta, perché la preoccupazione di essere approvati può interferire con l'autorealizzazione.
\emph{Karen Horney} arrivò a sostenere che le persone possiedono un \emph{vero sé}, che richiede condizioni ambientali favorevoli. In assenza di tali condizioni, può svilupparsi un'ansia di base che impedisce relazioni autentiche con gli altri, producendo movimenti verso, contro o lontano dagli altri.
Le teorie umanistiche sono state descritte come:
\begin{itemize}
    \item \emph{olistiche}, perché spiegano i singoli atti degli individui in relazione alla personalità complessiva;
    \item \emph{disposizionali}, perché si concentrano sulle qualità innate che influenzano la direzione del comportamento;
    \item \emph{fenomenologiche}, perché enfatizzano il punto di vista soggettivo dell'individuo sulla realtà.
\end{itemize}
\paragraph{Valutazione delle teorie umanistiche}
Le teorie umanistiche sono nate per valorizzare la personalità sana, orientata alla felicità e all'autorealizzazione. Tuttavia, i critici hanno sottolineato la vaghezza dei concetti chiave e la difficoltà di esaminarli empiricamente. Inoltre, gli psicologi umanistici tenderebbero a trascurare alcune importanti variabili ambientali che influenzano il comportamento.
\subsection{Teorie cognitive e dell'apprendimento sociale}
Le teorie considerate finora pongono l'accento su meccanismi interni ipotetici che sottendono il comportamento e formano la base della personalità. Tuttavia, esse non sempre stabiliscono un collegamento solido tra personalità e comportamenti specifici.
Da un'attenzione più diretta alle differenze individuali del comportamento è emersa una tradizione teorica centrata sulle contingenze ambientali, utile a prevedere gran parte del comportamento di una persona.
Gli psicologi dell'apprendimento analizzano le circostanze ambientali che controllano il comportamento.
Le teorie cognitive, invece, enfatizzano sia le circostanze ambientali sia i processi cognitivi e le differenze individuali. In questa prospettiva, esistono diverse personalità perché le persone hanno avuto esperienze differenti e perché partecipano attivamente alla costruzione della propria personalità.
I principali autori sono:
\begin{itemize}
    \item \emph{Rotter}, con la teoria dell'aspettativa;
    \item \emph{Mischel}, con la teoria cognitivo-affettiva della personalità;
    \item \emph{Bandura}, con la teoria cognitiva dell'apprendimento sociale.
\end{itemize}
\subsubsection{Teoria dell'aspettativa di Rotter}
\emph{Rotter} focalizzò la propria teoria sull'\emph{aspettativa}, cioè sulla misura in cui le persone credono che i loro comportamenti producano ricompense in situazioni specifiche.
\begin{snippetdefinition}{Aspettativa}
    L'\emph{aspettativa} è la credenza che un certo comportamento possa produrre una ricompensa in una determinata situazione.
\end{snippetdefinition}
Le aspettative derivano in parte dalla storia personale dei rinforzi.
Rotter sottolineò anche il \emph{valore della ricompensa}, cioè l'importanza soggettiva che un individuo attribuisce alla ricompensa stessa. Questo valore è un fattore determinante nella previsione del comportamento.
Rotter riteneva che le persone avessero aspettative specifiche rispetto alle molte situazioni affrontate nella vita, ma anche un'aspettativa più generale relativa alla misura in cui possono controllare le ricompense ottenute.
\begin{snippetdefinition}{Locus of control}
    Il \emph{locus of control}, o luogo di controllo, indica la credenza generale relativa al grado in cui i risultati ottenuti dipendono dalle proprie azioni o da fattori esterni.
\end{snippetdefinition}
Le persone con locus of control \emph{interno} credono che i risultati ottenuti dipendano soprattutto dalle proprie azioni. Le persone con locus of control \emph{esterno} credono invece che i risultati dipendano principalmente da fattori ambientali.
Il locus of control nell'infanzia è collegato al benessere psicofisico in età adulta.
\subsubsection{Teoria cognitiva dell'apprendimento sociale di Bandura}
La teoria di \emph{Bandura} combina i principi dell'apprendimento con le interazioni umane nei contesti sociali. In questa prospettiva sono centrali i processi cognitivi coinvolti nell'acquisizione e nel mantenimento dei pattern comportamentali e, quindi, della personalità.
La teoria mette in evidenza la complessa interazione tra fattori individuali, comportamento e stimoli ambientali. Ognuna di queste componenti può influenzare o cambiare le altre due.
\begin{snippetdefinition}{Determinismo reciproco}
    Il \emph{determinismo reciproco} è il principio secondo cui individuo, comportamento e ambiente interagiscono reciprocamente: ciascuna componente influenza ed è influenzata dalle altre.
\end{snippetdefinition}
\emph{Bandura} elaborò anche il costrutto di \emph{autoefficacia}.
\begin{snippetdefinition}{Autoefficacia}
    L'\emph{autoefficacia} è la credenza dell'individuo di poter eseguire una prestazione adeguata in una determinata situazione.
\end{snippetdefinition}
Il senso di autoefficacia influenza percezioni, motivazioni e prestazioni. Spesso non si tenta neppure di svolgere un compito o di assumersi un rischio quando ci si aspetta di essere inefficaci. Anche quando si possiedono capacità e desiderio di agire, si può rinunciare se si ritiene di non essere all'altezza.
Le principali fonti di informazione per formulare giudizi di autoefficacia sono:
\begin{enumerate}
    \item i \emph{risultati oggettivi};
    \item l'\emph{esperienza vicaria}, cioè l'osservazione delle prestazioni altrui;
    \item la \emph{persuasione}, da parte degli altri o di sé stessi;
    \item l'osservazione della propria \emph{attivazione emotiva} durante il compito o mentre lo si immagina.
\end{enumerate}
L'ansia può indicare basse aspettative di efficacia, mentre l'eccitazione può indicare aspettative di successo. Questi giudizi influenzano il livello di sforzo e la durata della persistenza di fronte alle difficoltà.
\begin{center}

    \includegraphics[width=0.55\textwidth]{media_psychology/bandura-self-efficacy-model.jpg}
    \captionof{figure}{Modello dell'autoefficacia di Bandura.}
\end{center}
Nel modello di Bandura, le aspettative di efficacia sono collocate tra la persona e il comportamento, mentre le aspettative di risultato sono collocate tra il comportamento e i risultati previsti.
\subsection{Teorie del sé}
L'interesse per l'analisi del sé ebbe come primo sostenitore \emph{William James}, che descrisse il sé come una struttura mentale dinamica capace di motivare, interpretare, organizzare, mediare e regolare i comportamenti e i processi intrapersonali e interpersonali.
James individuò tre componenti dell'esperienza di sé:
\begin{itemize}
    \item \emph{me materiale}, cioè il me corporeo e gli oggetti fisici che lo circondano;
    \item \emph{me spirituale}, cioè il sé che monitora pensieri e sentimenti intimi;
    \item \emph{me sociale}, cioè la consapevolezza di come gli altri ci vedono.
\end{itemize}
James riteneva che qualunque cosa associata alla propria identità diventasse, in un certo senso, parte del sé. Questo spiega perché le persone possono reagire in modo difensivo quando amici o parenti vengono attaccati.
\subsubsection{Concetto di sé e autostima}
Il concetto di sé contiene \emph{schemi di sé}, che permettono di organizzare le informazioni su sé stessi e influenzano il modo in cui vengono elaborate sia le informazioni personali sia quelle relative agli altri.
\begin{snippetdefinition}{Schema di sé}
    Uno \emph{schema di sé} è una struttura cognitiva che organizza le informazioni relative a sé stessi e guida il modo in cui tali informazioni vengono elaborate.
\end{snippetdefinition}
\begin{snippetdefinition}{Autostima}
    L'\emph{autostima} è una valutazione generale del sé, influenzata dall'interazione tra fattori genetici ed esperienze ambientali.
\end{snippetdefinition}
L'autostima può esercitare una forte influenza su pensieri, umori e comportamento.
Per preservare un'immagine positiva di sé, le persone ricorrono a strategie di \emph{autoaccrescimento}.
\begin{snippetdefinition}{Autoaccrescimento}
    L'\emph{autoaccrescimento} comprende strategie finalizzate a interpretare le proprie azioni e i propri comportamenti in modo costantemente positivo.
\end{snippetdefinition}
\begin{snippetdefinition}{Strategie autolesive}
    Le \emph{strategie autolesive} sono processi che, in previsione di un possibile fallimento, portano a produrre comportamenti o spiegazioni capaci di ridurre la probabilità di attribuire l'insuccesso alle proprie insufficienti capacità.
\end{snippetdefinition}
Lo scopo delle strategie autolesive è avere una scusa pronta per il fallimento, in modo che esso non implichi un giudizio di incapacità.
\subsubsection{Costruzione culturale del sé}
I ricercatori hanno iniziato a studiare il modo in cui il concetto di sé e lo sviluppo del sé sono influenzati dai condizionamenti culturali.
Il sé occidentale emerge da una cultura individualistica, che rappresenta una minoranza nella popolazione mondiale. La maggior parte delle culture è invece centrata sulla collettività.
Da ciascuna cultura derivano differenti interpretazioni del significato del sé, cioè differenti \emph{costrutti del sé}.
\begin{snippetdefinition}{Costrutto indipendente del sé}
    Il \emph{costrutto indipendente del sé}, tipico delle culture individualistiche, concepisce il sé come individuo il cui comportamento è organizzato principalmente attorno ai propri pensieri, sentimenti e azioni.
\end{snippetdefinition}
\begin{snippetdefinition}{Costrutto interdipendente del sé}
    Il \emph{costrutto interdipendente del sé}, tipico delle culture centrate sulla collettività, concepisce il sé come parte di una relazione sociale inclusiva, in cui il comportamento è organizzato anche dai pensieri, sentimenti e azioni altrui.
\end{snippetdefinition}
\subsection{Teorie a confronto}
È utile confrontare le principali teorie della personalità attraverso alcune contrapposizioni fondamentali.
\begin{enumerate}
    \item \emph{Ereditarietà vs. ambiente:} le teorie dei tratti sono divise su questo punto; la teoria freudiana si basa in larga misura sull'ereditarietà; le teorie umanistiche, dell'apprendimento sociale, cognitive e del sé pongono invece enfasi sull'ambiente o sull'interazione con esso.
    \item \emph{Apprendimento vs. leggi innate del comportamento:} la teoria freudiana esprime una visione deterministica interna, mentre le teorie umanistiche valorizzano il ruolo dell'esperienza. Le teorie dell'apprendimento sociale, cognitive e del sé sostengono che comportamento e personalità cambino per effetto delle esperienze apprese.
    \item \emph{Enfasi sul passato, sul presente o sul futuro:} le teorie dei tratti enfatizzano cause passate, innate o apprese; la teoria freudiana valorizza i processi inconsci infantili; le teorie dell'apprendimento sociale considerano rinforzi passati e contingenze presenti; le teorie umanistiche sottolineano realtà presente e obiettivi futuri.
    \item \emph{Conscio vs. inconscio:} la teoria freudiana pone l'accento sui processi inconsci, mentre le teorie umanistiche, dell'apprendimento sociale e cognitive valorizzano i processi consci.
    \item \emph{Disposizioni interne vs. situazioni esterne:} le teorie dell'apprendimento sociale sottolineano i fattori situazionali; le teorie dei tratti enfatizzano i fattori disposizionali; le altre prospettive propongono un'interazione tra variabili personali e situazionali.
\end{enumerate}
\subsection{Valutare la personalità}
Gli psicologi si interrogano sugli attributi che caratterizzano un individuo, distinguendolo dagli altri individui o da altri gruppi di persone.
Due assunti sono fondamentali:
\begin{itemize}
    \item esistono caratteristiche individuali che rendono coerente il comportamento di ciascuno;
    \item tali caratteristiche possono essere valutate o misurate.
\end{itemize}
I test di personalità devono soddisfare standard di \emph{affidabilità} e \emph{validità}. Inoltre, clinici e ricercatori ricevono una formazione specifica per poterli somministrare e interpretare.
\begin{snippetdefinition}{Inventario della personalità}
    Un \emph{inventario della personalità} è un test oggettivo basato su un questionario self-report, composto da item relativi a pensieri, sentimenti e azioni.
\end{snippetdefinition}
Chi si sottopone a un inventario della personalità legge una serie di affermazioni e indica se ciascuna è vera o tipica per sé.
Due famosi inventari della personalità sono:
\begin{itemize}
    \item \emph{Minnesota Multiphasic Personality Inventory}, o \emph{MMPI}, usato per diagnosi psichiatriche e dotato di scale cliniche per distinguere gruppi clinici da gruppi non psichiatrici;
    \item \emph{NEO-PI}, ideato per valutare caratteristiche di popolazioni adulte non cliniche e fondato sul modello della personalità a cinque fattori.
\end{itemize}
I test oggettivi possono presentare affermazioni a cui rispondere con opzioni semplici, come vero, falso o non so, oppure chiedere ai soggetti di valutare sé stessi rispetto a una certa dimensione. In entrambi i casi, il rispondente deve scegliere tra opzioni predeterminate.
\begin{snippetdefinition}{Test proiettivo}
    Un \emph{test proiettivo} è un test che presenta stimoli ambigui, come macchie, immagini incomplete o disegni interpretabili in più modi, chiedendo alla persona di descriverli, completarli o costruire storie a partire da essi.
\end{snippetdefinition}
Nei test proiettivi si assume che le risposte dell'individuo rivelino sentimenti, motivazioni e conflitti interiori. Poiché gli stimoli sono ambigui, le risposte dipendono in parte da ciò che la persona porta di sé nella situazione, cioè dalla propria esperienza di vita.
\begin{snippetdefinition}{Test di Rorschach}
    Il \emph{test di Rorschach} è un test proiettivo in cui vengono mostrate macchie d'inchiostro, alcune in bianco e nero e altre a colori, chiedendo alla persona che cosa potrebbero rappresentare.
\end{snippetdefinition}
Nel test di Rorschach, le risposte vengono registrate parola per parola, insieme al tempo impiegato e al modo in cui il soggetto manipola la tavola. Nella seconda fase, detta \emph{inchiesta}, viene chiesto al rispondente di precisare le risposte fornite. Il punteggio si basa su tre caratteristiche fondamentali:
\begin{itemize}
    \item la \emph{localizzazione};
    \item il \emph{contenuto} della risposta;
    \item i \emph{determinanti}.
\end{itemize}
\begin{snippetdefinition}{Test di appercezione tematica}
    Il \emph{test di appercezione tematica}, o \emph{TAT}, è un test proiettivo in cui vengono mostrate immagini ambigue e viene chiesto al soggetto di produrre storie che descrivano che cosa fanno e pensano i personaggi, che cosa ha originato gli eventi e come si concluderanno.
\end{snippetdefinition}
Chi somministra il TAT valuta la struttura e il contenuto delle narrazioni, oltre al comportamento dell'individuo che le racconta, nel tentativo di individuare preoccupazioni, motivazioni e caratteristiche della personalità.
Ogni tipo di test può fornire informazioni specifiche sulla personalità di un individuo. Per questo motivo, nelle valutazioni della personalità, i clinici usano spesso una combinazione di test.

\section{L'intelligenza}
L'\emph{intelligenza} è un costrutto teorico particolarmente ampio e difficile da caratterizzare. Le sue definizioni spaziano dalla capacità di ragionare astrattamente ai processi di adattamento all'ambiente.
Dopo diverse definizioni proposte da vari esperti, si è giunti alla conclusione che l'intelligenza è un'abilità mentale molto generale che coinvolge, tra le altre cose, capacità di ragionamento, pianificazione, \emph{problem solving}, pensiero astratto, comprensione di idee complesse, apprendimento rapido e apprendimento esperienziale.
\begin{snippetdefinition}{Intelligenza}
    L'\emph{intelligenza} è un'abilità mentale generale che comprende ragionamento, pianificazione, soluzione di problemi, pensiero astratto, comprensione di idee complesse, apprendimento rapido e apprendimento dall'esperienza.
\end{snippetdefinition}
\subsection{Teorie dell'intelligenza}
Sono state formulate diverse teorie sulle funzioni intellettive.
\emph{Sternberg} distingue tra:
\begin{itemize}
    \item \emph{teorie esplicite dell'intelligenza:} teorie formulate da esperti e fondate sulle risposte di gruppi di soggetti a compiti considerati idonei a misurare il comportamento intelligente;
    \item \emph{teorie implicite dell'intelligenza:} teorie che riguardano le idee che le persone hanno sulla propria intelligenza e su quella altrui. Esse sono alla base delle teorie esplicite, permettono di ripensarle e mettono in luce differenze culturali nelle opinioni sulle funzioni intellettive.
\end{itemize}
\subsubsection{Teorie fattoriali}
Alcuni autori, interrogandosi su come fosse organizzata l'intelligenza, utilizzarono procedure matematico-statistiche raggruppabili sotto il nome di \emph{analisi fattoriale}.
\begin{snippetdefinition}{Analisi fattoriale}
    L'\emph{analisi fattoriale} è un insieme di tecniche matematico-statistiche che permette di sintetizzare le informazioni espresse da molte variabili di partenza, tra loro non indipendenti, in un numero inferiore di variabili ipotetiche soggiacenti.
\end{snippetdefinition}
Queste tecniche hanno un importante valore \emph{economico}, perché riducono la complessità dei dati individuando fattori comuni sottostanti.
\begin{snippetexample}{Capacità induttiva}
    Se i componenti di un campione ottengono punteggi simili in test di ragionamento induttivo, come classificare, costruire serie logiche o rispondere a quesiti analogici, si può presumere che l'analisi fattoriale individui un'unica capacità induttiva di base.
\end{snippetexample}
\begin{center}

    \includegraphics[width=0.52\textwidth]{media_psychology/factor-analysis-intelligence.jpg}
    \captionof{figure}{Rappresentazione schematica dell'analisi fattoriale applicata alle funzioni intellettive.}
\end{center}
Secondo \emph{Spearman}, nella sua \emph{teoria bifattoriale}, l'intelligenza è espressione di:
\begin{enumerate}
    \item un fattore comune sottostante, definito \emph{fattore generale} \(g\);
    \item fattori specifici \(s\), che caratterizzano le diverse abilità.
\end{enumerate}
Se al fattore \(g\) va attribuita principalmente la capacità intellettiva di una persona, il rendimento in una particolare prova, ad esempio una prova di vocabolario, dipende sia dall'incidenza del fattore \(g\) sia dal fattore specifico dell'abilità richiesta per quella prova.
Altri ricercatori ottennero risultati differenti, pur utilizzando l'analisi fattoriale.
Secondo \emph{Thurstone}, alla base dei processi intellettivi vi sarebbero diverse \emph{abilità mentali primarie}, o \emph{PMA}. Le sette che hanno trovato maggiori conferme nelle ricerche sono:
\begin{itemize}
    \item \emph{comprensione verbale} \(V\);
    \item \emph{fluidità verbale} \(W\);
    \item \emph{abilità numerica} \(N\);
    \item \emph{relazioni spaziali} \(S\);
    \item \emph{memoria associativa} \(M\);
    \item \emph{velocità percettiva} \(P\);
    \item \emph{induzione} \(I\).
\end{itemize}
Tuttavia, applicando la propria batteria di test a un gruppo di bambini, Thurstone osservò che i sette fattori erano correlati tra loro e non potevano quindi essere considerati completamente indipendenti. Per questo motivo ipotizzò la presenza di un fattore \(g\) superiore, capace di spiegare la quota di interdipendenza tra le abilità primarie.
Secondo \emph{Cattell}, si possono distinguere due fattori principali: \(Gf\) e \(Gc\).
\begin{snippetdefinition}{Intelligenza fluida}
    L'\emph{intelligenza fluida}, o \(Gf\), è l'abilità di elaborare informazioni e ragionare in modo astratto. È connessa ad aspetti biologici ed è quindi maggiormente soggetta a deterioramento.
\end{snippetdefinition}
\begin{snippetdefinition}{Intelligenza cristallizzata}
    L'\emph{intelligenza cristallizzata}, o \(Gc\), è costituita dai contenuti che la persona ha appreso nel corso della propria esperienza, come vocabolario, comprensione verbale e informazioni di carattere generale.
\end{snippetdefinition}
Questa teoria fu poi ripresa da \emph{Horn}. Cattell e Horn ritengono che a una maggiore intelligenza fluida corrisponda un più alto grado di intelligenza cristallizzata, a condizione che la persona non viva in un ambiente socioculturale molto svantaggiato.
Horn ipotizza inoltre che le abilità cristallizzate possano aumentare nel corso del ciclo di vita, mentre le abilità fluide crescerebbero negli anni giovanili e mostrerebbero una tendenza a diminuire in età avanzata.
\begin{center}

    \includegraphics[width=0.52\textwidth]{media_psychology/fluid-crystallized-intelligence-lifespan.jpg}
    \captionof{figure}{Andamento dell'intelligenza fluida e cristallizzata nel corso del ciclo di vita.}
\end{center}
\subsubsection{La teoria triarchica dell'intelligenza di Sternberg}
\emph{Sternberg} ha sviluppato la \emph{teoria triarchica dell'intelligenza}, composta da tre subteorie.
\begin{snippetdefinition}{Teoria triarchica dell'intelligenza}
    La \emph{teoria triarchica dell'intelligenza} di Sternberg descrive l'intelligenza attraverso tre dimensioni: componenziale, esperienziale e contestuale.
\end{snippetdefinition}
Le tre subteorie sono:
\begin{enumerate}
    \item \emph{subteoria componenziale:} riguarda gli aspetti interni dell'intelligenza, cioè i meccanismi mentali che sottendono il comportamento intelligente. Comprende metacomponenti, componenti di prestazione e componenti di acquisizione di conoscenze;
    \item \emph{subteoria esperienziale:} riguarda l'intelligenza che l'individuo dimostra nell'uso della propria esperienza, e comprende la capacità di affrontare le novità e di automatizzare l'elaborazione;
    \item \emph{subteoria contestuale:} sostiene che il comportamento intelligente sia definito dal contesto socioculturale in cui si esprime, attraverso adattamento, selezione e modellazione.
\end{enumerate}
Le \emph{metacomponenti} sono le componenti di livello più elevato che le persone intelligenti userebbero per progettare i propri compiti. Le \emph{componenti di prestazione} sono processi mentali di tipo esecutivo, mentre le \emph{componenti di acquisizione di conoscenze} permettono di apprendere le informazioni necessarie alla soluzione dei problemi.
Sternberg ipotizza inoltre che l'attività intelligente sia di tre generi fondamentali:
\begin{itemize}
    \item \emph{analitica:} capacità di analizzare, giudicare, valutare, stabilire confronti e contrasti ed esaminare;
    \item \emph{creativa:} abilità di creare, scoprire, produrre, immaginare e supporre;
    \item \emph{pratica:} capacità di usare strumenti, applicare e attuare progetti e piani.
\end{itemize}
\begin{center}
    \includegraphics[width=0.58\textwidth]{media_psychology/sternberg-triarchic-intelligence.jpg}
    \captionof{figure}{Rappresentazione della teoria triarchica dell'intelligenza.}
\end{center}
\subsubsection{Teoria delle intelligenze multiple di Howard E. Gardner}
Nella revisione del 1999 della teoria delle intelligenze multiple, \emph{Howard E. Gardner} individua nove tipi di intelligenza, ciascuno funzionante in modo indipendente dagli altri e ancorato a un proprio substrato neurologico.
\begin{snippetdefinition}{Teoria delle intelligenze multiple}
    La \emph{teoria delle intelligenze multiple} sostiene che non esista una sola intelligenza generale, ma una pluralità di intelligenze relativamente indipendenti, combinate in modo peculiare in ogni individuo.
\end{snippetdefinition}
Gardner sostiene che ogni persona sia contraddistinta da una combinazione specifica di intelligenze, che possono funzionare in modo più o meno efficace a seconda dei casi.
\begin{center}

    \includegraphics[width=0.55\textwidth]{media_psychology/gardner-multiple-intelligences.jpg}
    \captionof{figure}{Rappresentazione delle intelligenze multiple secondo Gardner.}
\end{center}
\subsubsection{Intelligenza emotiva}
Secondo \emph{Goleman}, l'\emph{intelligenza emotiva} consiste in un insieme di abilità che includono il controllo dei propri impulsi, l'automotivazione, l'empatia e la competenza sociale nelle relazioni interpersonali.
\begin{snippetdefinition}{Intelligenza emotiva}
    L'\emph{intelligenza emotiva} è l'insieme di abilità legate alla gestione delle emozioni proprie e altrui, come controllo degli impulsi, automotivazione, empatia e competenza sociale.
\end{snippetdefinition}
In questo costrutto si intrecciano aspetti cognitivi e affettivi. Per questo motivo ci si può chiedere quanto un comportamento emotivamente intelligente rifletta una specifica capacità intellettiva e quanto invece sia il prodotto di una personalità ben adattata, emotivamente stabile e attenta ai bisogni propri e altrui.
\subsection{La valutazione dell'intelligenza}
I primi strumenti di misura dell'intelligenza risalgono alla seconda metà dell'Ottocento e consistevano in \emph{test mentali}.
\emph{Galton} riteneva che le famiglie britanniche dei ceti più elevati fossero intellettualmente più dotate e che tale dotazione fosse trasmessa alle generazioni successive. Il suo obiettivo era quindi dimostrare l'ereditarietà dell'intelligenza.
Le sue valutazioni avevano lo scopo di misurare caratteristiche fisiche e psicologiche come statura, peso, acutezza visiva e tempi di reazione visivi e uditivi, partendo dal presupposto che l'intelligenza fosse legata alla discriminazione sensoriale.
Questo approccio fu criticato da \emph{Alfred Binet}, che propose alternative diverse per valutare le capacità intellettive.
La valutazione standardizzata dell'intelligenza iniziò all'inizio del Novecento proprio con Binet. Il \emph{test di Binet-Simon} valutava abilità come comprensione, immaginazione, memoria e attenzione. Esso era finalizzato a identificare i bambini con difficoltà di apprendimento, così da poterli seguire con un percorso capace di porre rimedio alle loro carenze.
Il test fu applicato a bambini e preadolescenti dai \(3\) ai \(13\) anni e introdusse il concetto di \emph{età mentale}.
\begin{snippetdefinition}{Età mentale}
    L'\emph{età mentale}, o \(EM\), indica il livello di capacità cognitive di un individuo, equiparato al rendimento medio dei soggetti della stessa età cronologica.
\end{snippetdefinition}
\begin{snippetdefinition}{Età cronologica}
    L'\emph{età cronologica}, o \(EC\), è l'età anagrafica effettiva di un individuo.
\end{snippetdefinition}
\begin{snippetexample}{Età mentale ed età cronologica}
    Se un bambino di \(4\) anni di età cronologica supera le prove previste per la sua classe di età, si ipotizza che la sua età mentale coincida con la sua età anagrafica. Se invece risponde correttamente solo alle prove dei \(3\) anni, si suppone che la sua età mentale sia inferiore alla sua età cronologica.
\end{snippetexample}
Dopo la morte di Binet, \emph{Lewis M. Terman} rielaborò lo strumento nel 1916 alla Stanford University, creando la \emph{scala Stanford-Binet}. L'innovazione più nota fu l'introduzione del \emph{quoziente di intelligenza}, o \(QI\).
\begin{snippetdefinition}{Quoziente di intelligenza}
    Il \emph{quoziente di intelligenza}, o \(QI\), è un indice del funzionamento intellettivo. Nella formulazione originaria di Stern era ottenuto dal rapporto tra età mentale ed età cronologica, moltiplicato per \(100\).
\end{snippetdefinition}
\[
    QI = \frac{EM}{EC} \times 100
\]
Successivamente fu introdotto il \emph{QI-deviazione}, determinato in base alla distanza, espressa in unità di deviazione standard, del punteggio totale conseguito da un soggetto rispetto alla media aritmetica dei punteggi del proprio gruppo di età.
\begin{snippetdefinition}{QI-deviazione}
    Il \emph{QI-deviazione} è un indice calcolato in base alla distanza del punteggio di un individuo dalla media del proprio gruppo di età, espressa in unità di deviazione standard.
\end{snippetdefinition}
I punteggi di \(QI\) tendono a distribuirsi secondo la classica forma a campana.
\begin{center}

    \includegraphics[width=0.55\textwidth]{media_psychology/iq-normal-distribution.jpg}
    \captionof{figure}{Distribuzione dei punteggi di QI secondo una curva a campana.}
\end{center}
\subsubsection{La teoria Cattell-Horn-Carroll delle abilità cognitive}
La teoria \emph{Cattell-Horn-Carroll}, o \emph{CHC}, è il risultato dell'integrazione tra la teoria di Cattell-Horn, fondata su intelligenza fluida e cristallizzata, e la teoria a tre strati di Carroll.
\begin{snippetdefinition}{Teoria CHC}
    La \emph{teoria CHC} è una tassonomia delle abilità cognitive che integra la teoria dell'intelligenza fluida e cristallizzata con la teoria a tre strati di Carroll.
\end{snippetdefinition}
La teoria prevede tre strati:
\begin{itemize}
    \item \emph{strato III:} livello più elevato, corrispondente al fattore \(g\);
    \item \emph{strato II:} livello composto da abilità ampie, come intelligenza fluida, intelligenza cristallizzata, memoria, visualizzazione, velocità di elaborazione, recupero della memoria, percezione uditiva e fattore di sintesi di \(Gf\) e \(Gc\);
    \item \emph{strato I:} livello formato da numerosi fattori specifici.
\end{itemize}
\begin{center}
    \includegraphics[width=0.58\textwidth]{media_psychology/chc-cognitive-abilities-model.jpg}
    \captionof{figure}{Teoria CHC delle abilità cognitive.}
\end{center}
Questa teoria ha dato origine a una tassonomia articolata delle abilità cognitive, utilizzata per costruire o modificare importanti strumenti psicometrici.
Un importante strumento derivato da questa prospettiva è la \emph{scala d'intelligenza Wechsler per adulti}, formata da \(11\) subtest suddivisi in una subscala verbale e una subscala di performance, da cui si ottengono un \(QI\) verbale, un \(QI\) di performance e un \(QI\) totale.
La \emph{WISC-IV}, quarta versione della \emph{Wechsler Intelligence Scale for Children}, rivolta a soggetti tra \(6\) e \(16\) anni, è composta da \(15\) subtest divisi in quattro indici:
\begin{enumerate}
    \item \emph{indice di comprensione verbale}, o \(ICV\): capacità di utilizzare e formulare concetti verbali;
    \item \emph{indice di ragionamento visuo-percettivo}, o \(IVP\): ragionamento non verbale, pianificazione e uso di abilità visuo-spaziali o visuo-motorie;
    \item \emph{indice di memoria di lavoro}, o \(IML\): capacità di memorizzare nuove informazioni nella memoria a breve termine, mantenere l'attenzione focalizzata e usare tali informazioni per raggiungere una soluzione;
    \item \emph{indice di velocità di elaborazione}, o \(IVE\): capacità di focalizzare l'attenzione ed elaborare rapidamente gli stimoli.
\end{enumerate}
Questi quattro indici concorrono alla definizione del quoziente intellettivo.
\begin{center}
    \includegraphics[width=0.62\textwidth]{media_psychology/wisc-iv-cognitive-indexes.jpg}
    \captionof{figure}{Indici cognitivi della WISC-IV.}
\end{center}
Tra le prove collettive più note vi sono le \emph{Matrici Progressive di Raven}, costituite da stimoli figurativi geometrici. Esse richiedono di scegliere, tra più alternative, lo stimolo pertinente al quesito proposto.
\begin{center}
    \includegraphics[width=0.45\textwidth]{media_psychology/raven-progressive-matrices.jpg}
    \captionof{figure}{Matrici Progressive di Raven.}
\end{center}
\begin{snippetdefinition}{Matrici Progressive di Raven}
    Le \emph{Matrici Progressive di Raven} sono prove non verbali di intelligenza basate su stimoli geometrici, considerate strumenti relativamente \emph{culture free} perché non richiedono acquisizioni scolastiche specifiche.
\end{snippetdefinition}
\begin{center}
    \includegraphics[width=0.45\textwidth]{media_psychology/raven-nonverbal-matrix-item.jpg}
    \captionof{figure}{Esempio di quesito non verbale delle Matrici Progressive di Raven.}
\end{center}
\subsubsection{Affidabilità del QI}
L'affidabilità della misura del \(QI\), cioè la misura in cui essa è predittiva del quoziente intellettivo degli stessi individui a età diverse, varia in funzione dell'età: più l'età è alta, più la misura risulta affidabile.
Nei bambini molto piccoli è difficile ottenere valutazioni stabili, perché i test di intelligenza nei primi anni di vita fanno riferimento ad abilità differenti da quelle misurate in età successive. Nei primi anni, ad esempio, può essere valutato soprattutto lo sviluppo motorio, mentre non è ancora possibile valutare adeguatamente gli aspetti verbali.
I risultati dipendono anche dalla capacità del clinico di coinvolgere il bambino in modo proattivo, evitando che la madre o il caregiver interferiscano con la valutazione.
La valutazione diventa molto più affidabile dopo i \(7\)-\(8\) anni, anche per la maggiore uniformità dei test.
Il quoziente intellettivo è considerato un tratto relativamente stabile dell'individuo: un soggetto testato tende a mostrare valori simili nel corso del tempo. Tuttavia, questa stabilità diminuisce quando l'intervallo tra due misurazioni diventa molto lungo.
\begin{center}

    \includegraphics[width=0.58\textwidth]{media_psychology/iq-stability-over-time.jpg}
    \captionof{figure}{Stabilità del QI in funzione dell'intervallo temporale tra due misurazioni.}
\end{center}
Il \(QI\) può fluttuare verso il basso o verso l'alto a causa di variabili negative o positive. È inoltre sensibile alle influenze ambientali. Le fluttuazioni nel punteggio possono accompagnarsi anche a fluttuazioni a livello neurale, ad esempio nella connettività cerebrale o nel volume cerebrale.
\subsection{La disabilità intellettiva}
\begin{snippetdefinition}{Disabilità intellettiva}
    La \emph{disabilità intellettiva} è una condizione caratterizzata da significative limitazioni nel funzionamento intellettivo e nel comportamento adattivo, espresso in abilità concettuali, sociali e pratico-adattive. Ha origine prima dei \(18\) anni.
\end{snippetdefinition}
Si parla di disabilità intellettiva quando il valore del \(QI\) scende al di sotto di due deviazioni standard rispetto alla media.
Nel \emph{DSM-5} sono definiti diversi livelli di gravità, anche se la valutazione clinica non si basa soltanto sul calcolo del \(QI\):
\begin{enumerate}
    \item \emph{lieve:} \(QI\) da \(50\) a \(70\) punti;
    \item \emph{moderata:} \(QI\) da \(35\) a \(50\);
    \item \emph{grave:} \(QI\) da \(20\) a \(35\);
    \item \emph{gravissima:} \(QI\) inferiore a \(20\)-\(25\).
\end{enumerate}
Non sempre è facile distinguere tra questi gradi, perché può essere difficile coinvolgere pazienti con disabilità intellettiva nello svolgimento dei test. Anche l'abilità del clinico nel porre le domande ha un ruolo importante.
La diagnosi integra altre dimensioni, come gli aspetti emotivi, l'adattamento sociale e familiare e l'adattamento alla scuola o ad altri contesti. La valutazione serve infatti a fornire supporti, risorse e strategie finalizzate a promuovere educazione, interessi e benessere della persona, migliorandone il funzionamento individuale.
\subsection{Persone particolarmente dotate sul piano intellettivo}
Nel caso della \emph{plusdotazione}, o \emph{giftedness}, il \(QI\) supera di solito di due deviazioni standard la media generale, collocandosi intorno a \(130\)-\(140\) punti. Questa condizione deve essere osservabile già in età evolutiva.
\begin{snippetdefinition}{Plusdotazione}
    La \emph{plusdotazione} è una condizione in cui il quoziente intellettivo supera di circa due deviazioni standard la media generale e si accompagna al raggiungimento di traguardi di eccellenza nei settori di attività dell'individuo.
\end{snippetdefinition}
Le persone plusdotate si distinguono per il raggiungimento di traguardi di eccellenza nei settori di cui si occupano.
Un \(QI\) elevato non dipende soltanto da basi genetiche ed ereditarie, ma anche dalla componente educativo-ambientale.
Il \(QI\) può essere predetto o associato al livello di educazione parentale, cioè all'ambiente in cui il bambino cresce. In media, i bambini cresciuti da caregiver con un titolo di studio più elevato tendono a manifestare un \(QI\) più alto, probabilmente perché ricevono stimolazioni differenti e un diverso supporto educativo.
Anche l'aspetto socio-economico è influente: maggiore è il prestigio occupazionale dei genitori, maggiore è la probabilità che il bambino presenti un \(QI\) più elevato. Un esempio indiretto di investimento culturale familiare è il numero di libri presenti in casa.
\subsection{Intelligenza ed età}
Per quanto riguarda le capacità cognitive degli anziani, vi sono poche prove a sostegno dell'idea di un declino generalizzato nei soggetti in buona salute.
Si tende a pensare che gli anziani vadano incontro a un declino generale delle capacità e abilità, ma questa idea è troppo semplice. L'\emph{intelligenza cristallizzata}, ad esempio, non subisce un decremento significativo. In alcune sfere, invece, si può osservare un certo declino, come nel caso della memoria di lavoro o della memoria a lungo termine.
Alcune attività mostrano una relazione positiva con il \(QI\), favorendone una migliore conservazione. Le persone che nel corso della vita hanno svolto attività sociali, fisiche e intellettuali mostrano una più elevata velocità di processazione nei compiti cognitivi.
\begin{snippetdefinition}{Riserva cognitiva}
    La \emph{riserva cognitiva} è un bagaglio di risorse a livello neurale che consente di far fronte in modo più efficace a eventi traumatici o al declino cognitivo.
\end{snippetdefinition}
Esistono questionari che misurano la riserva cognitiva nel corso dello sviluppo, come il \emph{CRIq}. Per la terza età vengono poste domande sul lavoro pregresso, sulle attività svolte nel tempo libero e sulle attività ancora praticate negli ultimi anni.
Gli individui con punteggi più alti in questi questionari tendono a mostrare, nella terza età, un minore decremento cognitivo.
\subsection{Basi genetiche dell'intelligenza}
Alcuni studiosi sostengono che l'intelligenza sia in larga misura geneticamente determinata, ma non si dispone di certezze definitive a questo riguardo. Gli studi condotti per quantificare il peso delle componenti genetiche sull'intelligenza si sono basati soprattutto sul \(QI\).
\begin{snippetdefinition}{Indice di ereditabilità}
    L'\emph{indice di ereditabilità} fornisce una stima del grado in cui la variabilità di una caratteristica in una popolazione è determinata dalla variabilità genetica.
\end{snippetdefinition}
Gli studi sui gemelli mostrano che il \(QI\) è influenzato in modo importante dalla componente genetica. Ad esempio, gemelli monozigoti che condividono lo stesso ambiente domestico e il \(100\%\) dei geni presentano un \(QI\) altamente correlato.
Tuttavia, la componente genetica può essere modulata dalla condizione ambientale. Ad esempio, un genitore e un figlio biologico che non vivono nello stesso ambiente domestico possono presentare una correlazione molto più bassa nei punteggi di \(QI\).
\begin{center}

    \includegraphics[width=0.62\textwidth]{media_psychology/iq-genetics-environment-correlations.jpg}
    \captionof{figure}{Relazione tra grado di parentela, ambiente condiviso e correlazione nei punteggi di QI.}
\end{center}
Questi dati mettono in luce il peso elevato della componente genetica, ma mostrano anche che l'ambiente svolge un ruolo tutt'altro che trascurabile.

\section{Comunicazione e linguaggio}
La \emph{comunicazione} è un'esperienza che accompagna l'individuo per tutta l'esistenza. Comunicare è un'attività riconducibile a un comportamento di interazione, è una necessità umana e compare fin dai primissimi giorni di vita.
\begin{snippetexample}{Pianto del bambino}
    Il pianto di un bambino può essere distinto in base alla finalità comunicativa. Il bambino non piange sempre per dolore: può piangere per fame, noia o per un dolore specifico, come una colica. La \emph{sensibilità materna} consiste nella capacità di sintonizzarsi con il proprio figlio o la propria figlia, interpretando correttamente questi segnali.
\end{snippetexample}
Ricerche recenti hanno utilizzato l'intelligenza artificiale per evidenziare come sia possibile differenziare i diversi tipi di pianto sulla base delle loro caratteristiche acustiche. In questi studi veniva chiesto alla madre di indicare la possibile causa del pianto del bambino, anche quando i pianti apparivano simili.
A partire dai primissimi mesi di vita, il bambino manifesta la necessità di comunicare non solo attraverso il pianto, ma anche tramite sorrisi e vocalizzazioni. Questo repertorio comunicativo si arricchisce con l'avanzare dell'età, soprattutto attraverso l'acquisizione del linguaggio.
\subsection{I punti di vista sulla comunicazione}
\subsubsection{Il punto di vista matematico: comunicazione come scambio di informazioni}
Lo studio della comunicazione si è sviluppato a partire dagli anni Cinquanta anche grazie all'approccio matematico di \emph{Shannon}, fondato sullo studio dell'informazione.
\begin{snippetdefinition}{Informazione}
    Per Shannon, l'\emph{informazione} è una grandezza fisica, osservabile e misurabile, cioè una relazione tra due o più dati capace di generare ulteriori conoscenze.
\end{snippetdefinition}
L'informazione possiede alcune proprietà:
\begin{itemize}
    \item è \emph{espansiva} e \emph{diffusiva};
    \item è \emph{comprimibile};
    \item è facilmente \emph{trasformabile}.
\end{itemize}
Shannon identifica nel proprio modello non solo l'informazione, ma anche il suo inverso, cioè l'assenza di informazione: l'\emph{entropia}.
\begin{snippetdefinition}{Entropia}
    L'\emph{entropia} è l'assenza di informazione, che si verifica quando il segnale e le interferenze alla fonte sono equiprobabili ed equiparabili.
\end{snippetdefinition}
Shannon e Weaver elaborarono nel 1949 un modello della comunicazione basato sulla trasmissione di un segnale da una fonte a un destinatario.
\begin{center}

    \includegraphics[width=0.68\textwidth]{media_psychology/shannon-weaver-communication-model.jpg}
    \captionof{figure}{Modello matematico della comunicazione di Shannon e Weaver.}
\end{center}
Il requisito fondamentale del modello di Shannon è che la fonte e il destinatario possiedano lo stesso \emph{codice}.
\begin{snippetdefinition}{Codice}
    Il \emph{codice} è l'insieme di regole che associa in modo coerente e univoco gli elementi di un sistema con quelli di un altro sistema, come accade nella struttura sintattica o nelle associazioni semantiche.
\end{snippetdefinition}
Se gli individui non possiedono lo stesso codice, l'informazione non può essere veicolata correttamente. Queste regole sono quindi alla base dello scambio comunicativo.
\subsubsection{Il punto di vista semiotico e semantico: comunicazione come produzione di significati}
Comunicare riguarda anche la capacità di produrre segni dotati di significato e di interpretarli.
\begin{snippetdefinition}{Processo di significazione}
    Il \emph{processo di significazione} è il processo attraverso cui un segno rimanda a un significato e a un oggetto o evento del mondo.
\end{snippetdefinition}
Il processo di significazione prevede tre componenti:
\begin{itemize}
    \item \emph{referente:} ciò di cui si parla, cioè l'oggetto o evento;
    \item \emph{referenza:} l'immagine o rappresentazione mentale corrispondente;
    \item \emph{simbolo:} l'immagine sonora o iconica dell'oggetto o evento.
\end{itemize}
Tra simbolo e referente esiste una relazione indiretta, mediata dalla rappresentazione mentale, cioè dalla referenza. Da questa relazione nasce il \emph{triangolo della significazione}.
\begin{center}

    \includegraphics[width=0.55\textwidth]{media_psychology/triangle-of-meaning.jpg}
    \captionof{figure}{Triangolo della significazione: relazione tra simbolo, referenza e referente.}
\end{center}
L'individuo, per comprendere il referente ed elaborarlo, deve passare attraverso uno stadio di rappresentazione mentale.
\begin{snippetdefinition}{Semiotica}
    La \emph{semiotica} è la disciplina che studia il processo di significazione.
\end{snippetdefinition}
La semiotica assume che il segno sia \emph{arbitrario} e \emph{oppositivo}: per denominare un oggetto è necessario rappresentarlo mentalmente, ma non esiste una relazione evidente e naturale tra significante e significato. Tale relazione è il risultato di un accordo comunicativo.
\emph{Peirce} definisce il segno come qualcosa \(A\) che conduce a un interpretante \(B\) in corrispondenza con qualcosa \(C\). Il segno assume quindi una funzione di rimando, perché rimanda a qualcosa di diverso da sé.
Peirce individua tre tipi di segni:
\begin{enumerate}
    \item \emph{icone:} segni fondati su una relazione di somiglianza con le proprietà del referente;
    \item \emph{indici:} segni fondati su un rapporto di analogia o connessione con l'oggetto a cui si riferiscono;
    \item \emph{simboli:} segni in cui la connessione con il referente ha natura convenzionale e arbitraria.
\end{enumerate}
\subsubsection{Il punto di vista pragmatico: comunicazione come pratica}
Per comunicare non è sufficiente disporre di una rete di significati: occorre anche saperli utilizzare in modo appropriato nelle diverse situazioni. Questo è l'obiettivo di studio della \emph{pragmatica}.
\begin{snippetdefinition}{Pragmatica}
    La \emph{pragmatica} studia l'uso concreto dei segni, delle parole, delle frasi e dei gesti in un dato contesto comunicativo.
\end{snippetdefinition}
Il significato è una realtà dinamica e dipende dal modo concreto in cui si usa una parola, una frase o un gesto in una determinata situazione.
La pragmatica esamina i rapporti tra un testo e il contesto in cui esso è manifestato. Inoltre, prende in esame i processi impliciti della comunicazione, attraverso cui si inferisce dal contesto il contenuto del testo.
\emph{Austin} propone nel 1962 la \emph{teoria degli atti linguistici}, secondo cui dire qualcosa significa anche fare qualcosa.
\begin{snippetdefinition}{Atto linguistico}
    Un \emph{atto linguistico} è un'azione compiuta attraverso il linguaggio.
\end{snippetdefinition}
Nel dire qualcosa si compiono tre tipi di azione:
\begin{enumerate}
    \item \emph{atto locutorio:} l'atto stesso del parlare;
    \item \emph{atto illocutorio:} l'intenzione comunicativa espressa dall'enunciato;
    \item \emph{atto perlocutorio:} gli effetti prodotti sull'interlocutore.
\end{enumerate}
Il modo in cui viene interpretato un enunciato dipende dalla forza dell'atto e dai suoi effetti sull'interlocutore. La forza dell'atto illocutorio può essere modulata tramite pause, accento, intonazione e altri aspetti espressivi.
Secondo \emph{Grice}, comunicare significa rispettare il \emph{principio di cooperazione}.
\begin{snippetdefinition}{Principio di cooperazione}
    Il \emph{principio di cooperazione} afferma che l'interlocutore dovrebbe fornire il proprio contributo alla conversazione in modo opportuno.
\end{snippetdefinition}
Il principio di cooperazione si articola in quattro massime:
\begin{enumerate}
    \item \emph{massima di quantità:} soddisfare la richiesta di informazioni senza eccedere;
    \item \emph{massima di qualità:} affermare solo ciò che si ritiene vero e di cui si hanno prove;
    \item \emph{massima di relazione:} essere pertinenti;
    \item \emph{massima di modo:} essere chiari, brevi e ordinati.
\end{enumerate}
Negli scambi reali, Grice distingue tra ciò che è detto e ciò che è significato. Tra questi due livelli esiste uno scarto che deve essere colmato attraverso l'\emph{implicatura conversazionale}.
\begin{snippetdefinition}{Implicatura conversazionale}
    L'\emph{implicatura conversazionale} è un processo di inferenza mentale che permette di andare oltre il significato letterale di un enunciato e di comprendere l'intenzione comunicativa di chi parla.
\end{snippetdefinition}
\subsubsection{Il punto di vista psicologico: comunicazione come relazione interpersonale}
Nella comunicazione modifichiamo le relazioni interpersonali che costituiscono il nostro mondo sociale e che contribuiamo a costruire.
La comunicazione è una sequenza ininterrotta e una spirale di messaggi in cui stimolo, risposta e rinforzo si sovrappongono e si fondono. Ogni scambio comunicativo è infatti, nello stesso tempo, una risposta a un messaggio precedente.
Questa condizione è spesso alla base dei conflitti interpersonali, perché tendiamo a \emph{linearizzare} e segmentare in modo arbitrario il processo circolare e continuo della comunicazione.
\begin{snippetdefinition}{Giochi senza fine}
    I \emph{giochi senza fine} sono dinamiche comunicative e psicologiche che tendono a rinnovarsi continuamente, perché sono alimentate dagli stessi interlocutori coinvolti.
\end{snippetdefinition}
\begin{center}

    \includegraphics[width=0.52\textwidth]{media_psychology/circular-communication-process.jpg}
    \captionof{figure}{Rappresentazione della comunicazione come processo circolare e continuo.}
\end{center}
Grazie a questi processi, la comunicazione è la matrice della nostra identità personale e sociale, che ha natura dialogica. L'identità non è una realtà oggettiva, ma relazionale: ciascuno è ciò che è anche perché gli altri non sono ciò che è lui.
\subsection{Le funzioni della comunicazione}
La comunicazione è un vincolo per la specie umana, data la sua natura sociale. Gli esseri umani vivono in gruppi composti mediamente da circa \(140\)-\(150\) persone: questo valore è noto come \emph{numero di Dunbar}.
\begin{snippetdefinition}{Numero di Dunbar}
    Il \emph{numero di Dunbar} indica la grandezza media dei gruppi umani, stimata intorno a \(140\)-\(150\) persone.
\end{snippetdefinition}
Per questa ragione, la specie umana è definita \emph{ultrasociale} e ha posto la comunicazione al centro dell'esistenza individuale e sociale.
La comunicazione svolge diverse funzioni contemporaneamente. Le principali sono:
\begin{itemize}
    \item la \emph{funzione proposizionale};
    \item la \emph{funzione relazionale};
    \item la \emph{funzione espressiva}.
\end{itemize}
\begin{center}

    \includegraphics[width=0.58\textwidth]{media_psychology/communication-functions.jpg}
    \captionof{figure}{Principali funzioni della comunicazione.}
\end{center}
\begin{snippetdefinition}{Funzione proposizionale}
    La \emph{funzione proposizionale} della comunicazione consente di elaborare, organizzare, trasmettere e condividere conoscenze facendo ricorso a proposizioni.
\end{snippetdefinition}
La funzione proposizionale, esclusiva della nostra specie, è di natura simbolica e permette di raggiungere elevati livelli di astrazione. Grazie a essa è possibile accrescere il patrimonio delle conoscenze attraverso apprendimento e memoria.
Le conoscenze dichiarative rendono comunicabile il pensiero nelle sue varie espressioni: esse sono processi di mappatura della realtà in un'ampia rete di informazioni-proposizioni collegate tra loro.
\begin{snippetdefinition}{Funzione relazionale}
    La \emph{funzione relazionale} della comunicazione consiste nel dare forma alla trama delle relazioni interpersonali nei rapporti tra noi e gli altri.
\end{snippetdefinition}
La comunicazione partecipa a:
\begin{itemize}
    \item \emph{generare e sviluppare una relazione}, attraverso segnali capaci di avviare un nuovo rapporto;
    \item \emph{mantenere e rinnovare una relazione}, alimentandola una volta costruita;
    \item \emph{cambiare una relazione}, modificando la rete dei rapporti;
    \item \emph{restaurare una relazione}, ristabilendo un rapporto logorato;
    \item \emph{estinguere una relazione}, riducendo progressivamente o improvvisamente i contatti.
\end{itemize}
\begin{snippetdefinition}{Funzione espressiva}
    La \emph{funzione espressiva} della comunicazione è la modalità attraverso cui si manifestano pensieri, emozioni, valori e ideali.
\end{snippetdefinition}
La funzione espressiva è alla base della creatività umana nelle sue diverse forme.
\subsubsection{Verso una definizione di comunicazione}
\begin{snippetdefinition}{Comunicazione}
    La \emph{comunicazione} è uno scambio interattivo osservabile fra due o più partecipanti, dotato di un certo grado di consapevolezza e intenzionalità reciproca, attraverso cui si condivide un percorso di significati sulla base di sistemi convenzionali propri della cultura di riferimento.
\end{snippetdefinition}
Da questa definizione deriva che la comunicazione non coincide con il comportamento.
\begin{snippetdefinition}{Comportamento}
    Il \emph{comportamento} è qualsiasi azione motoria di un individuo osservabile da un altro.
\end{snippetdefinition}
Comunicazione e comportamento sono categorie distinte ma in rapporto di inclusione: ogni comunicazione è un comportamento, ma non ogni comportamento è comunicazione.
Occorre distinguere anche comunicazione e interazione.
\begin{snippetdefinition}{Interazione}
    L'\emph{interazione} è qualsiasi contatto, fisico o virtuale, fra due o più individui, anche involontario, capace di modificare lo stato preesistente delle cose fra loro.
\end{snippetdefinition}
Anche in questo caso esiste un rapporto di inclusione: ogni comunicazione è un'interazione, ma non ogni interazione è comunicazione.
\subsubsection{Intenzionalità e comunicazione}
Ogni scambio comunicativo presuppone un certo grado di \emph{intenzionalità}.
\begin{snippetdefinition}{Intenzionalità}
    L'\emph{intenzionalità} è sia una proprietà essenziale della coscienza sia la proprietà di un'azione compiuta in modo deliberato e volontario per raggiungere uno scopo.
\end{snippetdefinition}
L'intenzione si suddivide in due ambiti:
\begin{itemize}
    \item \emph{intenzione antecedente:} proposito di fare qualcosa, progettazione e pianificazione;
    \item \emph{intenzione in azione:} capacità di intervenire in modo consapevole durante lo svolgimento di un'attività.
\end{itemize}
L'\emph{intenzione comunicativa} consiste nel rendere manifesto al destinatario e al parlante che il parlante vuole dire certe cose al destinatario.
\begin{snippetdefinition}{Pars pro toto}
    La \emph{pars pro toto} è il principio secondo cui, nel manifestare un messaggio, il parlante può esprimere solo una parte dei propri contenuti mentali.
\end{snippetdefinition}
Questo principio può portare al fenomeno dell'\emph{opacità comunicativa}.
\begin{snippetdefinition}{Opacità comunicativa}
    L'\emph{opacità comunicativa} è la condizione per cui il messaggio espresso non rende completamente trasparenti i contenuti mentali e le intenzioni del parlante.
\end{snippetdefinition}
\begin{snippetdefinition}{Totum ex parte}
    Il \emph{totum ex parte} è il principio secondo cui il ricevente attribuisce un'intenzione completa al messaggio del parlante sulla base di un numero limitato di indizi comunicativi.
\end{snippetdefinition}
\begin{snippetdefinition}{Assumere per garantito}
    L'\emph{assumere per garantito} è il principio secondo cui il ricevente tende ad accogliere il primo senso del messaggio che gli viene in mente, purché non sia immediatamente contraddetto da un altro significato.
\end{snippetdefinition}
\subsection{Il significato è il fondamento della comunicazione}
Il \emph{significato} è il cuore della comunicazione: senza significato non sarebbe possibile comunicare.
I significati non sono dati né esistono nella realtà come generi naturali. Sono invece prodotti culturali, esiti contingenti dell'intreccio degli scambi tra parlanti. Essi sono caratterizzati da un insieme di tratti interconnessi.
\begin{snippetdefinition}{Semantica a tratti}
    La \emph{semantica a tratti} studia come i significati si compongono e si articolano attraverso insiemi di tratti interconnessi.
\end{snippetdefinition}
\subsubsection{Il significato come insieme di condizioni necessarie e sufficienti}
\begin{snippetdefinition}{Semantica vero-condizionale}
    La \emph{semantica vero-condizionale} sostiene che il significato di una parola o di una frase sia dato dal rapporto tra linguaggio e realtà, cioè dalle condizioni in base a cui un enunciato può essere considerato vero.
\end{snippetdefinition}
Le \emph{condizioni di verità} sono proprietà che un certo stato di cose, all'interno di un dato mondo, deve possedere affinché un determinato enunciato possa essere considerato vero in quel mondo.
Questa concezione del significato come insieme di proprietà intrinseche e immodificabili è stata criticata. Essa presuppone infatti che il significato consista in un insieme finito di condizioni necessarie e sufficienti, dotato di una dimensione oggettiva.
\subsubsection{Il significato come valore linguistico}
\begin{snippetdefinition}{Semantica strutturale}
    La \emph{semantica strutturale} studia il linguaggio come entità autonoma di relazioni interne, cioè come rete di dipendenze in cui i fatti linguistici esistono l'uno in ragione dell'altro.
\end{snippetdefinition}
In questa prospettiva, il significato è determinato esclusivamente dal \emph{valore linguistico}. Il significato di una parola nasce dalle relazioni con le altre parole e consiste nel sistema di differenze che separa quella parola da tutte le altre.
\subsubsection{Il significato come comprensione dell'esperienza e come prototipo}
\begin{snippetdefinition}{Semantica cognitiva}
    La \emph{semantica cognitiva}, o \emph{interpretativa}, sostiene che il significato consista nel modo in cui gli individui comprendono la propria esperienza.
\end{snippetdefinition}
L'esperienza non coincide con la realtà in sé, ma è l'incontro tra il proprio punto di vista soggettivo e i dati esterni provenienti dal mondo.
Il problema dell'individualità consiste nel fatto che i significati dipendono dall'elaborazione e dall'uso che ne fanno i parlanti.
\subsubsection{Stabilità e variabilità del significato}
Il significato è una realtà complessa, capace di cambiare in funzione della situazione e, allo stesso tempo, di restare abbastanza stabile da evitare continue incomprensioni.
La radice di questo equilibrio instabile è l'interdipendenza tra \emph{testo} e \emph{contesto}: non vi è testo senza contesto e non vi è contesto senza testo.
\begin{snippetdefinition}{Contesto}
    Il \emph{contesto} è l'insieme delle restrizioni e delle opportunità biologiche, spaziotemporali, relazionali, istituzionali e culturali dell'\emph{hic et nunc} che, insieme a un dato testo, genera un messaggio dotato di senso.
\end{snippetdefinition}
Esiste un numero indefinito di contesti disponibili in riferimento allo stesso testo, e questi contesti sono incapsulati gerarchicamente l'uno dentro l'altro. L'inclusione dei contesti può cambiare radicalmente il significato del messaggio, anche a parità di altre circostanze.
La connessione tra testo e contesto consente di comprendere la stabilità e la variabilità del significato. Il testo influenza il contesto nello stesso tempo in cui il contesto influenza il testo.
\subsection{Comunicare con le parole}
\emph{Chomsky} sostiene che il linguaggio, pur manifestandosi in oltre \(7000\) lingue naturali, abbia una struttura unica e universale, di natura innata e biologica, grazie alla presenza di un organo del linguaggio.
In opposizione alla posizione universalista, la prospettiva relativista, avanzata tra gli altri da \emph{Sapir} e \emph{Whorf}, sostiene che esista un intreccio intrinseco tra pensiero, lingua e cultura: nascere in una certa cultura implica necessariamente apprenderla.
Oggi molti studiosi propongono una versione debole del \emph{relativismo linguistico}.
\begin{snippetdefinition}{Fono}
    Un \emph{fono} è un suono che l'apparato vocale umano è in grado di emettere.
\end{snippetdefinition}
Tra tutti i foni, studiati dalla \emph{fonetica}, solo alcuni assumono valore linguistico. Questi sono i \emph{fonemi}.
\begin{snippetdefinition}{Fonema}
    Il \emph{fonema} è un'unità fonica indivisibile e astratta, dotata di valore linguistico distintivo e oppositivo. Di per sé è privo di significato, ma permette di distinguere significati lessicali diversi.
\end{snippetdefinition}
Ciascuna lingua naturale è caratterizzata da uno specifico insieme di fonemi.
Nella costruzione di una lingua, le componenti fonetiche vengono aggregate dai parlanti per formare parole. Le parole possono essere segmentate in unità linguistiche più piccole, portatrici di una porzione del significato complessivo: i \emph{morfemi}.
\begin{snippetdefinition}{Morfema}
    Il \emph{morfema} è un'unità grammaticale minima, non ulteriormente scomponibile, isolabile e dotata di un proprio significato.
\end{snippetdefinition}
I morfemi si distinguono in:
\begin{itemize}
    \item \emph{morfemi liberi:} costituiscono una parola autonoma, come \texttt{il}, \texttt{infatti}, \texttt{adesso};
    \item \emph{morfemi legati:} richiedono la combinazione con altri morfemi, come in \texttt{gatt-o}, dove \texttt{gatt-} è morfema lessicale e \texttt{-o} è morfema grammaticale.
\end{itemize}
\begin{snippetdefinition}{Lessico}
    Il \emph{lessico} è l'insieme delle parole di una lingua.
\end{snippetdefinition}
Il lessico è un insieme finito ed eterogeneo, composto da nove categorie grammaticali di parole: nome, verbo, aggettivo, avverbio, pronome, articolo, preposizione, congiunzione e interiezione.
Qualsiasi lingua possiede una grande creatività espressiva, poiché permette di produrre un numero infinito di enunciati. Tale creatività, tuttavia, è governata da regole.
\begin{snippetdefinition}{Sintassi}
    La \emph{sintassi} è l'insieme organico e formale delle regole e dei procedimenti computazionali che consentono di disporre in ordine gerarchico gli elementi del lessico.
\end{snippetdefinition}
\subsection{Comunicare senza parole}
Per molti anni i nostri antenati hanno comunicato in modo non verbale attraverso azioni e movimenti corporei, come gesti e sguardo. Su questa piattaforma non verbale, nel corso dell'evoluzione umana, si è radicata e sviluppata la facoltà del linguaggio.
Esiste un'asimmetria strutturale tra linguaggio verbale e sistemi non verbali di comunicazione: il linguaggio verbale, nella sua manifestazione, non è indipendente, ma è sempre supportato e veicolato da qualche dispositivo non verbale.
Il linguaggio verbale non è autonomo, perché richiede sempre il supporto di un mezzo e di un dispositivo extralinguistico. I sistemi non verbali, invece, sono autonomi nel proprio ambito di riferimento, poiché sono di natura \emph{ostensiva}.
\begin{snippetexample}{Mimica facciale}
    La \emph{mimica facciale} è un insieme coerente e articolato di espressioni visibili, capace di mostrare segnali di diversa natura.
\end{snippetexample}
La potenza, la flessibilità e l'adattabilità del linguaggio sono connesse alla sua \emph{convenzionalità} e \emph{arbitrarietà}. Queste proprietà assegnano a ogni lingua un valore digitale, fondato sulla presenza o assenza dei tratti che caratterizzano le sue componenti.
I sistemi non verbali, al contrario, hanno un valore \emph{motivato} e \emph{iconico}: trattengono al proprio interno aspetti della realtà che intendono evocare e possiedono quindi un valore analogico, fatto di variazioni continue e graduate.
I sistemi non verbali svolgono in modo efficace una funzione relazionale, poiché permettono di avanzare e condividere bisogni, stati motivazionali, emozioni, desideri e richieste.
\subsubsection{Principali sistemi di comunicazione non verbale}
\paragraph{Il sistema vocale}
La \emph{voce} manifesta e trasmette numerose componenti di significato oltre alle parole stesse. Essa va intesa come una sostanza fonica, composta da una serie di fenomeni e processi vocali che permettono la trasmissione e la ricezione dei segnali a distanza.
Contrapposto alla voce vi è il \emph{silenzio}. In quanto assenza di parola, il silenzio costituisce una modalità strategica di comunicazione.
\begin{snippetdefinition}{Silenzio}
    Il \emph{silenzio} è una forma di comunicazione non verbale fondata sull'assenza di parola. Il suo valore comunicativo dipende dalla sua ambiguità e dalle regole culturali che stabiliscono dove, quando, come e perché usarlo.
\end{snippetdefinition}
Nelle culture occidentali il silenzio è spesso considerato una minaccia o una mancanza di cooperazione nella conversazione. Nelle culture orientali, invece, può essere interpretato come indicatore di fiducia, confidenza e armonia.
\paragraph{Espressioni facciali}
I movimenti del volto costituiscono un sistema semiotico privilegiato, perché il volto è una regione del corpo particolarmente capace di attirare attenzione e interesse.
Le espressioni facciali manifestano stati mentali, esperienze emotive e altri segnali comunicativi.
Secondo la \emph{prospettiva emotiva}, i movimenti facciali sono destinati a esprimere le emozioni provate dall'individuo.
\begin{snippetdefinition}{Teoria dei programmi affettivi}
    La \emph{teoria dei programmi affettivi}, sostenuta da Ekman, afferma che esiste un'universalità delle espressioni facciali specifiche per ciascuna emozione di base.
\end{snippetdefinition}
Secondo questa prospettiva esisterebbe una sorta di isomorfismo tra emozione ed espressione facciale: a ogni emozione corrisponderebbe una specifica espressione.
In alcune circostanze, però, gli individui possono intervenire volontariamente sui muscoli facciali per governare le espressioni. A questo proposito, Ekman parla di \emph{regole di esibizione}.
\begin{snippetdefinition}{Regole di esibizione}
    Le \emph{regole di esibizione} sono regole attraverso cui le persone modulano volontariamente l'espressione delle proprie emozioni in funzione della situazione.
\end{snippetdefinition}
Le principali regole di esibizione sono:
\begin{enumerate}
    \item \emph{accentuazione};
    \item \emph{attenuazione};
    \item \emph{neutralizzazione};
    \item \emph{simulazione}.
\end{enumerate}
In contrapposizione alla prospettiva emotiva, \emph{Fridlund} propone la \emph{prospettiva comunicativa}.
\begin{snippetdefinition}{Prospettiva comunicativa}
    La \emph{prospettiva comunicativa} sostiene che le espressioni facciali abbiano valore sociale, poiché permettono di comunicare agli altri, in modo flessibile, i propri obiettivi e desideri.
\end{snippetdefinition}
Secondo questa prospettiva, le espressioni facciali:
\begin{itemize}
    \item hanno un valore sociale;
    \item sono governate da una socialità implicita, perché è sempre presente un uditorio implicito;
    \item non sono necessariamente legate a stati emotivi interni;
    \item rendono meno netta la distinzione tra espressione autentica ed espressione artificiale.
\end{itemize}
La \emph{prospettiva relazionale}, invece, assume che le espressioni facciali siano orientate verso altre persone e implichino varie forme di coinvolgimento interpersonale, finalizzate a una condizione soddisfacente di aggiustamento interattivo.
\begin{snippetdefinition}{Effetto audience}
    L'\emph{effetto audience} è l'influenza della presenza reale o implicita di altri individui sulla comparsa, il mantenimento e il cambiamento delle espressioni facciali.
\end{snippetdefinition}
\paragraph{Gesti}
Gli scambi comunicativi sono di norma accompagnati dalla presenza di \emph{gesti}.
\begin{snippetdefinition}{Gesto}
    Un \emph{gesto} è un'azione motoria che partecipa all'elaborazione del significato di un messaggio.
\end{snippetdefinition}
I principali tipi di gesti sono:
\begin{enumerate}
    \item \emph{gesti iconici:} accompagnano l'azione del parlare;
    \item \emph{pantomime:} imitazioni motorie di azioni;
    \item \emph{emblemi:} gesti convenzionali, come il segnale \texttt{OK};
    \item \emph{gesti deittici:} movimenti compiuti per indicare;
    \item \emph{gesti motori};
    \item \emph{lingua dei segni}.
\end{enumerate}
Gesti e parole sono strettamente interdipendenti. Essi sono generati dalla stessa rappresentazione mentale, manifestano la medesima intenzione comunicativa, sono pianificati dallo stesso processo cognitivo e vengono realizzati in sintonia rispetto a un dato contesto d'uso.
In caso di \emph{afasia}, si osserva una scomparsa simultanea del linguaggio e dei gesti iconici associati. In modo simile, nelle balbuzie la riduzione della fluenza del parlato è associata a una riduzione analoga dei gesti.
\paragraph{Prossemica e aptica}
\emph{Prossemica} e \emph{aptica} sono sistemi di contatto.
\begin{snippetdefinition}{Prossemica}
    La \emph{prossemica} riguarda la percezione, l'organizzazione e l'uso dello spazio, della distanza e del territorio nei confronti degli altri.
\end{snippetdefinition}
La gestione dello spazio implica un equilibrio instabile tra avvicinamento e allontanamento. Da un lato abbiamo bisogno di creare e mantenere contatti con gli altri; dall'altro abbiamo bisogno di proteggere la nostra privatezza.
La regolazione di queste oscillazioni è mediata dalla \emph{territorialità}.
\begin{snippetdefinition}{Territorialità}
    La \emph{territorialità} è la regolazione dello spazio personale e sociale attraverso cui gli individui definiscono distanza, vicinanza e confini nelle relazioni.
\end{snippetdefinition}
Si possono distinguere diverse zone:
\begin{itemize}
    \item \emph{zona intima};
    \item \emph{zona personale};
    \item \emph{zona sociale};
    \item \emph{zona pubblica}.
\end{itemize}
Esistono importanti differenze culturali nella gestione della distanza. Nella cultura araba, ad esempio, è presente una cultura della vicinanza, in cui la distanza può essere valutata come freddezza o ostilità.
\begin{snippetdefinition}{Aptica}
    L'\emph{aptica} è il sistema di significazione e segnalazione che riguarda le azioni di contatto corporeo nei confronti degli altri.
\end{snippetdefinition}
Il contatto corporeo è un atto comunicativo non verbale primario e risponde a un bisogno fondamentale della specie umana. Anche in questo caso esistono differenze culturali: accanto a culture del contatto, come quella araba e quella latina, vi sono culture del non contatto, come quelle nordiche o giapponesi.

\section{Psicologia della salute e del benessere della persona}
\subsection{Salute e benessere}
La costituzione dell'\emph{OMS} del 1948 ha definito la salute come uno stato di completo benessere fisico, sociale e mentale, e non soltanto come assenza di malattia o infermità.
Nella promozione della salute, la salute non viene considerata una condizione astratta, ma un mezzo finalizzato a un obiettivo: in termini operativi, essa è una risorsa che permette alle persone di condurre una vita produttiva sul piano individuale, sociale ed economico.
\begin{snippetdefinition}{Salute}
    La \emph{salute} è una risorsa per la vita quotidiana e non lo scopo dell'esistenza. È un concetto positivo che valorizza le risorse sociali e personali, oltre alle capacità fisiche.
\end{snippetdefinition}
In base a questa definizione, una persona sana, cioè non affetta da malattia, non è necessariamente in salute. L'assenza di malattia non garantisce infatti uno stato di completo benessere.
Il \emph{benessere fisico} e il \emph{benessere psichico} sono condizioni ottimali che devono essere considerate in relazione al contesto socioculturale in cui la persona vive e si esprime.
Della felicità e del benessere si è occupata ampiamente la \emph{psicologia positiva}, a partire da due approcci:
\begin{itemize}
    \item \emph{approccio edonico:} concepisce la felicità come puro soddisfacimento di bisogni e desideri;
    \item \emph{approccio eudaimonico:} concepisce la felicità come autorealizzazione virtuosa e collettiva.
\end{itemize}
Questi due approcci condividono un aspetto comune: in entrambi, la felicità non è legata al possesso di ricchezze o beni materiali.
\begin{center}
    \includegraphics[width=0.62\textwidth]{media_psychology/hedonic-well-being.jpg}
    \captionof{figure}{Approccio edonico al benessere.}
\end{center}
\begin{snippetdefinition}{Edonismo}
    L'\emph{edonismo} è una concezione della felicità come stato di appagamento connesso al raggiungimento di un equilibrio omeostatico ottenuto attraverso il soddisfacimento dei desideri.
\end{snippetdefinition}
\begin{snippetdefinition}{Eudaimonia}
    L'\emph{eudaimonia} è una concezione della felicità come processo di costruzione continua, crescita di complessità, autorealizzazione e armonizzazione del benessere individuale con il benessere collettivo.
\end{snippetdefinition}
\begin{center}
    \includegraphics[width=0.62\textwidth]{media_psychology/eudaimonic-well-being.jpg}
    \captionof{figure}{Approccio eudaimonico al benessere.}
\end{center}
Eudaimonia ed edonismo si distinguono per tre aspetti principali:
\begin{enumerate}
    \item l'eudaimonia è un processo di costruzione continua e di crescita, mentre l'edonismo è uno stato di appagamento legato al soddisfacimento dei desideri;
    \item l'eudaimonia è correlata a interesse, apertura verso le novità, opportunità di crescita, impegno e sforzo, mentre l'edonismo è legato a piacevolezza e affettività positiva;
    \item nell'eudaimonia il benessere individuale si armonizza con il benessere collettivo, in un processo di reciproca influenza.
\end{enumerate}
Secondo la \emph{teoria del set point}, particolari condizioni di gioia o tristezza tendono a sfumare tornando verso l'abituale stato affettivo della persona.
\begin{snippetdefinition}{Set point}
    Il \emph{set point} è l'abituale stato affettivo di una persona, modificato solo transitoriamente da intense emozioni.
\end{snippetdefinition}
\subsection{Psicologia della salute: sviluppo storico, obiettivi e temi di ricerca}
\begin{snippetdefinition}{Psicologia della salute}
    La \emph{psicologia della salute} è un ambito di studio e di ricerca della psicologia finalizzato a fornire un contributo alla comprensione e al mantenimento della salute psicologica e fisica.
\end{snippetdefinition}
Secondo Matarazzo, la psicologia della salute comprende i contributi specifici della psicologia, scientifici, professionali e formativi, alla promozione e al mantenimento della salute, alla prevenzione e al trattamento della malattia e all'identificazione dei correlati eziologici e diagnostici della salute, della malattia e delle disfunzioni associate.
L'\emph{APA}, cioè \emph{American Psychological Association}, è un'organizzazione che rappresenta la psicologia negli Stati Uniti e promuove il contributo della psicologia alla comprensione della salute e della malattia.
I principali fattori che hanno contribuito alla nascita della psicologia della salute sono:
\begin{enumerate}
    \item il cambiamento delle cause principali di mortalità;
    \item i fattori economici;
    \item il contributo della psicologia rispetto all'approccio medico tradizionale.
\end{enumerate}
I principali settori della psicologia della salute sono:
\begin{itemize}
    \item \emph{promozione e mantenimento della salute};
    \item \emph{prevenzione e trattamento della malattia};
    \item \emph{correlati eziologici e diagnostici della salute e della malattia};
    \item \emph{sistema di cura della salute e linee di azione politica}.
\end{itemize}
Per quanto riguarda la ricerca in psicologia della salute, si possono distinguere tre ampie categorie di studi:
\begin{itemize}
    \item studi \emph{descrittivi};
    \item studi \emph{predittivi};
    \item studi \emph{basati su interventi}.
\end{itemize}
\begin{snippetexample}{Adolescenti e fumo}
    Immaginiamo di voler studiare i fattori psicosociali che inducono gli adolescenti a fumare. Gli studi descrittivi permettono di descrivere le relazioni tra variabili, ad esempio tra stato di salute e attività funzionali o disfunzionali al suo mantenimento. Gli studi predittivi mirano invece a prevedere la relazione tra comportamenti funzionali o disfunzionali e stato di salute nel tempo.
\end{snippetexample}
Gli studi basati sugli interventi cercano di tradurre le conoscenze teoriche in approcci mirati a prevenire comportamenti disfunzionali per lo stato di salute.
\subsection{Il modello biopsicosociale}
\begin{snippetdefinition}{Modello biomedico}
    Il \emph{modello biomedico} spiega la malattia esclusivamente sulla base di cause di natura organica, ben definite e identificabili.
\end{snippetdefinition}
\emph{Michael Balint}, nel volume \emph{Medico, paziente e malattia}, evidenzia un modello di cura fondato sul rapporto a due tra curante e paziente, in cui il farmaco più comunemente prescritto è il medico stesso.
\begin{snippetdefinition}{Modello biopsicosociale}
    Il \emph{modello biopsicosociale} è un orientamento concettuale che sostiene un approccio multidimensionale alla salute e alla malattia, considerate come il risultato di una complessa articolazione di piani biologici, psichici e sociali.
\end{snippetdefinition}
\begin{center}
    \includegraphics[width=0.62\textwidth]{media_psychology/biopsychosocial-health-model.jpg}
    \captionof{figure}{Modello biopsicosociale della salute.}
\end{center}
Questo modello risente della \emph{teoria generale dei sistemi} di von Bertalanffy, secondo cui:
\begin{itemize}
    \item il sistema è un'unità unica composta da parti in relazione tra loro;
    \item le parti tendono all'equilibrio;
    \item un cambiamento in una parte influenza la totalità del sistema.
\end{itemize}
\emph{Engel} attribuisce importanza al vissuto soggettivo della malattia: la sua gravità non dipende solo dalla condizione oggettiva, ma anche dal processo interpretativo dell'individuo.
\subsubsection{La medicina psicosomatica}
\begin{snippetdefinition}{Medicina psicosomatica}
    La \emph{medicina psicosomatica} studia i rapporti tra il funzionamento della psiche e quello del corpo in un'ottica terapeutica.
\end{snippetdefinition}
La medicina psicosomatica presuppone una visione olistica dell'uomo e del paziente, secondo cui i meccanismi di natura emozionale possono incidere negativamente sulla salute fisica.
In quanto metodo di studio della patologia umana, essa tiene conto di variabili a livelli differenti, psicologico, sociale e biologico, e riconosce un'\emph{eziopatogenesi multipla complessa}.
I conflitti psichici possono comportare un'iperattivazione del sistema neurovegetativo, favorendo l'insorgenza di malattie di varia gravità.
\subsection{La comunicazione al servizio della salute}
\begin{snippetdefinition}{Health literacy}
    La \emph{health literacy} è l'insieme delle conoscenze di base di area medica attraverso cui l'individuo mette in atto comportamenti virtuosi nella gestione del proprio benessere psicofisico.
\end{snippetdefinition}
La health literacy comprende conoscenze che devono essere apprese attraverso iniziative mirate.
Bassi livelli di health literacy influenzano negativamente lo stato di salute delle persone, rappresentano un fattore di rischio, aumentano le disuguaglianze e incrementano i costi a carico del sistema sanitario.
Alla base di questo concetto vi è il presupposto che esista una proporzionalità diretta tra alfabetizzazione, acculturazione e correttezza nella gestione della propria salute.
Le conseguenze negative di un grado insufficiente di health literacy sono molteplici: possono andare dall'assunzione di comportamenti e abitudini dannose per il proprio benessere psicofisico fino all'adesione a credenze magico-fatalistiche sulla salute.
Anche l'accesso improprio ai servizi di medicina di emergenza-urgenza rappresenta un problema attuale.
\subsection{Personalità e benessere}
La felicità appare un obiettivo difficile da ottenere; tuttavia, non mancano studi e ricerche che hanno cercato di fare luce su questo argomento.
\begin{snippetdefinition}{Psicologia positiva}
    La \emph{psicologia positiva} è un approccio allo studio della personalità che si concentra sulle caratteristiche di coloro che tendono a rispondere costruttivamente alle difficoltà della vita e mostrano una forte motivazione all'autorealizzazione.
\end{snippetdefinition}
\begin{snippetdefinition}{Felicità}
    La \emph{felicità} è la capacità di provare emozioni positive.
\end{snippetdefinition}
\begin{snippetdefinition}{Ottimismo}
    L'\emph{ottimismo} è una predisposizione interiore di fiducia verso ciò che può accadere nella vita.
\end{snippetdefinition}
Il metodo più utilizzato dai ricercatori consiste nel chiedere alle persone di riferire i propri vissuti di felicità e infelicità e di valutare quanto siano soddisfacenti le loro vite.
Altri metodi utilizzano immagini oppure affermazioni che devono essere valutate scegliendo tra alternative di risposta predeterminate.
Si può affermare che chi vive in nazioni ricche sia, in linea generale, più felice di chi vive in nazioni meno ricche. Tuttavia, la differenza non è di entità particolarmente elevata.
Quattro tratti sembrano caratterizzare le persone felici:
\begin{enumerate}
    \item sentirsi contenti di sé stessi;
    \item avere consapevolezza di possedere il pieno controllo della propria vita;
    \item essere ottimisti;
    \item essere estroversi.
\end{enumerate}
\subsection{Eventi stressanti e modalità di risposta}
\begin{snippetdefinition}{Stress}
    Lo \emph{stress} è un elevato tasso di attivazione fisica e psichica della persona in risposta a eventi interni ed esterni che la coinvolgono.
\end{snippetdefinition}
Lo stress coinvolge diversi sistemi e apparati dell'organismo, come il sistema respiratorio, cardiaco e urogenitale.
Partendo dallo stress come risposta aspecifica, è possibile definire la \emph{sindrome generale di adattamento}, o \emph{SGA}, come lo sviluppo delle modalità di risposta dell'individuo a qualsiasi tipo di evento stressante.
\begin{snippetdefinition}{Sindrome generale di adattamento}
    La \emph{sindrome generale di adattamento}, o \emph{SGA}, è il modello che descrive lo sviluppo delle risposte dell'individuo di fronte a eventi stressanti.
\end{snippetdefinition}
La sindrome generale di adattamento si articola in tre fasi:
\begin{enumerate}
    \item \emph{allarme:} l'individuo diventa consapevole della novità della situazione;
    \item \emph{resistenza:} l'evento stressante perdura nel tempo e l'organismo cerca di adattarsi;
    \item \emph{esaurimento:} si verifica una riduzione della resistenza.
\end{enumerate}
Lo stress non deve essere inteso necessariamente come esperienza negativa. È infatti opportuno distinguere tra \emph{eustress} e \emph{distress}.
\begin{snippetdefinition}{Eustress}
    L'\emph{eustress} è un processo di attivazione in cui la persona, di fronte a un evento che richiede l'utilizzo di energie, risponde in modo costruttivo.
\end{snippetdefinition}
\begin{snippetdefinition}{Distress}
    Il \emph{distress} è un processo conseguente a un evento stressante percepito come minaccioso e destabilizzante, rispetto al quale l'individuo si sente privo di mezzi di difesa.
\end{snippetdefinition}
In risposta agli eventi stressanti, l'organismo mette in atto meccanismi di \emph{coping} per affrontare positivamente le diverse situazioni.
\begin{snippetdefinition}{Coping}
    Il \emph{coping} è l'insieme dei meccanismi di fronteggiamento di eventi interni ed esterni alla persona, percepiti come minacciosi.
\end{snippetdefinition}
Fra le strategie di coping si distinguono:
\begin{itemize}
    \item strategie \emph{emotion-focused}, cioè centrate sulle emozioni;
    \item strategie \emph{problem-focused}, cioè centrate sui problemi.
\end{itemize}
Le strategie \emph{emotion-focused} hanno per obiettivo gli stati emotivi interni e sono dirette a modificare le risposte emotive. Le strategie \emph{problem-focused}, invece, consistono in tentativi di risoluzione dei problemi.
\begin{center}
    \includegraphics[width=0.58\textwidth]{media_psychology/stress-coping-resilience.jpg}
    \captionof{figure}{Stress, coping e resilienza.}
\end{center}
\begin{snippetdefinition}{Resilienza}
    La \emph{resilienza} è uno stile di coping che consiste nella capacità di gestire positivamente eventi dolorosi che si protraggono a lungo.
\end{snippetdefinition}

\end{document}
